\documentclass[11pt]{ctexart}
\usepackage[a4paper,margin=1in]{geometry}
\usepackage{graphicx}
\usepackage{xcolor}
\usepackage{hyperref}
\definecolor{refgray}{gray}{0.45}
\title{\textbf{市场最新观点汇总}\\[0.4em]{\large 宏观叙事切换：从地缘风险定价转向库存重建、AI供给再定价与结构性资金洪流}}
\author{KC桌面 自动生成}
\date{260617}
\begin{document}
\maketitle
\tableofcontents
\newpage
\section{一页摘要}
\begin{itemize}
  \item 市场对美联储鹰派定价已较充分，美元短期区间震荡，真正的交易机会在于地缘风险缓和后的美元走弱。
  \item 全球石油市场已消耗约10亿桶库存，即便霍尔木兹恢复通航，补库需求将在未来两年每天吸收约300万桶供给，油价底部坚实。
  \item AI正从聊天机器人迈向智能体时代，CPU与GPU配比从1:8回归1:1，服务器CPU市场2030年TAM上调至2230亿美元。
  \item 中国5月经济数据显示出口是唯一结构性亮点，内需疲软拉长政策等待期，三季度财政与货币宽松窗口正在打开。
  \item 中国房地产市场并非全面回暖，一线城市二手房连续3-4个月环比正增长，但全国总量疲软，定价权正从新房转向二手房。
  \item 美国债券市场持有者结构发生不可逆重置：外国占比降至30\%，一级交易商库存激增至5500亿美元，国内机构成为边际买家。
  \item 日本央行6月加息25bp至1.0\%，但鸽派委员反对票与前瞻指引变化暗示加息节奏可能慢于市场预期。
  \item 全球资本支出超级周期正在重塑资产定价逻辑，回报驱动力从估值扩张转向盈利增长，选股能力变得更为关键。
  \item 中国AI供应链在英伟达下一代GB300/Vera Rubin架构中将实现份额跃迁，技术升级与响应速度是核心驱动力。
  \item 美国证券化市场正经历需求冲击定价：全年发行量首破1万亿美元，但净供给仅2700亿，保险与年金资金洪流持续压缩利差。
\end{itemize}
\section{宏观与利率：美联储鹰派定价充分，美元区间待突破}
\textbf{市场已为鹰派声明充分定价，沃什的模糊表态才是大概率事件，美元短期区间震荡，中期看涨但需等待地缘风险缓和触发突破。}

\begin{itemize}
  \item Citi认为市场对美联储鹰派声明的定价已非常充分，包括删除宽松倾向、上调核心PCE预测、点阵图中位数指向2026年不降息，这些变化不足以引发显著外汇反应。
  \item 新任主席沃什的新闻发布会是真正变量，其减少前瞻指引的一贯立场可能给出模棱两可的回答，从而温和打压当前鹰派定价。
  \item 美元多头仓位已大幅回撤，DXY接近短期公允价值，短期可能继续区间震荡，但中期看涨逻辑基于相对增长差异。
  \item DB指出美国61万亿美元债券市场的持有者结构正经历不可逆重置：外国持有占比降至30\%，一级交易商国债库存激增至5500亿美元。
  \item 全球外汇储备美元占比持续走低，外国增持意愿减弱并非抛售，而是美国国债发行增长更快，国内机构必须成为边际买家。
  \item DB证券化团队判断市场定价逻辑已从信用风险转向需求冲击：2026年证券化发行量首破1万亿美元，但净供给仅2700亿。
  \item 固定年金和IG债券基金两个渠道约有5000亿美元资金涌入信用市场，供给克制与需求爆发正在改写利差定价方程。
  \item AAA级证券化利差预计将进一步收窄，非代理RMBS和CLO的净供给稀缺性最为突出。
\end{itemize}
\begin{figure}[htbp]
\centering
\includegraphics[width=0.88\linewidth]{figures/fig_001.jpg}
\caption{Figure 1 - Despite significant interest for Kevin Warsh's first meeting, we do not expect any big surprises this week. Market expectations already lean towards a hawkish statement and SEP: (1) the statement is likely to drop the easing bias; (2) the core PCE forecast sho}
\end{figure}
\begin{figure}[htbp]
\centering
\includegraphics[width=0.88\linewidth]{figures/fig_002.jpg}
\caption{Figure 2 - © 2026 Citi Inc. No redistribution without Citi's written permission. Figure 2. DXY is close to its short-term fair value}
\end{figure}
\begin{figure}[htbp]
\centering
\includegraphics[width=0.88\linewidth]{figures/fig_004.jpg}
\caption{Figure 4 - © 2026 Citi Inc. No redistribution without Citi's written permission. Figure 4. DXY range at the base of the uptrend continues}
\end{figure}
{\small\color{refgray}\textbf{References:} [R001] Citi：市场对美联储的“鹰派”定价已足够充分，真正的交易机会在于美元走弱时的战术性做空; [R006] DB：美国债券市场的真正变化不在需求端，而在持有者结构的不可逆重置; [R023] DB：市场真正低估的不是风险，而是结构性资金洪流对资产定价的重塑}

\section{AI/算力：智能体时代CPU复兴，中国供应链份额跃迁}
\textbf{AI从聊天机器人转向智能体范式，CPU与GPU配比从1:8回归1:1，中国AI基础设施供应商在下一代GPU架构中将实现结构性份额跃升。}

\begin{itemize}
  \item Bernstein将2030年服务器CPU市场TAM从此前1370亿美元大幅上调至2230亿美元，牛市情景下可达3300亿美元。
  \item 智能体AI改变了CPU与GPU配比逻辑：传统推理集群中GPU与CPU配比曾达8:1，但智能体工作流需要CPU承担任务调度、工具调用和内存管理。
  \item Bernstein预计CPU在AI工作负载中的价值占比将从传统LLM的约14\%升至50\%，推动CPU从配角回归核心地位。
  \item Arm被看作结构性受益者，其架构在功耗效率上有优势，且正从只卖IP转向自己做CPU，预计Arm AGI CPU销售额FY31达220亿美元。
  \item Citi调研显示中国AI基础设施供应商在英伟达GB200中份额有限，但在GB300和Vera Rubin中将实现显著跃迁，驱动力是技术曲线升级和响应速度。
  \item 上游材料是最大瓶颈也是最大机会：AI玻纤布和电子玻纤布产能最为紧张，某CCL龙头AI-CCL月产能占比从约10\%升至约15\%。
  \item 数据中心互联正从附加功能演变为高利润产品，毛利率高达70\%-90\%，Equinix以222个设施和51.3万条创收连接遥遥领先。
  \item 推理成本正成为AI商业化的第一约束，Revolut通过模型路由实现8倍成本节约，ElevenLabs在固定硬件上实现70倍吞吐提升。
\end{itemize}
\begin{figure}[htbp]
\centering
\includegraphics[width=0.88\linewidth]{figures/fig_018.jpg}
\caption{Exhibit 3 - EXHIBIT 2: Arm argues that agentic AI shifts more work back to the CPU: accelerators generate tokens, but CPUs orchestrate the agents, memory and workflows needed to deliver answers, making Arm's efficient CPU architecture increasi}
\end{figure}
\begin{figure}[htbp]
\centering
\includegraphics[width=0.88\linewidth]{figures/fig_020.jpg}
\caption{EXHIBIT 4 - EXHIBIT 4: Agentic AI shifts compute balance toward CPUs, with CPU share rising from \textbackslash{}\textasciitilde{}14\% in Traditional LLMs to 50\%, highlighting CPUs' growing orchestration role alongside GPUs in AI workloads at scale CPU:GPU ratio shift in Age}
\end{figure}
\begin{figure}[htbp]
\centering
\includegraphics[width=0.88\linewidth]{figures/fig_022.jpg}
\caption{EXHIBIT 6 - EXHIBIT 6: We expect the server CPU market to grow from US\textbackslash{}\$37bn in 2025 to US\textbackslash{}\$223bn in 2030, accelerating at a 43\% CAGR, driven by Agentic AI adoption.}
\end{figure}
\begin{figure}[htbp]
\centering
\includegraphics[width=0.88\linewidth]{figures/fig_029.jpg}
\caption{Exhibit 16 - EXHIBIT 16: We estimate a 2030 server CPU TAM of \textbackslash{}\$223bn, of which we believe Arm will be around \textbackslash{}\$123bn. EXHIBIT 17: We estimate Nvidia-related Arm CPU units to grow by 5x by 2030.}
\end{figure}
\begin{figure}[htbp]
\centering
\includegraphics[width=0.88\linewidth]{figures/fig_173.jpg}
\caption{Figure 1 - \#\# Upstream CCL and AI-fabric Material Greatest capacity shortages are in Al-fabric and e-glass fabric in CCL / PCB supply chain, as cited by Shengyi Tech and Grace Fabric Shengyi Tech and Grace Fabric both cited the greatest shortage in terms of capacity bein}
\end{figure}
\begin{figure}[htbp]
\centering
\includegraphics[width=0.88\linewidth]{figures/fig_210.jpg}
\caption{EXHIBIT 1 - EXHIBIT 1: Global Registered Network Connections by Provider (Units) Global Registered Network Connections by Provider (Units)}
\end{figure}
\begin{figure}[htbp]
\centering
\includegraphics[width=0.88\linewidth]{figures/fig_220.jpg}
\caption{EXHIBIT 13 - EXHIBIT 13: EQIX Interconnection Revenue Overview (\textbackslash{}\$M, \%) EQIX Interconnection Revenue Overview (\$M, \%)}
\end{figure}
{\small\color{refgray}\textbf{References:} [R003] Bernstein：AI从聊天走向智能体，CPU正在经历一场被低估的复兴; [R016] Citi：市场低估了中国AI供应链在下一代GPU架构中的份额跃迁; [R021] Bernstein：数据中心互联的“网络效应”才是真正的护城河; [R032] Citi：AI产业真正的分水岭不是模型能力，而是推理成本的精细化控制}

\section{能源与大宗：库存重建锚定油价，铜关税赌局6月底见分晓}
\textbf{全球石油市场消耗约10亿桶库存，补库需求将吸收未来两年供给增量，油价底部坚实；铜价受美国关税政策扭曲，6月30日截止日决定方向。}

\begin{itemize}
  \item Bernstein估算全球石油市场已累计消耗约10亿桶库存，包括约3亿桶SPR、约4.5亿桶中国库存和约1.4亿桶海上浮仓。
  \item 即便霍尔木兹海峡恢复通航，第一批恢复的油流将首先进入库存重建通道，而非直接满足终端需求，未来两年每天需额外吸收约300万桶。
  \item GS将2026Q4布伦特预测从90美元下调至80美元，2027年均价从80美元下调至75美元，但强调安全溢价的结构性重置使油价底部更坚实。
  \item GS预计2027年全球石油市场将出现每天320万桶过剩，但布伦特仍维持在75美元长期公允价值，市场愿为地缘风险支付结构性溢价。
  \item 铜价受美国关税预期扭曲：Comex库存自2025年4月以来持续攀升，交易员押注2027年1月开征15\%关税。
  \item Bernstein给出两条路径：基准情景无关税则铜价回落至11000美元/吨，看涨情景若延期或加征15\%+30\%则铜价可能重返历史高点。
  \item 对铜价敏感度最高的公司是FCX和ANTO（beta约1.5x和1.4x），BHP相对温和（0.7x）。
  \item 全球除卡塔尔外LNG供给增长超预期，GS将2026年6月预测上调10百万吨/年至384百万吨/年，欧洲TTF气价2027年预计降至30欧元/兆瓦时。
\end{itemize}
\begin{figure}[htbp]
\centering
\includegraphics[width=0.88\linewidth]{figures/fig_041.jpg}
\caption{EXHIBIT 2 - EXHIBIT 2: Oil prices to remain supported near \textbackslash{}\textasciitilde{}\textbackslash{}\$90/bbl in 2026 before easing to \textbackslash{}\textasciitilde{}\textbackslash{}\$78/bbl as supply normalises and inventories rebuild}
\end{figure}
\begin{figure}[htbp]
\centering
\includegraphics[width=0.88\linewidth]{figures/fig_058.jpg}
\caption{EXHIBIT 22 - EXHIBIT 22: Oil prices to average \textbackslash{}\textasciitilde{}\textbackslash{}\$90/bbl in 2026 before easing to \textbackslash{}\textasciitilde{}\textbackslash{}\$78/bbl in 2027 as inventories rebuild}
\end{figure}
\begin{figure}[htbp]
\centering
\includegraphics[width=0.88\linewidth]{figures/fig_131.jpg}
\caption{Exhibit 1 - Exhibit 1: We Are Reducing Our 2026Q4 Brent Forecast to \textbackslash{}\$80 and to \textbackslash{}\$75 for the 2027 Average as We Now Assume an Earlier Recovery in Mideast Supply Following a Deal to Reopen Hormuz}
\end{figure}
\begin{figure}[htbp]
\centering
\includegraphics[width=0.88\linewidth]{figures/fig_139.jpg}
\caption{Exhibit 9 - Exhibit 9: We See Risks to Our Brent Price Forecast as Two-Sided but Tilted to the Upside on Net}
\end{figure}
\begin{figure}[htbp]
\centering
\includegraphics[width=0.88\linewidth]{figures/fig_005.jpg}
\caption{EXHIBIT 1 - EXHIBIT 1: Net COT positioning had climbed significantly in late Q4 2025, helping copper price to reach \textbackslash{}\$13,000/t in December. Copper Prices vs. COT Investment Funds Positioning}
\end{figure}
\begin{figure}[htbp]
\centering
\includegraphics[width=0.88\linewidth]{figures/fig_006.jpg}
\caption{EXHIBIT 2 - EXHIBIT 2: Comex copper inventories have steadily risen since April 2025 with an average of c.9kt incremental inventories per week. Comex Copper Inventories (short ton) vs. Weekly Change}
\end{figure}
\begin{figure}[htbp]
\centering
\includegraphics[width=0.88\linewidth]{figures/fig_169.jpg}
\caption{Exhibit 1 - Exhibit 1: Low European gas storage and an ongoing Asia pull for LNG cargoes at the expense of Europe would likely lead to a further rise in TTF without a Hormuz flow normalization NW European Storage (Ihs) vs Global Ex-Qatar LNG L}
\end{figure}
\begin{figure}[htbp]
\centering
\includegraphics[width=0.88\linewidth]{figures/fig_171.jpg}
\caption{Exhibit 3 - Exhibit 3: We expect European gas prices to balance the market in the low 40 EURs/MWh for the remainder of summer Realized and forward TTF prices and fuel switching ranges, and TTF GS fcast}
\end{figure}
{\small\color{refgray}\textbf{References:} [R004] Bernstein：市场低估的不是供给恢复的速度，而是库存重建的深度; [R010] GS：市场低估的不是供需平衡，而是安全溢价的结构性重置; [R002] 0002-02-Bernstein-Global Metals - Mining- Copper Taco or Boco Season Finale - U.S. policy is distorting copper price but; [R014] GS：市场低估了全球LNG供给的结构性改善，而非地缘政治风险}

\section{地产与消费：一线城市结构性再定价，消费疲软拉长政策等待期}
\textbf{中国房地产市场并非全面回暖，一线城市二手房连续3-4个月环比正增长，但全国总量疲软；5月社零同比转负，消费情绪短期难逆转。}

\begin{itemize}
  \item GS报告显示5月一线城市新房价格环比上涨0.2\%，连续第4个月正增长；二手房价格环比上涨0.3\%，连续第3个月正增长。
  \item 深圳和上海几乎所有产品类型二手房成交价格都在改善，北京高端项目跑赢大盘，但驱动力高度集中且依赖局部政策。
  \item 摩根斯坦利指出新房成交的脉冲式反弹正在退潮，50城新房成交同比从+21\%骤降至-1\%，而二手房成交韧性更强，定价权正从新房转向二手房。
  \item 5月社零同比下跌0.6\%，为2022年底以来首次月度负增长，远低于市场预期的微增0.1\%，餐饮增速放缓至0.6\%。
  \item NOM认为消费疲软正在从个别品类扩散为系统性现象，房地产回暖带来的财富效应和AI行情对消费信心的支撑均未在数据中验证。
  \item GS判断5月经济数据低于预期，但政策窗口正在打开，7月政治局会议是潜在政策调整的战略节点，三季度财政发力和降息概率上升。
  \item NOM指出5月数据的核心特征是需求弱、供给更紧，石油供应冲击对工业生产的传导比预期更深更持久，原油加工量同比下滑9.1\%。
  \item 白酒端午动销预期转弱，经销商预期从20\%-30\%增长下修至10\%-20\%，飞天茅台批价在1635-1670元窄幅波动，渠道库存周期正在钝化。
\end{itemize}
\begin{figure}[htbp]
\centering
\includegraphics[width=0.88\linewidth]{figures/fig_184.jpg}
\caption{Exhibit 1 - Exhibit 2: May-26 70-city ASP index was -0.2\%/-0.3\% mom (vs. -0.2\%/-0.2\% mom for Apr), with 16/10 cities (vs. 14/12 in Apr) recording sequentially improved ASPs and 2/3 cities (vs. 7/4 in Apr) recording sequentially flattened ASPs}
\end{figure}
\begin{figure}[htbp]
\centering
\includegraphics[width=0.88\linewidth]{figures/fig_191.jpg}
\caption{Exhibit 9 - Exhibit 9: 15 major cities' secondary sales volume on average fell \$12\textbackslash{}\%\$ mom and was \$+8\textbackslash{}\%\$ yoy in May-26, and recorded \$-3\textbackslash{}\%\$ yoy in 5M26 15 major cities' secondary market transaction volume tracker 15 major cities secondary sale}
\end{figure}
\begin{figure}[htbp]
\centering
\includegraphics[width=0.88\linewidth]{figures/fig_206.jpg}
\caption{Exhibit 1 - Exhibit 1: In the primary market, the sequential decline in 70-city weighted average property price eased to 1.8\% mom annualized in May}
\end{figure}
\begin{figure}[htbp]
\centering
\includegraphics[width=0.88\linewidth]{figures/fig_209.jpg}
\caption{Exhibit 1 - Exhibit 1: Primary weekly unit sales and four-week moving average – Total 50 cities}
\end{figure}
\begin{figure}[htbp]
\centering
\includegraphics[width=0.88\linewidth]{figures/fig_164.jpg}
\caption{Exhibit 1 - Exhibit 1: Industrial production rose in month-on-month terms in May after seasonal adjustment}
\end{figure}
\begin{figure}[htbp]
\centering
\includegraphics[width=0.88\linewidth]{figures/fig_166.jpg}
\caption{Exhibit 3 - Exhibit 3: FAI growth declined further in May}
\end{figure}
{\small\color{refgray}\textbf{References:} [R017] GS：中国房地产市场的真实信号并非全面回暖，而是核心城市的结构性再定价; [R019] GS：五月新房价格跌幅收窄，但市场真正的信号藏在二手房与库存之间; [R020] 摩根斯坦利：中国房地产市场真正的拐点信号，不是成交量的反弹，而是成交结构的断裂; [R025] NOM：消费数据跌破疫后低点，市场正在低估“无政策刺激”下的结构性分化; [R013] GS：5月经济数据低于预期，但政策窗口正在打开; [R015] NOM：五月经济数据印证了一个被市场低估的核心矛盾——供给侧的约束比需求侧更顽固; [R012] GS：白酒端午动销预期转弱，市场低估了渠道库存周期的“钝化”}

\section{区域市场：日本央行鸽派加息，银行股定价偏保守}
\textbf{日本央行6月加息25bp至1.0\%，但鸽派委员反对票与前瞻指引变化暗示加息节奏可能慢于市场预期，银行股当前股价隐含的利率假设偏低。}

\begin{itemize}
  \item DB指出6月会议以7:1通过加息，唯一反对票来自3月上任的浅田委员，就职不到三个月公开反对行长提案在日本央行历史上极为罕见。
  \item 另一位新委员佐藤将于6月加入委员会，若也采取鸽派立场，委员会投票格局将从鹰派主导转向鸽派力量显著增强。
  \item 前瞻指引从实际利率处于极低水平改为金融环境一直宽松，暗示央行不打算连续激进加息，也不打算将利率推入紧缩区间。
  \item GS认为市场对日本银行股的定价仍然偏保守，当前股价隐含的政策利率假设约在0.75\%-1.0\%，而央行传递的加息周期延续信号需要重新校准。
  \item 主要银行存款贝塔系数维持在40\%，意味着每100个基点加息，银行只需传导40个基点给存款人，剩余60个基点转化为净息差扩张。
  \item 横滨金融集团率先披露加息对下财年利润正面影响，预计税后利润增加约100亿日元，ROE提升0.7\%，其他银行大概率跟进。
  \item 日本银行股仓位较轻，参考去年12月加息后走势，加息后往往迎来估值修复或空头回补，P/B仍有上行空间。
\end{itemize}
\begin{figure}[htbp]
\centering
\includegraphics[width=0.88\linewidth]{figures/fig_060.jpg}
\caption{Figure 1 - Figure 1: Outlook for outstanding of fund-provisioning measure to stimulate bank lending}
\end{figure}
\begin{figure}[htbp]
\centering
\includegraphics[width=0.88\linewidth]{figures/fig_061.jpg}
\caption{Figure 2 - Figure 2: Forecasts for BoJ's JGB holding Gross purchase}
\end{figure}
\begin{figure}[htbp]
\centering
\includegraphics[width=0.88\linewidth]{figures/fig_291.jpg}
\caption{Exhibit 2 - Exhibit 2: Banks' product rate hikes reflecting BOJ policy rate hike Exhibit 3: Topix banks' share price performance in 30 days before and after previous rate hikes}
\end{figure}
{\small\color{refgray}\textbf{References:} [R005] DB：日本央行鸽派转向比加息本身更值得关注; [R007] GS：日本央行加息至1\%，市场真正低估的不是终点利率，而是加息节奏的重新定价; [R035] GS：日本银行股正在定价的利率路径，可能比市场意识到的更陡}

\section{新能源车：渗透率突破67\%，竞争逻辑从抢增量转向抢存量}
\textbf{新能源车周订单回落是新品脉冲退潮的正常表现，渗透率突破67\%才是真正信号，行业竞争逻辑正从抢增量转向抢存量。}

\begin{itemize}
  \item GS追踪的9家主要新能源品牌6月第二周合计订单15.36万辆，环比下降13\%，同比下降14\%，但回落是5月底新品上市脉冲退潮的正常节奏。
  \item 6月第一周新能源零售渗透率66.7\%，批发渗透率67.2\%，相比5月的63\%/61.1\%明显提升，燃油车市场份额正被不可逆侵蚀。
  \item 吉利、比亚迪、特斯拉相对抗跌，周订单环比分别+8\%/-4\%/-8\%，蔚来YTD订单同比+93\%表现亮眼。
  \item 新能源经销商折扣小幅收窄至7.54\%，燃油车折扣继续扩大至19.56\%，燃油车压力不减。
  \item 电池级碳酸锂涨至17.45万/吨（环比+3.6\%），但电芯价格稳定，成本传导暂未发生。
  \item 进口车市场5月零售销量同比下滑38\%，环比零增长，雷克萨斯和奔驰相对韧性，品牌护城河正从豪华认知转向供需纪律。
\end{itemize}
\begin{figure}[htbp]
\centering
\includegraphics[width=0.88\linewidth]{figures/fig_204.jpg}
\caption{Exhibit 2 - Exhibit 2: NEV dealer discount trend}
\end{figure}
\begin{figure}[htbp]
\centering
\includegraphics[width=0.88\linewidth]{figures/fig_205.jpg}
\caption{Exhibit 3 - Exhibit 3: ICE dealer discount trend}
\end{figure}
\begin{figure}[htbp]
\centering
\includegraphics[width=0.88\linewidth]{figures/fig_234.jpg}
\caption{Figure 1 - Figure 1. Imported Sales by Brand in China (units) © 2026 Citi Inc. No redistribution without Citi's written permission. Figure 2. May-26 and 5M26 YoY Growth of Imported Sales by Brand}
\end{figure}
{\small\color{refgray}\textbf{References:} [R018] GS：新能源车周订单回落，但渗透率突破67\%才是市场真正该关注的信号; [R027] Citi：进口车市场正在经历的不是周期性波动，而是品牌溢价能力的根本性分化}

\section{资本市场：资本支出超级周期重塑资产定价}
\textbf{过去四十年的现代周期已不可逆终结，资本支出超级周期正在将回报驱动力从估值扩张转向盈利增长，选股能力变得更为关键。}

\begin{itemize}
  \item GS将1982年至疫情前定义为现代周期，其核心特征低通胀、低利率、低波动、全球化、低资本密集度和高利润率正在全面逆转。
  \item 支撑现代周期的所有支柱都在瓦解：从通缩转向再通胀，从零利率转向高资金成本，从全球化转向区域化，从低资本密集度转向大规模基础设施投资。
  \item AI革命带动私营部门投资激增，能源安全和地缘政治催生公共投资，供应链从效率优先转向韧性优先。
  \item 市场回报将从估值驱动转向盈利增长驱动，资产之间收益差异拉大，传统买指数躺赢策略可能失效。
  \item AI主题的逢低买入策略回测显示显著超额收益，其深层含义是AI基础设施供给端存在结构性瓶颈，每一次情绪回调都被供需基本面修正所救赎。
  \item 摩根斯坦利欧洲行业模型显示半导体排第一，金属矿业升至第二，银行从第六跃升至第三，AI、银行和矿业被加倍下注。
  \item 国防产业正经历去空心化的结构性重塑，中型国防企业崛起，风险资本涌入，估值逻辑越来越像成长股。
\end{itemize}
\begin{figure}[htbp]
\centering
\includegraphics[width=0.88\linewidth]{figures/fig_252.jpg}
\caption{Exhibit 1 - Exhibit 1: Corporate tax rates fell from 1980s to 2020 Effective corporate tax rate. Historical constituents kept constant before year: 1989 (S\&P 500), 1996 (FTSE 350), 2000 (SBF 120 and CDAX)}
\end{figure}
\begin{figure}[htbp]
\centering
\includegraphics[width=0.88\linewidth]{figures/fig_253.jpg}
\caption{Exhibit 2 - Exhibit 2: The ‘Modern’ super cycle was characterised by lower volatility in macro variables Data for the US}
\end{figure}
\begin{figure}[htbp]
\centering
\includegraphics[width=0.88\linewidth]{figures/fig_269.jpg}
\caption{Exhibit 18 - Exhibit 18: The emergence of large language models has spawned a renewed requirement for capital spending Capex as a \% of cash flow from operations, dashed line represents 2026 consensus}
\end{figure}
\begin{figure}[htbp]
\centering
\includegraphics[width=0.88\linewidth]{figures/fig_273.jpg}
\caption{Exhibit 22 - Exhibit 22: AI hyperscalers are expected to spend c.\textbackslash{}\$755 billion on capex in 2026 Hyperscaler annual capex (\textbackslash{}\$bn)}
\end{figure}
\begin{figure}[htbp]
\centering
\includegraphics[width=0.88\linewidth]{figures/fig_113.jpg}
\caption{Exhibit 1 - Exhibit 1: Buy-the-dip strategy cumulative performance since 2023 - "AI Infrastructure" sub-theme}
\end{figure}
\begin{figure}[htbp]
\centering
\includegraphics[width=0.88\linewidth]{figures/fig_118.jpg}
\caption{Exhibit 6 - Exhibit 6: Semis stand out on fundamental momentum indicators; Banks jump back into the top 3}
\end{figure}
{\small\color{refgray}\textbf{References:} [R031] GS：市场低估了资本支出超级周期对资产定价的深远重塑; [R009] 摩根斯坦利：市场真正低估的不是AI泡沫，而是供给侧的定价权转移; [R036] 摩根斯坦利：国防产业正经历一场“去空心化”的结构性重塑，而非简单的周期上行}

\section{人形机器人：7.5万亿市场今年启动，供应链再定价}
\textbf{人形机器人正从技术可行性验证进入小规模商业化阶段，2026年是转折点，供应链从精密零部件到边缘计算芯片将出现价值重估。}

\begin{itemize}
  \item 摩根斯坦利预测到2050年全球人形机器人市场规模达7.5万亿美元，保有量超过10亿台。
  \item 小规模商业化已在2026年开始：前4个月全球VC融资超过2025年全年，Figure估值390亿美元，多家公司已提交IPO申请。
  \item 到2040年核心零部件市场总规模预计达7800亿美元，其中电机2073亿、边缘计算1233亿、减速器1370亿、摄像头316亿。
  \item 中国供应链在零部件环节比整机更有吸引力，某精密零部件公司已向美国头部机器人公司批量出货。
  \item Optimus Gen3预计今年产量1-2万台，明年可能10倍增长到10万+。
  \item 工业场景（搬运、质检、装配）是率先落地的应用领域，家庭服务场景商业化路径更长。
\end{itemize}
\begin{figure}[htbp]
\centering
\includegraphics[width=0.88\linewidth]{figures/fig_178.jpg}
\caption{Figure 8 - In China, UBTECH (9880.HK) will soon debut several new robot products, including bionic robots, industrial humanoid robot Walker S3, commercial humanoid robot new Walker C, and wheel-based Cruzr Y1 to drive its top line growth. Mgmt. guided that the total robo}
\end{figure}
{\small\color{refgray}\textbf{References:} [R024] 摩根斯坦利：人形机器人市场真正被低估的不是需求，而是供应链的再定价}

\section{锂电与焦煤：库存认知分裂，供应冲击暴露结构性脆弱}
\textbf{锂电产业链上下游对库存水平出现方向性分歧，转化企业悲观情绪蔓延；中国焦煤供应冲击的真正影响不在中国，而在印度钢企利润结构。}

\begin{itemize}
  \item JEF调研显示转化企业预计未来三个月锂盐价格下跌超10\%，而电池企业预期仅为个位数跌幅，上下游价格预期差距在扩大。
  \item 转化企业认为电池厂商库存过低，但电池厂商认为自身客户库存适度过剩，库存认知分裂正在重塑产业链议价机制。
  \item 转化环节产能过剩更为严重，产品同质化程度高，转化企业承受比下游更重的价格压力。
  \item NOM指出山西煤矿事故后137座焦煤矿停产，截至6月11日仅77座恢复，恢复进程不均衡，部分矿井复产几天后再次停产。
  \item 国内焦煤现货价已涨超10\%至每吨270美元，市场定价的不仅是单一事件，而是国内钢铁原料供应整体收紧的预期。
  \item 澳大利亚5月对华焦煤出口环比激增114\%，但蒙古和俄罗斯受物流和品质限制难以完全替代。
  \item 焦煤每涨10美元/吨，对印度综合钢企EBITDA影响约7-9美元/吨，JSW Steel和Tata Steel暴露度最高。
\end{itemize}
\begin{figure}[htbp]
\centering
\includegraphics[width=0.88\linewidth]{figures/fig_030.jpg}
\caption{Exhibit 18 - For hyperscaler in-house CPUs, such as AWS Graviton and Google Axion, we model it as a function of server unit assumptions as well as ODM-direct server exposures, and applying Arm penetration within the CSP server base. In our model, we assume Arm penetration }
\end{figure}
{\small\color{refgray}\textbf{References:} [R033] JEF：锂电产业链的真实信号不是价格下跌，而是库存认知的分裂; [R030] NOM：中国焦煤供应冲击的真正影响，不在中国，而在印度钢厂的利润结构}

\section{油运与能源基建：供给侧弹性被低估，补库周期是真正定价锚}
\textbf{油轮运输板块过去三周跑输大盘17-27个百分点，但市场对需求萎缩的担忧可能方向搞反，海峡重开后的补库周期将带来持续性运价支撑。}

\begin{itemize}
  \item 摩根斯坦利指出市场抛售油轮股的逻辑链条是海峡关闭导致运输需求下降，但忽略了被压制的需求可能在短时间内集中释放。
  \item 海峡关闭期间全球原油库存持续消耗，进口国补库意愿被压制，一旦重开，进口国重建库存和出口国清空储油两股力量叠加。
  \item 供给侧真实弹性远低于市场假设：波斯湾内虽仍有油轮但多为满载状态，短期内放不出运力。
  \item 招商轮船预计二季度净利润42亿，环比增长超50\%，主要受益于3月VLCC即期运价大涨和干散货运价走强。
  \item 能源基础设施的定价权正从宏观叙事转向微观资产：霍尔木兹重开可能成为资金重新回流能源板块的转折点。
  \item AI/数据中心股票回调连带天然气基建股走弱，但摩根斯坦利认为超跌过度，下半年仍可能有新项目公告。
  \item KMI的Gulf Coast Express管道扩容已投运，推动Waha天然气价格反弹，但新产能可能很快填满。
\end{itemize}
\begin{figure}[htbp]
\centering
\includegraphics[width=0.88\linewidth]{figures/fig_292.jpg}
\caption{Exhibit 1 - Exhibit 1: Relative share performance: Tankers}
\end{figure}
\begin{figure}[htbp]
\centering
\includegraphics[width=0.88\linewidth]{figures/fig_293.jpg}
\caption{Exhibit 2 - Exhibit 2: Tanker Spot Earnings have corrected mildly since May 2026}
\end{figure}
\begin{figure}[htbp]
\centering
\includegraphics[width=0.88\linewidth]{figures/fig_296.jpg}
\caption{Exhibit 5 - Exhibit 5: 20\% of crude tankers are "sanctioned", 24\% of global tankers are "non-compliant"}
\end{figure}
\begin{figure}[htbp]
\centering
\includegraphics[width=0.88\linewidth]{figures/fig_235.jpg}
\caption{Exhibit 2 - Exhibit 2: Hyperscaler AI Capex}
\end{figure}
\begin{figure}[htbp]
\centering
\includegraphics[width=0.88\linewidth]{figures/fig_238.jpg}
\caption{Exhibit 5 - Exhibit 5: Waha Natural Gas Prices (\textbackslash{}\$/Mcf)}
\end{figure}
{\small\color{refgray}\textbf{References:} [R037] 摩根斯坦利：市场低估的不是油运需求，而是供给侧的再定价; [R029] 摩根斯坦利：能源基础设施的定价权正在从宏观叙事转向微观资产}

\section{银行与信贷：社融回暖但需求根基未稳，选股逻辑转向定价权}
\textbf{5月社融回暖但结构揭示真实实体融资意愿未修复，银行股选股逻辑正从规模扩张转向定价权和资产质量优势。}

\begin{itemize}
  \item DB数据显示5月新增社融2.0万亿，人民币贷款5200亿，较4月低点回升，但存量社融和贷款同比增速分别放缓至7.7\%和5.5\%历史新低。
  \item 政府债券仍是社融绝对主力，企业中长期贷款连续为负（-200亿），居民贷款延续收缩（-1410亿），真实融资需求疲弱。
  \item 存款结构变化显示资金在往理财和资本市场迁移，居民和企业存款减少，非银金融机构存款增加1.1万亿。
  \item DB偏好具有较大企业业务敞口、经营韧性较强、股息率合理的大型银行，首选为建设银行H股和中国银行H股。
  \item M1增速5.5\%（前值5.0\%），M2增速8.6\%（前值8.4\%），剪刀差略有收窄但仍在低位。
\end{itemize}
\begin{figure}[htbp]
\centering
\includegraphics[width=0.88\linewidth]{figures/fig_227.jpg}
\caption{Figure 1 - Figure 1: New increased TSF - May 2026}
\end{figure}
\begin{figure}[htbp]
\centering
\includegraphics[width=0.88\linewidth]{figures/fig_228.jpg}
\caption{Figure 2 - Figure 2: New increased RMB loans - May 2026}
\end{figure}
\begin{figure}[htbp]
\centering
\includegraphics[width=0.88\linewidth]{figures/fig_230.jpg}
\caption{Figure 4 - Figure 4: Money supply vs TSF (\%YoY growth)}
\end{figure}
{\small\color{refgray}\textbf{References:} [R022] DB：五月社融回暖但需求根基未稳，银行股的选股逻辑正在从规模转向定价权}

\section{结语}
综合来看，市场正处于宏观叙事切换的关键窗口：地缘风险溢价正在被库存重建和结构性资金洪流所取代，AI从训练走向推理和智能体时代正在重塑产业链价值分配，而中国经济的结构性分化要求投资者从总量思维转向精细化选股。
\newpage
\section{报告覆盖清单}
\begin{itemize}
  \item [R001] 宏观与利率：美联储鹰派定价充分，美元区间待突破 - Citi：市场对美联储的“鹰派”定价已足够充分，真正的交易机会在于美元走弱时的战术性做空
  \item [R002] 能源与大宗：库存重建锚定油价，铜关税赌局6月底见分晓 - 0002-02-Bernstein-Global Metals - Mining- Copper Taco or Boco Season Finale - U.S. policy is distorting copper price but
  \item [R003] AI/算力：智能体时代CPU复兴，中国供应链份额跃迁 - Bernstein：AI从聊天走向智能体，CPU正在经历一场被低估的复兴
  \item [R004] 能源与大宗：库存重建锚定油价，铜关税赌局6月底见分晓 - Bernstein：市场低估的不是供给恢复的速度，而是库存重建的深度
  \item [R005] 区域市场：日本央行鸽派加息，银行股定价偏保守 - DB：日本央行鸽派转向比加息本身更值得关注
  \item [R006] 宏观与利率：美联储鹰派定价充分，美元区间待突破 - DB：美国债券市场的真正变化不在需求端，而在持有者结构的不可逆重置
  \item [R007] 区域市场：日本央行鸽派加息，银行股定价偏保守 - GS：日本央行加息至1\%，市场真正低估的不是终点利率，而是加息节奏的重新定价
  \item [R008] 未被主题摘要引用 - JPM：出口仍是唯一结构性亮点，但国内需求疲软正在拉长政策等待期
  \item [R009] 资本市场：资本支出超级周期重塑资产定价 - 摩根斯坦利：市场真正低估的不是AI泡沫，而是供给侧的定价权转移
  \item [R010] 能源与大宗：库存重建锚定油价，铜关税赌局6月底见分晓 - GS：市场低估的不是供需平衡，而是安全溢价的结构性重置
  \item [R011] 未被主题摘要引用 - GS：市场低估的不是电商增速放缓，而是云计算在AI需求下的重新定价
  \item [R012] 地产与消费：一线城市结构性再定价，消费疲软拉长政策等待期 - GS：白酒端午动销预期转弱，市场低估了渠道库存周期的“钝化”
  \item [R013] 地产与消费：一线城市结构性再定价，消费疲软拉长政策等待期 - GS：5月经济数据低于预期，但政策窗口正在打开
  \item [R014] 能源与大宗：库存重建锚定油价，铜关税赌局6月底见分晓 - GS：市场低估了全球LNG供给的结构性改善，而非地缘政治风险
  \item [R015] 地产与消费：一线城市结构性再定价，消费疲软拉长政策等待期 - NOM：五月经济数据印证了一个被市场低估的核心矛盾——供给侧的约束比需求侧更顽固
  \item [R016] AI/算力：智能体时代CPU复兴，中国供应链份额跃迁 - Citi：市场低估了中国AI供应链在下一代GPU架构中的份额跃迁
  \item [R017] 地产与消费：一线城市结构性再定价，消费疲软拉长政策等待期 - GS：中国房地产市场的真实信号并非全面回暖，而是核心城市的结构性再定价
  \item [R018] 新能源车：渗透率突破67\%，竞争逻辑从抢增量转向抢存量 - GS：新能源车周订单回落，但渗透率突破67\%才是市场真正该关注的信号
  \item [R019] 地产与消费：一线城市结构性再定价，消费疲软拉长政策等待期 - GS：五月新房价格跌幅收窄，但市场真正的信号藏在二手房与库存之间
  \item [R020] 地产与消费：一线城市结构性再定价，消费疲软拉长政策等待期 - 摩根斯坦利：中国房地产市场真正的拐点信号，不是成交量的反弹，而是成交结构的断裂
  \item [R021] AI/算力：智能体时代CPU复兴，中国供应链份额跃迁 - Bernstein：数据中心互联的“网络效应”才是真正的护城河
  \item [R022] 银行与信贷：社融回暖但需求根基未稳，选股逻辑转向定价权 - DB：五月社融回暖但需求根基未稳，银行股的选股逻辑正在从规模转向定价权
  \item [R023] 宏观与利率：美联储鹰派定价充分，美元区间待突破 - DB：市场真正低估的不是风险，而是结构性资金洪流对资产定价的重塑
  \item [R024] 人形机器人：7.5万亿市场今年启动，供应链再定价 - 摩根斯坦利：人形机器人市场真正被低估的不是需求，而是供应链的再定价
  \item [R025] 地产与消费：一线城市结构性再定价，消费疲软拉长政策等待期 - NOM：消费数据跌破疫后低点，市场正在低估“无政策刺激”下的结构性分化
  \item [R026] 未被主题摘要引用 - NOM：电商渗透率回升背后，消费结构的分化比总量更值得关注
  \item [R027] 新能源车：渗透率突破67\%，竞争逻辑从抢增量转向抢存量 - Citi：进口车市场正在经历的不是周期性波动，而是品牌溢价能力的根本性分化
  \item [R028] 未被主题摘要引用 - 摩根斯坦利：市场对日本汽车股的悲观定价，忽略了三个尚未被充分理解的变量
  \item [R029] 油运与能源基建：供给侧弹性被低估，补库周期是真正定价锚 - 摩根斯坦利：能源基础设施的定价权正在从宏观叙事转向微观资产
  \item [R030] 锂电与焦煤：库存认知分裂，供应冲击暴露结构性脆弱 - NOM：中国焦煤供应冲击的真正影响，不在中国，而在印度钢厂的利润结构
  \item [R031] 资本市场：资本支出超级周期重塑资产定价 - GS：市场低估了资本支出超级周期对资产定价的深远重塑
  \item [R032] AI/算力：智能体时代CPU复兴，中国供应链份额跃迁 - Citi：AI产业真正的分水岭不是模型能力，而是推理成本的精细化控制
  \item [R033] 锂电与焦煤：库存认知分裂，供应冲击暴露结构性脆弱 - JEF：锂电产业链的真实信号不是价格下跌，而是库存认知的分裂
  \item [R034] 未被主题摘要引用 - Bernstein：液冷CDU的真正分化不在技术路线，而在推理与训练的两极定价
  \item [R035] 区域市场：日本央行鸽派加息，银行股定价偏保守 - GS：日本银行股正在定价的利率路径，可能比市场意识到的更陡
  \item [R036] 资本市场：资本支出超级周期重塑资产定价 - 摩根斯坦利：国防产业正经历一场“去空心化”的结构性重塑，而非简单的周期上行
  \item [R037] 油运与能源基建：供给侧弹性被低估，补库周期是真正定价锚 - 摩根斯坦利：市场低估的不是油运需求，而是供给侧的再定价
\end{itemize}
\section{逐篇报告摘录}
\subsection{[R001] Citi：市场对美联储的“鹰派”定价已足够充分，真正的交易机会在于美元走弱时的战术性做空}
{\small\color{refgray}归类：宏观与利率：美联储鹰派定价充分，美元区间待突破}

Citi：市场对美联储的“鹰派”定价已足够充分，真正的交易机会在于美元走弱时的战术性做空

这份Citi研报的标题直接点明了核心判断：Fade EUR rallies if Warsh disappoints hawks。翻译过来，就是“如果新任美联储主席沃什未能满足鹰派预期，则逢高做空欧元”。这看似是一个具体的交易建议，但其背后隐藏着一个对当前宏观交易格局更根本的判断：市场已经提前消化了美联储转向鹰派的大部分预期，美元指数长达12个月的区间震荡格局，其突破的触发点并不在本次会议本身，而在于沃什在新闻发布会上的措辞，以及更重要的——地缘政治风险的缓和。

这不再是一个简单的“加息与否”的二元博弈。报告揭示了一个更微妙的定价环境：美元的多头头寸已经大幅回撤，美元指数也已接近其短期公允价值。这意味着，即便美联储如市场所愿释放鹰派信号，其边际效用也在递减。真正的变量，是那些尚未被充分定价的“非对称风险”。

1. 市场已为“鹰派声明”充分定价，沃什的“模糊”才是大概率事件

报告明确指出，市场对于本次FOMC声明的鹰派调整已有充分预期：声明可能删除宽松倾向、核心PCE预测将被上调、点阵图中位数将指向2026年不再降息，甚至可能出现一两个加息点。这些变化，在Citi看来，不足以引发显著的外汇市场反应，因为它们已经“price in”了。

真正的悬念在于沃什的新闻发布会。作为新任主席，他的沟通风格和具体措辞将成为市场关注的焦点。Citi的基准判断是，沃什将延续其“减少前瞻指引”的一贯立场，对当前市场定价（即年内加息18个基点，2027年3月首次加息完全定价）给出模棱两可的回答。这种“模糊”本身，就是对当前鹰派定价的一种温和打压。

这个判断的核心含义在于：市场对“鹰派美联储”的叙事已经过度交易。如果沃什没有给出比市场预期更鹰派的惊喜，美元将缺乏向上突破的催化剂。对于投资者而言，这意味着需要警惕“利好出尽”式的美元回调，而不是追涨。

2. 真正可能打破美元区间震荡的是“沃什意外鹰派”，而非地缘政治缓和

报告并未将地缘政治风险（如伊朗协议）视为压垮美元的因素。恰恰相反，Citi认为，一个潜在的美伊协议，以及沃什可能对此的“倾斜”，反而可能成为市场忽略商品通胀冲击的额外理由，从而利好风险资产，不利于美元。

那么，什么才能真正打破美元的12个月区间？报告给出了一个清晰的“非基准情景”：如果沃什在新闻发布会上认可了市场定价，或者暗示加息和降息的可能性是“同等可能”，这将成为一个重大的美元上行催化剂。因为市场会预期加息提前，并可能为第二次加息预留溢价。这将推动美元指数突破其长期上升通道底部的盘整区间。

这个分析框架的洞察在于：地缘政治缓和带来的美元下行风险是“膝跳反射”式的，且可能被日本财务省的干预所放大（如USDJPY），但其持续性有限。而“沃什意外鹰派”带来的美元上行风险，则具有更强的趋势性和可交易性，因为这意味着整个利率定价中枢的上移。因此，当前交易的核心不是押注方向，而是管理这个非对称风险敞口。

3. 美元多头的“拥挤”已大幅缓解，这为战术性做空提供了安全垫

报告引用了两个关键图表来支持其判断：一是美元多头头寸在FOMC会议前已大幅回撤（Figure 1），二是美元指数接近其短期公允价值（Figure 2）。这两个数据点合在一起，揭示了一个重要的市场结构特征。

当美元多头头寸极度拥挤时，任何不及预期的鹰派信号都可能导致剧烈的踩踏式下跌。而现在，拥挤度已经显著下降，这意味着美元的下行风险更多是“有序回调”而非“崩盘”。同时，接近公允价值意味着美元的下行空间相对有限，特别是在 的欧元兑美元阻力位附近，风险回报比变得非常具有吸引力。

这个结论对投资者的指导意义是：当前不是构建长期美元空头头寸的时机，而是可以利用F

[... middle omitted ...]

向一个“低波动、利差交易回归”的环境切换。

这个判断的关键在于：它解释了Citi为什么建议用欧元和加元作为利差交易的融资货币。在一个风险偏好改善、波动率下降的环境里，借入低收益货币（如欧元）去投资高收益货币（如美元）的利差交易将重新变得有利可图。而欧元和加元，恰恰因为其各自的经济增长前景和收益率劣势，成为理想的融资货币。

5. 报告尚未完全回答的关键问题：日本财务省的干预阈值与美债收益率曲线的陡峭化

尽管报告逻辑清晰，但仍有几个关键问题未被完全解答，而这些正是决定交易成败的细节。

首先，关于日本财务省的干预。报告提到了USDJPY可能因美元膝跳反射式下跌而出现更大反应，如果日本财务省利用这个机会进行干预。但报告没有给出明确的干预触发阈值或条件。历史上，日本财务省的干预往往具有“出其不意”的特征，其具体操作逻辑和容忍区间，是市场难以精确量化的变量。这构成了一个重要的尾部风险：如果美元下跌，USDJPY的下行幅度可能被干预行为放大，从而影响整体利差交易组合的风险。

\subsection{[R002] 0002-02-Bernstein-Global Metals - Mining- Copper Taco or Boco Season Finale - U.S. policy is distorting copper price but}
{\small\color{refgray}归类：能源与大宗：库存重建锚定油价，铜关税赌局6月底见分晓}

DeepSeek 生成 WeChat article 失败：Response ended prematurely

Personal reading notes and learning share only. Not investment advice.

1️⃣ 矿山供应中断：Kamoa-Kakula和Grasberg等大矿接连出事，去年9月Freeport的Grasberg泥浆事故直接推铜价破\$10,000/t。

2️⃣ 金融投机加码：去年Q4期货净多头仓位飙升，铜价从\$11,000一路冲到\$13,000。

3️⃣ 美国囤货潮：交易员押注2027年1月开征15\%关税、2028年可能再加到30\%，提前往美国仓库塞货。Comex库存去年涨了344kt，今年又加了152kt。

4️⃣ 风险情绪回暖：美股创新高，LME与Comex价差扩大，铜价5月中旬摸到历史新高\$14,109/t。

📌 基准情景（无关税）：宣布精炼铜不征关税。美国库存可能回流全球，铜价逐步回落到\$11,000/t。

📌 看涨情景（三种可能）：延期决定 → 铜继续涨；只征15\%且无后续指引 → 更涨；15\%+30\%两步走 → 涨到新高。即便矿山不出事，铜价也可能重返历史高点。

风险回报比很诱人：从西欧运铜到美国的成本约\$400-754/t，但一旦征15\%关税，价差可能超过\$2,000/t。交易员有足够动力继续往Comex塞货。

对铜价敏感度最高的公司是FCX和ANTO（beta约1.5x和1.4x），其次是AAL（1.1x），BHP相对温和（0.7x）。

Global Metals \& Mining: Copper Taco or Boco Season Finale - U.S. policy is distorting copper price but what can we do about it?

We approach the deadline (June 30th) for the White House to decide on tariffs on refined copper. In previous notes, we expressed our TACO view of no tariffs. Trade deals were promised (Trump Directs Negotiations to Adjust Imports of Processed Critical Minerals) but haven't arrived from, say, Chile or Indonesia. In this note, we evaluate potential paths for copper prices. Bernstein Occasionally Chickens Out?

Copper prices have surged over the past year due to: (1) supply disruptions at major mines, (2) increased financial speculation (Exhibit 1), (3) additional demand from US inventory stockpiling, as traders anticipate a 15\% tariff on refined copper by January 2027, and potentially 30\% by January 2028 (Exhibit 2 to Exhibit 3).

The potential tariff stems from the T

[... middle omitted ...]

0/t due to incremental supply released from Comex warehouse.

Bull case scenario could unfold through three pathways, listed in order of increasing bullishness for copper prices: (1) extended review period (i.e., decision will be made in the future), (2) Trump introduces 15\% tariff, with no firm guidance on potential tariff increases in 2028, (3) Trump announces both 15\% tariff starting January 2027 and 30\% tariff starting January 2028.

\subsection{[R003] Bernstein：AI从聊天走向智能体，CPU正在经历一场被低估的复兴}
{\small\color{refgray}归类：AI/算力：智能体时代CPU复兴，中国供应链份额跃迁}

Bernstein：AI从聊天走向智能体，CPU正在经历一场被低估的复兴

过去两年，AI基础设施的投资逻辑几乎被GPU完全定义。市场默认的叙事是：算力即GPU算力，数据中心建设等于加速器集群的扩张。CPU被降格为“配角”，在AI服务器中的价值占比被不断压缩，甚至出现一台服务器搭配8颗GPU、只配1颗CPU的极端配置。

但这份Bernstein研报提出了一个值得所有产业决策者重新审视的判断：随着AI从1.0的聊天机器人范式，转向2.0的智能体范式，CPU正在经历一场结构性复兴。 而市场对这一转变的定价，可能才刚刚开始。

Bernstein将2030年服务器CPU市场空间从此前的1370亿美元大幅上调至2230亿美元，并明确指出——这个数字在牛市情景下可能达到3300亿美元。驱动这一变化的，不是CPU自身的性能突破，而是AI工作负载的底层逻辑正在发生根本性转变。

1. 智能体AI改变了CPU与GPU的配比逻辑，从1:8回归到1:1甚至更高

理解这份报告的核心，首先要理解一个关键比率：CPU与GPU的配对比例。

在传统的AI推理集群中，GPU是绝对的主角。以Google的TPU v6e和Meta的Grand Teton为例，GPU与CPU的配比曾达到8:1。这背后的逻辑是：推理任务本质上是“一次模型调用、一次返回结果”，CPU只需要承担数据加载和结果传递的轻量级工作，GPU负责最密集的矩阵运算。CPU被视作“税”，是必须支付但越少越好的开销。

智能体AI彻底改变了这一假设。当一个AI系统不再只是回答问题，而是需要自主执行任务时，工作流变成了一个循环：用户请求触发检索，系统调用工具、进行中间推理、再次调用模型、执行动作、反馈结果。在这个过程中，GPU仍然负责密集的数学运算，但CPU决定了整个系统能否高效编排——调度任务、管理内存、协调工具调用、避免加速器空闲。

Bernstein的判断是，智能体AI将推动GPU与CPU的配比从当前的8:1快速收窄至1:1甚至更高。报告引用了2026年的硬件路线图作为佐证：AMD的Venice平台是1颗CPU对4颗GPU，NVIDIA的Vera超级芯片是1颗CPU对2颗Rubin GPU，Google的TPU 7x是1颗CPU对4颗TPU芯片。这些设计都在指向同一个方向——CPU不再是附属品，而是与GPU平起平坐的计算核心。

这意味着什么？如果GPU与CPU的配比从8:1回归到1:1，即使AI数据中心的总资本支出不变，CPU的市场空间也将增长数倍。而如果AI资本支出本身仍在高速扩张，CPU将成为叠加了“配比提升”和“总量增长”双重杠杆的最大受益者。

2. 2030年服务器CPU市场空间可达2230亿美元，是当前规模的6倍

基于上述逻辑，Bernstein构建了一套新的服务器CPU市场空间模型。在基准情景下，假设2030年AI资本支出累计达到3.5万亿美元，推理环节的CPU与GPU配比为1:1，训练环节为0.5:1，CPU相对GPU的成本占比从当前的10\%提升至13\%，那么2030年服务器CPU市场空间将达到2230亿美元。

这个数字有多大的想象力？作为对比，2025年服务器CPU市场空间仅约370亿美元。这意味着未来五年，这个市场将以43\%的年复合增长率扩张。而Bernstein此前的预测（1370亿美元）现在被降级为“熊市情景”。在牛市情景下——假设AI资本支出达到4万亿美元、CPU与GPU配比升至1.5:1——市场空间可以高达3300亿美元。

这份模型的关键假设值得推敲。首先，推理环节的CPU配比从当前的0.25倍提升至1倍，是最大的增量来源。这背后是智能体AI推理工作负载从“单次调用”变为“循环调用”的假设。其次，CPU成本占比从10\%提升至13\%，反映了

[... middle omitted ...]

纯的峰值性能。Arm架构在性能功耗比上的优势，在这一场景下被放大。

第二，Arm正在从IP授权商转型为CPU制造商。这一战略转型意味着Arm不再只是收取版税，而是直接参与CPU芯片的设计与销售，能够捕获更大的价值链份额。Bernstein此前预计Arm到2030年可实现150亿美元营收，但基于新的CPU市场空间模型，这一预测被上调至220亿美元。

第三，Arm的2030财年每股收益预测从9.83美元上调至11.79美元，基于42倍市盈率，目标价从300美元上调至500美元。这对应21\%的上行空间。同时，由于Arm是软银的核心资产，Bernstein也将软银的目标价从8200日元上调至11200日元，隐含58\%的上行空间。

需要注意的是，500美元的目标价对应的是2026年180倍以上的市盈率。这显然不是一个基于当前盈利的估值，而是基于2030年远景的折现。投资者需要判断的是：Arm能否在CPU复兴中兑现其市场领导地位，以及市场是否愿意为这一远景支付如此高的溢价。

\subsection{[R004] Bernstein：市场低估的不是供给恢复的速度，而是库存重建的深度}
{\small\color{refgray}归类：能源与大宗：库存重建锚定油价，铜关税赌局6月底见分晓}

Bernstein：市场低估的不是供给恢复的速度，而是库存重建的深度

这份报告的标题“Let the oil flow”看似直白，但真正有穿透力的判断藏在数字背后：全球石油市场在过去一段时间内累计消耗了约10亿桶库存。当所有人都在关注霍尔木兹海峡何时重新开放、供给何时恢复时，Bernstein的分析指向了一个更根本的问题——即便供给恢复正常，市场真正需要的不是“增量”，而是“补量”。重建这10亿桶库存，意味着未来两年内每天需要额外吸收约300万桶的供给。这个数字，才是理解油价底部和资产定价的关键。

这份Bernstein研报发布于一个微妙的时点：布伦特原油年初至今均价约87美元/桶，市场正处于对地缘政治缓和的乐观预期与物理供需现实之间的拉锯中。报告没有给出廉价的看多或看空结论，而是构建了一个“短期偏紧、中期逐渐宽松、长期受边际成本锚定”的分阶段框架。对于亚太油气行业的决策者和投资者而言，这份报告提供的不是简单的方向判断，而是一套理解库存周期如何重塑竞争格局的分析工具。

以下是我们从这份报告中提炼出的四个核心洞察，以及一个报告尚未完全回答的关键问题。

市场普遍将油价维持在80-90美元/桶归因于地缘政治风险溢价。Bernstein的框架提供了一个更坚实的解释：物理库存的消耗。报告详细拆解了这约10亿桶的消耗构成——约3亿桶来自战略石油储备（SPR），约4.5亿桶来自中国库存，约1.4亿桶来自海上浮仓，其余来自商业库存。这不是一个抽象的数字，而是一个已经被市场吸收的供给冲击。

这意味着什么？意味着即便霍尔木兹海峡明天就完全恢复通航，第一批恢复的油流并不会直接满足终端需求，而是会首先进入库存重建通道。Bernstein预计，这个过程需要大约6个月，原因包括扫雷作业、船舶重新部署、战争险保险重新定价、港口和装载计划重建等一系列物理和制度性约束。供给恢复的斜率将是平缓的，而非陡峭的。

这对投资者的含义是：短期油价的下行空间被库存底部封住。报告维持2026年布伦特均价90美元/桶的预测，这一判断的核心依据不是地缘政治，而是“库存极低水平为价格提供了强地板”。任何因“危机结束”预期引发的估值回调，可能更多是买入机会而非卖出信号。

2. 2027年的2.8百万桶/日过剩将被库存重建吸收，市场并非走向失衡

报告预测到2027年，市场将转为约280万桶/日（约10亿桶/年）的过剩。乍看这是一个利空信号。但Bernstein的关键洞察在于，这个所谓的“过剩”大部分将被库存重建所吸收。报告明确写道：“我们预计大部分过剩将被库存重建吸收。”

这个判断的逻辑链条是：经历了本轮大幅消耗后，买家有强烈的动机重建库存。不仅是为了回到正常运营水平，更是为了应对未来可能再次出现的供给中断风险。Bernstein预计，到2027年底，全球库存将基本恢复到冲突前水平。这意味着，2027年的“过剩”在性质上更接近“补库需求”，而非真正的供需失衡。

这一判断对油价的含义是深远的。如果市场将2027年的过剩简单解读为“供给过剩导致油价下跌”，可能低估了库存重建对价格的支撑作用。报告对2027年布伦特原油的预测为78美元/桶，呈现上半年偏强、下半年偏软的走势，长期锚定在75美元/桶——这远高于许多市场参与者基于“供给过剩”叙事所预期的水平。

在油价80美元/桶的假设下，Bernstein指出，亚太油气公司提供了两位数的自由现金流收益率。这是一个值得注意的信号。当市场因为地缘政治缓和预期而抛售这些股票时，估值压缩本身可能创造了一个有吸引力的入场点。

报告中的估值表提供了具体的参照：CNOOC（883.HK）被给予“跑赢大盘”评级，目标价31.70港元，基于2026年预期盈利的市盈率仅为5.2倍；PetroChina（857.HK）同样被给

[... middle omitted ...]

的过程性。报告指出，操作细节仍不明确，包括船舶返回的时间表、战争险保险的可用性和定价、过境船舶的安全协议等。扫雷进展和监督机制同样是关键未知数。

这意味着，即便美伊谅解备忘录是一个建设性的第一步，物理贸易的恢复不会一蹴而就。报告预计恢复过程需要约6个月。这一时间框架与市场可能预期的“几周内恢复正常”存在显著差距。如果市场定价已经开始反映快速正常化，而物理现实却滞后于预期，那么油价和相关资产价格可能面临阶段性的预期差修复。

Bernstein的报告对库存重建的数量给出了清晰的预测，但一个尚未完全展开的问题是：库存重建的“质量”和“分布”如何影响市场结构？

报告提到，一些国家可能会选择大幅增加库存，以应对未来类似事件重演的可能性。这意味着，库存重建不仅是一个总量问题，更是一个结构问题。如果OECD国家优先补充战略石油储备，而中国继续维持高库存水平，那么商业库存的恢复速度可能慢于预期。此外，不同地区、不同品级的原油库存恢复速度不同，可能导致特定品级的价差出现剧烈波动。

\subsection{[R005] DB：日本央行鸽派转向比加息本身更值得关注}
{\small\color{refgray}归类：区域市场：日本央行鸽派加息，银行股定价偏保守}

日本央行在6月会议上加息25个基点，这本该是一条中规中矩的紧缩路径。但DB最新研报揭示了一个被市场低估的信号：一位新任委员在就职不到三个月就投下反对票，这不仅打破了惯例，更预示着未来政策委员会的构成将发生根本性变化。

这份报告的核心判断是：日本央行正在用一种“鸽派加息”的方式完成紧缩，其政策工具的退出节奏存在显著不对称性，而这种不对称性恰恰是市场未来定价的盲点。

投资者习惯盯着利率终点在哪里，但真正决定资产定价的，往往是央行如何到达那里。DB的这份报告，提供了一套理解日本央行“行为逻辑”的新框架。

6月会议以7比1的投票结果通过加息决议，唯一反对票来自今年3月才加入委员会的浅田委员。DB指出，在就职不到三个月内就公开反对行长提案，这在日本央行历史上极为罕见。

这一票的意义不在于它阻止了加息——毕竟加息还是通过了。它的真正意义在于，它暴露了政策委员会内部正在形成的裂痕。浅田委员由首相高市早苗任命，而高市首相在货币政策立场上偏向鸽派。

更关键的是，DB在研报中提醒：另一位新委员佐藤将于6月加入委员会。如果他也采取类似的鸽派立场，那么委员会的投票格局将从“鹰派主导”转向“鸽派力量显著增强”。而到了2027年7月，目前最鹰派的两位委员高田和田中的任期将届满。如果高市首相继续主导继任者的任命，委员会的鸽派倾向将进一步加深。

这意味着什么？市场当前对加息路径的定价，可能严重低估了未来政策决策的政治化风险。一个越来越鸽派的委员会，即便继续加息，其节奏和力度也会受到内部制约。

DB特别关注了前瞻指引的修改。4月声明中“实际利率处于显著低位”的表述被删除，取而代之的是“金融环境一直保持宽松”。

这不仅仅是措辞的修饰。DB解释，这一变化与日本央行3月发布的关于“自然利率与货币宽松程度评估”的论文结论一致。它意味着日本央行正在放弃使用难以准确衡量的自然利率或实际利率作为直接基准，转而基于对“宽松金融环境”的更广泛评估来制定政策。

问题在于，判断“金融环境是否宽松”的指标是多样且模糊的。这给了日本央行更大的裁量权。DB直言，即使在加息之后，要确定金融环境是否仍然“宽松”也需要相当长的时间。这可以被理解为一个隐含信号：日本央行不打算进行激进的紧缩，比如连续加息。它也暗示，央行无意将政策利率提升至限制性区间。

此外，报告指出一个逻辑上的矛盾：如果日本央行预测从2026年10月起经济状况将实现物价稳定目标，那么届时金融环境仍被描述为“宽松”而非“中性”就不合逻辑了。所以，前瞻指引的进一步修订将是一个持续存在的讨论焦点。

对投资者来说，这条信息的含义很明确：日本央行正在为自己留出“走走停停”的空间，加息路径不会是一条直线。

尽管看到了鸽派信号，DB并未改变其加息预测。他们继续预计下一次加息在10月，之后每个季度加息一次，到2027年4月将政策利率提升至1.75\%。

但报告同时指出，市场部分参与者所定价的“每半年加息一次、最终达到2\%”的情景，在浅田委员投下反对票后面临了更高的门槛。DB认为，1.75\%的路径即便在委员会转向鸽派后也是可行的，但2\%的目标需要更鹰派的委员会来支撑。

这里有一个值得深思的细节：报告提到，尽管委员会出现了鸽派分歧，但鹰派委员田中和高田对物价展望的反对意见并未升级。他们仍然认为“基础通胀已经达到物价稳定目标”，但没有主张通胀已经超标。DB判断，这恰恰是为什么本次会议没有提出更激进的50个基点加息的原因之一。

所以，当前日本央行的内部状态是：鸽派在人数上正在增长，鹰派在主张上并未加码。这种力量对比的渐变，比某一次会议的结果更重要。

如果说利率政策是“温和鸽派”，那么日本央行在国债购买退出上的态度，DB用“格外谨慎”来形容。

日本央行决定维持当前缩减计划至2027年第一季度，之后从2027

[... middle omitted ...]

央行在国债市场上给自己留了一扇随时可以“重新介入”的后门。

DB对此提出了质疑：考虑到此前“债券市场参与者会议”上呼吁停止缩减购买的意见属于少数派，央行对停止缩减购买的决定给出的解释并不充分。这种不透明性，可能会在未来市场波动时引发更大的不确定性。

5. 各政策工具退出节奏的不对称，是未来市场与央行对话的核心焦点

报告中最具洞察力的部分，是DB将国债购买退出与其他政策工具的退出进行了对比，揭示了一个显著的不对称性。

以“支持银行贷款的基金供应操作”为例，这个特殊贷款便利在2024年3月收紧了条件，2025年1月决定在2025年6月底停止新增贷款。到2028年3月，其71万亿日元的余额预计将降至零。这一决策推进迅速，没有像债券市场会议那样的公开咨询论坛。DB指出，尽管该操作对金融体系有预期影响，比如减少对国债作为抵押品的需求、迫使银行寻找替代稳定资金来源，但退出的速度并未放缓。

与此同时，对于ETF，日本央行在2025年9月会议上宣布了超过100年的卖出计划。

\subsection{[R006] DB：美国债券市场的真正变化不在需求端，而在持有者结构的不可逆重置}
{\small\color{refgray}归类：宏观与利率：美联储鹰派定价充分，美元区间待突破}

DB：美国债券市场的真正变化不在需求端，而在持有者结构的不可逆重置

这份来自DB、由Matthew Luzzetti等五位分析师联合撰写的研报，提供了一个观察美国固定收益市场的新框架。报告的核心判断并非关于利率走势或通胀预期，而是关于一个更深层的问题：美国61万亿美元债券市场的持有者结构，正在经历一场不可逆的重置。

这份报告的价值在于，它没有停留在“谁在买”的简单统计上，而是试图回答一个更根本的问题：当美联储不再是最大买家、外国央行趋于保守、而国内私人部门开始重新定价风险时，市场运行逻辑会发生怎样的变化。

1. 外国持有者占比已降至30\%，这个数字的背后是结构性而非周期性的退出

报告显示，外国持有美国国债的比例在2026年第一季度降至30\%，接近2023年第三季度29\%的多年低点。这个数字本身并不惊人，但结合全球外汇储备总额的变化来看，趋势值得关注。

全球外汇储备总额在2025年第三季度约为13万亿美元，虽然总量仍在增长，但美元占比持续走低。报告没有给出具体的美元占比数据，但从上下文可以推断，外国央行的增持意愿正在减弱。原因并非简单的“去美元化”，而是更务实的选择：多极储备体系下，各国央行在边际上更倾向于分散配置。

这里的关键洞察是：外国持有者占比下降，不是因为他们抛售，而是因为美国国债发行量增长更快。这意味着，美国必须找到新的国内买家来吸收增量供给。谁将成为这些增量国债的边际买家，将直接影响利率定价和金融市场的稳定性。

2. 一级交易商的国债库存激增至5500亿美元，市场缓冲垫正在变薄

报告中一个容易被忽视但极其重要的数据点：一级交易商的美国国债净持仓自2022年以来急剧上升，到2026年已达到约5500亿美元。这是历史最高水平。

一级交易商本质上扮演的是市场流动性提供者的角色。当他们的库存处于高位时，意味着市场的“天然买家”正在减少，交易商不得不自己持有更多存货。这并非因为他们看好国债，而是因为他们在履行做市义务。

这种现象的潜在含义有两个。第一，市场对冲击的缓冲能力下降。当交易商库存满仓时，他们无法在市场恐慌时大量吸收抛盘。第二，交易商的融资成本上升，这会传导至整个回购市场，进而影响短端利率的稳定性。

报告没有完全展开的是：如果一级交易商的库存继续上升，监管层面的资本约束是否会成为新的瓶颈。这里仍需验证。

3. 银行系统正在重新增持国债，但这次逻辑与量化宽松时期完全不同

报告显示，商业银行持有的国债和机构证券在经历美联储紧缩周期的收缩后，2026年再次回升至约4.8万亿美元，接近2022年的峰值。

但这次的增持逻辑与 年完全不同。当时银行增持国债是因为存款激增，需要配置高流动性资产。而2026年的增持，更多是因为银行在贷款利率下行、信贷需求疲软的环境下，主动寻求收益替代。

这是一个重要的结构性转变。当银行把国债当作“收益资产”而非“流动性资产”来配置时，他们对利率的敏感度会发生变化。如果利率上升导致国债价格下跌，银行可能面临更大的未实现损失，从而影响其信贷投放能力。

报告没有明确讨论的是，这种转变对银行净息差的长期影响。如果银行用国债替代贷款，其盈利能力将更直接地受制于利率曲线形态。

报告的一个重要图表显示，在美联储持有的国债之外，私人持有的可流通国债中，期限在5年以内的占比最大。这意味着，美国国债市场的“长端”正在变得越来越依赖少数几个大型机构投资者。

美联储目前仍持有约30\%的10年期以上国债。随着量化紧缩的推进，美联储在长端的退出将迫使私人投资者吸收更多长期国债。但私人投资者对长端国债的需求并非无限弹性。

报告没有完全展开的是：养老金和保险公司作为长端国债的传统买家，其需求是否足以弥补美联储的退出。从报告数据看，保险公司在信用债市场的配置比例已从1950年代的6

[... middle omitted ...]

地取决于其是否被纳入指数，而非发行人的基本面。

这对高收益债市场的影响尤为显著。报告没有直接讨论的是，如果信用利差在指数配置的推动下被压缩至不合理的低位，一旦市场转向，被动资金可能加速抛售，造成比主动管理时代更剧烈的价格波动。

6. 报告尚未完全回答的一个关键问题：谁将填补美联储退出后的长端缺口

这是整份报告最核心但尚未完全解答的问题。美联储目前仍持有约30\%的10年期以上国债。随着美联储继续缩表，这个比例将持续下降。

从报告数据看，一级交易商和银行系统可能成为长端国债的边际买家。但这两个主体的资产负债表都受到监管限制。外国央行的增持意愿有限。养老金和保险公司的久期需求虽然存在，但其对收益率的敏感度正在提高。

这意味着，长端国债收益率可能需要在更高的水平上才能找到足够的买家。而这将反过来影响抵押贷款利率、企业融资成本和财政部的利息支出。

报告没有给出明确的答案，但提供了一套分析框架：谁在买、为什么买、还能买多少。这三个问题的答案，决定了美国利率的长期中枢。

\subsection{[R007] GS：日本央行加息至1\%，市场真正低估的不是终点利率，而是加息节奏的重新定价}
{\small\color{refgray}归类：区域市场：日本央行鸽派加息，银行股定价偏保守}

GS：日本央行加息至1\%，市场真正低估的不是终点利率，而是加息节奏的重新定价

6月16日，日本央行以多数票通过将政策利率上调25个基点至1.0\%，并维持了从明年4月起每月约2万亿日元的国债购买计划。这看起来是一次符合预期的“温和加息”。然而，GS在第一时间发布的研报中揭示了一个被市场普遍低估的深层信号：日本央行对通胀上行风险的承认，正在悄然改变未来加息的节奏预期。

市场习惯于将此次加息视为“一锤子买卖”的终点，但GS的分析指向了一个不同的结论。真正需要关注的，不是1.0\%这个数字本身，而是日本央行在措辞中从“实际利率极低”转向“金融环境宽松”这一表述变化背后蕴含的政策逻辑。这意味着，日本央行正在为后续的连续加息构建新的叙事框架，而市场对此的定价远未充分。

这份报告的核心价值，不在于复述加息决定，而在于它精准地指出了市场与央行之间在“加息速度”和“终端利率”两个维度上的认知鸿沟。对于持有日本资产或关注全球利率环境的决策者而言，忽略这个信号可能会付出代价。

1. 加息本身不是新闻，但“承认通胀上行风险”才是真正的政策转向

GS研报中最值得细读的，不是加息25个基点这一动作，而是日本央行在声明中明确写入的一句话：“存在基础CPI通胀偏离2\%价格稳定目标上行的风险。” 这句话在4月的展望报告中尚未出现，而此次被正式纳入政策声明，意味着日本央行内部对通胀路径的评估已经发生了实质性变化。

在过去的加息周期中，日本央行始终强调通胀的“暂时性”和“成本推动”特征，以此安抚市场对过快加息的担忧。但现在，措辞从“通胀正在放缓”调整为“存在上行风险”，这等同于承认：工资上涨向价格传导的机制可能比预期更持久、更广泛。

GS的分析师敏锐地捕捉到，这一表述变化直接关联到后续加息紧迫性的判断。如果通胀风险上行，那么下一次加息的时间窗口就会从“明年某个时候”提前到“今年下半年”。市场当前隐含的加息路径，显然还没有充分计入这一风险重估。这意味着，日元利率的上升速度和幅度，都可能超过当前债券市场的定价。

2. 中性利率区间给出“1.1\%至2.5\%”，但市场只看到了下限

日本央行在今年3月首次发布了中性利率的官方估算，区间为1.1\%至2.5\%。此次加息至1.0\%后，政策利率刚好触及该区间的下沿。市场的一个普遍解读是：加息任务已基本完成，后续空间有限。

GS的报告对此提出了一个关键的、但尚未被充分讨论的质疑：中性利率是一个区间，而非一个点。市场选择性地关注1.1\%的下限，可能忽略了2.5\%上限所隐含的政策空间。更重要的是，中性利率本身并非固定不变，它会随着经济增长潜力和结构性改革进展而动态调整。

此次副行长内田真一的新闻发布会，GS特别关注其对中性利率的看法。如果日本央行认为当前经济产出缺口已经转正、且工资-通胀螺旋正在形成，那么其对中性利率的估算就可能向上修正。换句话说，1.0\%的利率可能不是终点，而是新一轮利率正常化的起点。那些认为“日本利率将长期停留在1\%附近”的资产定价逻辑，可能建立在错误的假设之上。

3. 国债购买计划的“不变”之中，藏着一个关于退出节奏的伏笔

日本央行决定维持现有国债购买计划，即到明年一季度末将月度购买量降至约2万亿日元，并在2027年3月底前维持这一水平。这被市场解读为“鸽派”信号——央行没有加速缩减购债。

但GS的分析提供了一个不同的视角：维持购买计划不变，恰恰是为了给未来的加息创造更稳定的市场环境。如果日本央行在加息的同时大幅缩减购债，可能引发国债市场的剧烈波动，反而不利于利率传导。维持购债规模，是为了确保加息本身能够被市场平稳吸收，而不是央行对宽松立场的留恋。

更值得关注的是，日本央行宣布将停止对购债计划的年度中期评估。这意味着，未来购债规模的调整将不再遵循预设的时间表，而是完全取决于

[... middle omitted ...]

日元升值。这个反馈循环的速度和幅度，将直接影响日本央行能否避免“落后于曲线”。

当前市场对日本利率的定价，隐含了一个假设：日本央行有能力在通胀失控之前稳步加息。但如果日元贬值导致的输入性通胀加速到来，日本央行可能被迫在短时间内连续加息，从而对日本国债和全球套利交易产生冲击。报告提到“落后于曲线的风险”，但并未给出这一风险的具体触发条件或概率估算。这正是需要读者在完整报告中深入挖掘的部分。

5. 一个被忽略的观察框架：从“实际利率”到“金融条件”的表述切换意味着什么

此次声明中，日本央行将之前关于“实际利率处于极低水平”的表述，修改为“金融条件一直保持宽松”。GS指出，这一调整部分是为了纳入此次加息的影响。

但这一措辞变化的含义远不止于此。“实际利率”是一个客观的经济学概念，而“金融条件”是一个更宽泛、更具主观判断色彩的表述。它意味着，日本央行未来在决定是否继续加息时，将不再仅仅盯住实际利率与中性利率的差距，而是会综合评估信贷、汇率、资产价格等更广泛的金融环境。

\subsection{[R008] JPM：出口仍是唯一结构性亮点，但国内需求疲软正在拉长政策等待期}
{\small\color{refgray}归类：未被主题摘要引用}

JPM：出口仍是唯一结构性亮点，但国内需求疲软正在拉长政策等待期

这份由JPM新兴市场亚洲经济与政策研究团队发布的第57期中国高频数据追踪报告，传递了一个核心信号：5-6月的经济活动数据并未形成清晰的复苏趋势。出口侧的运量数据在6月出现了显著的同比加速，但国内生产、信贷、汽车销售和房地产市场的信号依然偏弱。JPM团队通过高频指标与官方月度数据的映射，给出了一个值得产业决策者关注的关键判断——二季度增长动能的重新校准，将高度依赖6月剩余时间的数据表现，而政策的主动发力窗口可能正在向后推迟。

报告中最引人注目的变化来自出口侧。截至6月中，港口跟踪数据显示，离港集装箱船和散货船的载重吨位分别环比增长8.9\%和10.8\%，同比增速更是从5月的7.7\%跃升至17.9\%。这意味着6月的出口量增长正在与近期运价上涨形成共振。然而，这一亮点的背面是国内需求的持续承压。JPM团队指出，4月国内活动指标的疲软已经将二季度增长风险向下倾斜，而5-6月的混合数据并未提供充分的逆转证据。对于投资者而言，真正的挑战不在于判断出口能否维持，而在于理解国内需求疲软何时会触发更实质性的政策响应，以及哪些资产类别会对这一响应率先定价。

1. 出口的强劲并非孤例，但运力与运价的同步上行正在重塑贸易利润分配

JPM的高频追踪显示，6月以来的出口表现超出了季节性规律。不仅集装箱船离港量增长加速，散货船离港量同比增速更是达到了21.7\%，环比增长10.8\%。与此同时，中国出口集装箱运价指数(CCFI)至美东和美西航线分别较两周前上涨6.5\%和9.1\%，至波斯湾/红海航线也上涨了8.6\%。运量与运价的同步上行，意味着出口企业面临的并非单纯的“量”的扩张，而是运输成本端也在快速抬升。

这里需要关注一个容易被忽视的逻辑：运价上涨在短期内可能改善航运企业的利润预期，但对于以FOB(离岸价)结算为主的制造业出口商而言，若运价涨幅无法完全转嫁，其实际利润率可能受到挤压。JPM报告没有直接讨论这一利润分配问题，但CCFI数据与出口量数据的叠加走势，为投资者提供了一个重要的二阶观察维度——当运价上涨速度超过出口量增速时，贸易环节的利润可能正在从生产端向物流端转移。这一判断对于评估出口相关产业链的盈利质量具有前置意义。

2. 国内生产端的分化正在加剧：原油加工收缩深化，钢铁边际改善但基础不稳

JPM通过沥青厂开工率、轮胎厂开工率和螺纹钢开工率等高频指标，分别映射了原油加工、汽车和钢铁三大工业领域的生产走势。数据显示，原油加工量的收缩可能在5月进一步加深，并在6月延续恶化，这与此前4月同比下降5.8\%的趋势一致。汽车工业增加值在5月或许有温和改善，但6月又出现放缓，4月同比已是下降2.6\%。钢铁生产的收缩在5月可能加深，但6月存在收窄的可能性。

这种分化背后反映的是不同行业面临的供需矛盾差异。原油加工量的持续收缩与石油沥青厂开工率下降直接相关，这指向了基建和道路施工需求的疲软。汽车生产的反复则与终端销售低迷密切相关——5月乘用车零售同比下降22\%，新能源汽车虽降幅较窄为7.5\%，但6月第一周乘用车零售降幅仍达23\%。钢铁生产的边际改善可能更多来自供应端的主动调整而非需求端的强劲复苏，其可持续性仍存疑。对于投资者而言，这三个高频指标的组合，实际上提供了一个比单一PMI数据更精细的工业生产景气度评估框架。

报告披露了一个值得关注的变化：6月前两周，30个大中城市的新房和二手房销售表现均好于去年同期，新房销售同比增速从5月的下降1.4\%转为增长4.0\%。然而，JPM团队在报告中明确点出了不确定性——“这是否标志着房地产市场触底仍有不确定性”。

这里的核心矛盾在于，月度销售数据的改善是否具备可持续性。过去一年中，中国房地产市场曾多次出现单月销售反弹，但随后均被证

[... middle omitted ...]

明显放缓，当月发行量仅为6160亿元，远低于5月的11400亿元。其中，国债发行降至1640亿元，而专项地方债发行仅温和回升至2120亿元。报告指出，除非月底出现跳升，否则6月的财政交付将保持“不太前置”的状态。

在货币政策端，央行通过质押式逆回购净投放了4100亿元流动性，但同时通过买断式逆回购回笼了3000亿元，整体净投放规模有限。JPM团队认为，如果增长阻力持续加剧并超过通胀风险，降息在2025年下半年会变得更有可能。

这些数据组合在一起，描绘了一幅与市场预期可能不同的政策图景。市场此前普遍预期二季度会有更积极的财政扩张和货币宽松，但实际数据表明，政策制定者可能正在等待更多经济数据来确认疲软的程度和持续性。这种“观察等待”的姿态本身就是一个重要信号——它意味着短期内政策超预期的概率在下降，而三季度才是更可能看到政策加速落地的窗口。对于债券投资者而言，这意味着利率曲线的平坦化趋势可能延续，而对于权益投资者来说，依赖政策预期驱动的板块可能需要重新定价其等待成本。

\subsection{[R009] 摩根斯坦利：市场真正低估的不是AI泡沫，而是供给侧的定价权转移}
{\small\color{refgray}归类：资本市场：资本支出超级周期重塑资产定价}

摩根斯坦利：市场真正低估的不是AI泡沫，而是供给侧的定价权转移

这份研报的标题叫“Charts That Caught My Eye”，出自摩根斯坦利全球研究总监Katy Huberty之手。它看起来像是一封内部备忘录——用七张图表串联起当前最值得关注的宏观与主题信号。但读完之后你会发现，这根本不是一份轻松的“看图说话”。它试图回答一个更根本的问题：当市场还在争论AI是不是泡沫、关税有没有效果、中国经济何时见底时，哪些结构性力量已经在静悄悄地改变资产定价的底层逻辑？

我们的判断是：这份报告隐含的主线，并非AI主题的短期涨跌，而是全球资本正在经历一次“供给侧定价权”的再分配。从AI基础设施的“逢低买入”策略回测，到欧洲行业模型的重新排序，再到美国回流数据的“名实不符”，再到中国K型经济的固化——每一张图表都在指向同一个方向：那些掌握稀缺供给能力的企业和资产，正在获得前所未有的议价权。而市场对这一变化的定价，仍然不够充分。

1. AI主题的“逢低买入”策略并非情绪交易，而是对供给稀缺性的结构性套利

报告中最引人注目的实证，是摩根斯坦利对AI基础设施和电力AI两个子主题开发的“逢低买入”策略回测。其规则很清晰：当主题组合下跌5\%时加仓至150\%，下跌15\%时加仓至200\%，加仓持仓三个月后恢复基准。回测结果显示，这一策略自2023年以来在两个子主题上都产生了显著的超额收益。

这意味着什么？不是“AI跌了就该买”这么简单。它的深层含义是：AI基础设施的供给端存在结构性瓶颈，每一次市场情绪导致的回调，都会被后续的供需基本面修正所“救赎”。换句话说，这个主题的波动更多来自需求侧的预期摇摆，而非供给侧的实际松动。当市场因为某个短期担忧抛售时，真正稀缺的算力、电力、冷却资源并不会因此增加供给。因此，逢低买入本质上是在对“供给刚性”进行套利。

这对投资者的启发是：在AI主题中，真正值得关注的不是那些需求故事讲得最动听的标的，而是那些在供给端拥有不可复制优势的企业——无论是芯片制造能力、电力接入权，还是数据中心的地理位置稀缺性。这些才是“逢低买入”策略能够持续有效的底层支撑。

2. 通胀阈值3\%是资产定价的“分水岭”，市场对宏观数据的解读正在发生非线性跳跃

报告对美国经济数据的分析揭示了一个被广泛忽视的非线性特征：当CPI高于3\%时，通胀超预期对利率市场的影响显著放大；而当CPI在2\%-3\%之间时，只有股市对数据有显著反应；一旦CPI低于2\%，无论利率还是股市都不再对数据惊喜做出可靠反应。

这个发现的重要性在于，它否定了“数据越强，资产越涨”的线性思维。当前美国CPI正处于3\%附近的敏感区间。如果通胀数据持续高于这个阈值，市场将进入一个“坏数据才是好消息”的反直觉状态——因为只有经济放缓才能让通胀回落，从而为利率下行打开空间。反之，如果通胀意外回落至3\%以下，市场焦点会迅速从通胀切换至增长，届时就业数据的强弱将重新主导资产定价。

这意味着，未来几个月的CPI读数，不仅是通胀本身的问题，更是资产定价框架的切换开关。投资者需要为两种截然不同的市场环境做好准备，而不是简单地押注“通胀回落利好所有资产”。

3. 欧洲行业模型的重排揭示了一个关键信号：AI的资本支出正从“概念”变为“现金流”

报告对欧洲行业模型的更新值得仔细推敲。半导体仍居榜首，但金属与矿业跃升至第二，资本品被上调至超配，银行升至第三——这些变化看似分散，实则围绕同一个核心：AI资本支出的传导链条正在从上游的芯片设计，向下游的电力、冷却、工业设备和融资需求延伸。

半导体排名第一并不意外。真正值得关注的是金属与矿业的跃升。这意味着，AI基础设施的建设已经进入了需要大量实物资源的阶段——铜、铝、稀土等关键矿产的需求正在从预期变为现实。同样，资本品被上调

[... middle omitted ...]

重组

报告对美国“回流”数据的解读是一记清醒剂。虽然2025年名义产出数据看起来不错，但报告明确指出，这主要是价格驱动的。进口渗透率在2025年不降反升，表明美国对进口的依赖并未减弱。真正发生的，是供应链的“重新定向”——从中国转向东南亚、墨西哥等地，而非大规模回流美国本土。

更关键的是，报告指出目前看到的有限投资是“窄而长期”的，主要集中于AI产能建设。这意味着，当前的制造业投资并非传统意义上的“再工业化”，而是AI基础设施建设的副产品。如果政策环境（尤其是USMCA的延续性）发生变化，这一轮投资能否从AI生态扩展到更广泛的制造业，仍然存在很大不确定性。

这个判断对资产定价的含义是：不要将当前制造业投资的增长简单解读为“美国制造业复兴”的开始。它更像是一次由AI驱动的、高度选择性的资本支出周期。那些与AI基础设施建设无关的传统制造业，可能并未从中受益。投资者需要区分“AI带动的结构性需求”和“关税引发的短期库存调整”，前者可持续，后者则可能随着政策变化而逆转。

\subsection{[R010] GS：市场低估的不是供需平衡，而是安全溢价的结构性重置}
{\small\color{refgray}归类：能源与大宗：库存重建锚定油价，铜关税赌局6月底见分晓}

当特朗普宣布达成临时协议、霍尔木兹海峡即将重新开放的消息传出，多数市场参与者的第一反应是油价将大幅回落。这份直觉没有错，但可能过于简单。GS在最新发布的研报中做出了一个看似矛盾、实则深刻的判断：即便在供应提前恢复正常的情景下，油价底部比大多数人想象的更坚实。

GS将2026年第四季度布伦特原油预测从90美元下调至80美元，2027年均价从80美元下调至75美元。表面看这是看空调整，但隐含的结论是：即便全球石油市场在2027年面临每天320万桶的过剩，油价仍将维持在每桶75美元左右。这个数字恰好是GS的长期公允价值。

这组预测的真正张力在于：一个被历史性供应冲击打乱的市场，在冲击解除后，并不会回到冲击前的定价逻辑。因为市场已经经历了一次“结构性的安全溢价重定价”。

1. 供应恢复的时间差被压缩，但价格弹性的不对称性更值得关注

GS将波斯湾出口恢复至战前水平的时间节点从8月底提前到7月底，这个一个月的提前意味着2026年第四季度和2027年油价公允价值分别下调约10美元和5美元。这个调整的幅度本身并不令人意外，但需要追问的是：为什么下调幅度不是对称的？

答案藏在供需弹性的不对称中。当供应冲击发生时，价格上行的幅度远超基本面隐含的均衡水平——这是稀缺溢价。但当冲击解除时，价格下行却受到多重因素的托底，不会跌回冲击前的水平。GS称之为“安全溢价”。

这种不对称在2027年的预测中表现得尤为明显。尽管GS预计2027年全球石油市场将出现每天320万桶的过剩，但布伦特原油价格预测仍维持在75美元，与长期公允价值一致。这意味着，即便在供过于求的背景下，市场仍愿意为地缘政治风险支付一个结构性溢价。

这个判断对资产定价的含义是：传统基于库存和供需平衡的油价预测模型需要重新校准。安全溢价不再是零均值波动的短期变量，而可能成为中长期定价的一个固定组成部分。

2. 2027年的过剩为何不会压垮油价：库存结构与战略储备的范式转变

每天320万桶的过剩，放在任何历史时期都足以让油价大幅承压。但GS认为这个过剩不会产生预期的价格压力，理由有两个。

第一个理由是库存结构。OECD商业石油库存在2026年上半年经历了大规模消耗，这意味着即便2027年出现过剩，库存水平也很难回升到高位。GS特别指出，全球战略储备的结构性增长趋势达到每天超过100万桶——这个数字本身就是一个重要的结构性变化。各国在经历2026年的供应冲击后，对战略石油储备的重视程度显著提升，这种“补库存”行为吸收了大量的过剩供应。

第二个理由更值得政策制定者和企业决策者关注：安全溢价的存在意味着市场对供应中断风险给出了一个持续的定价。这个溢价不会因为霍尔木兹海峡重新开放就立即消失，因为市场需要时间来验证“重新开放”是否可持续。

但GS也给出了一个重要的约束条件：这个安全溢价的上限可能有限。原因在于，市场刚刚经历了一次史无前例的每天1400万桶的石油产量冲击，而实际供需缺口只有每天500万桶。这意味着全球市场（尤其是中国需求）展现出了惊人的弹性。这种弹性既证明了市场调节能力，也意味着一旦中东出口恢复正常，安全溢价可能迅速收窄。

对于投资者而言，关键问题不是2027年的过剩有多大，而是这些过剩会被谁吸收、以什么价格吸收。战略储备的补库需求为油价提供了一个隐性底部。

GS研报中一个容易被忽略但可能带来重大影响的变量是伊朗。报告指出，伊朗产量可能因制裁放松而升至战前水平以上。这个判断如果成立，将改变整个中东供应恢复的节奏和幅度。

伊朗在2026年的供应冲击中是一个“受害者”——其出口受到霍尔木兹海峡关闭的直接影响。但如果海峡重新开放，且伴随制裁的实质性放松，伊朗可能成为供应恢复中最具弹性的力量。GS没有给出具体数字，但暗示伊朗产量的上行空间可能超出市场预期。

[... middle omitted ...]

极端情景。在上行情景中，假设霍尔木兹海峡持续关闭到2027年底，布伦特原油可能在2026年底升至130美元以上，2027年均价达到105美元。在下行情景中，假设出口在7月初就恢复正常，叠加需求损失和更强的供应，布伦特原油2026年第四季度均价可能略低于70美元，2027年均价略低于60美元。

这两个情景的对比揭示了一个重要的不对称性：上行空间（从75美元到105美元，涨幅40\%）远大于下行空间（从75美元到60美元，跌幅20\%）。GS明确表示，尽管风险是双向的，但净方向偏向上行。

这个不对称性背后有一个逻辑支撑：供应恢复是一个可逆性存疑的过程。即便协议签署，实际执行可能面临排雷、船东风险规避、以及伊朗可能再次关闭海峡等不确定性。而需求方面，GS预计2026年下半年到2027年将出现较为强劲的复苏，原因是油价回落后可负担性改善。

对于产业决策者而言，这个不对称性意味着：在制定2027年的采购或对冲策略时，应该给予上行风险更高的权重，而不是简单地按照基准预测来配置。

\subsection{[R011] GS：市场低估的不是电商增速放缓，而是云计算在AI需求下的重新定价}
{\small\color{refgray}归类：未被主题摘要引用}

GS：市场低估的不是电商增速放缓，而是云计算在AI需求下的重新定价

中国互联网板块正在经历一次罕见的估值分化。一方面，电商GMV增速在5月仅录得3\%的同比增长，低于市场预期，零售数据疲软直接压低了以阿里、京东为代表的电商平台股价。另一方面，云计算与AI基础设施的投资叙事却持续升温，GS在最新报告中明确将“云与数据中心”列为其最看好的子板块。这种分化背后隐藏着一个更深层的判断：市场对互联网公司的定价逻辑，正在从“电商交易规模”转向“AI计算能力的供给与议价权”。真正被低估的，不是电商需求何时复苏，而是云计算在AI需求爆发中获得的重新定价空间。

这份GS于2026年6月发布的《Navigating China Internet: Ecommerce Tracker》报告，看似是一份常规的电商追踪月报，实则是对阿里近期股价疲软的全面回应，并在此基础上重构了整个中国互联网的投资排序。报告的核心信号是：电商的短期压力是真实的，但云计算的结构性机会正在被低估，而阿里恰好是这两个叙事交汇点上的关键标的。

1. 五月零售数据的“坏消息”已被市场消化，电商竞争格局正在从价格战转向运营效率

五月全国线上零售商品GMV同比增长3\%，相较于四月的零增长有所改善，但整体零售数据低于彭博一致预期。珠宝、家电、体育用品等品类分别录得-8.9\%、-15.6\%和-8\%的同比下滑，显示出消费端的谨慎情绪并未消散。家电品类的进一步下滑尤其值得关注，GS指出这与去年“以旧换新”政策带来的高基数有关，手机品类自2025年1月被纳入换新补贴后，基数效应仍在拖累同比增速。

对于市场最关心的阿里CMR（客户管理收入）问题，GS的判断是：市场已经预期了二季度CMR的低个位数同比下滑，而包裹量的健康增长可能会缩小GMV与CMR之间的缺口。这意味着，商户返利政策虽然短期压制了CMR，但只要GMV目标达成，返利规模自然会收缩。更关键的是，行业竞争正在从“烧钱抢份额”转向“更理性的运营”。618大促期间，平台开始简化促销机制、延长统计窗口，监管层也在收紧不正当竞争和虚假促销的规则。这种转变的直接含义是：电商平台的利润底部正在形成，而不是继续恶化。

对于投资者而言，电商板块的真正拐点不在于GMV增速何时反弹，而在于市场何时开始用“利润稳定性和竞争格局改善”来重新定价这些平台。GS将京东列为2H26的关键标的，核心逻辑正是“二季度将是最后一个困难的同比比较季度，之后将迎来收入和利润的同比复苏”。

2. 阿里股价疲软背后是五重担忧，但GS认为其中三个是“伪命题”

阿里近期股价从年初至今下跌约28\%，EPS下调和估值收缩各贡献了一半跌幅。GS将市场的主要担忧归纳为五个方面：电商CMR疲软、AI云计算面临运营商竞争、AI模型层竞争加剧、AI资本开支资金来源、以及SOTP折价扩大。但报告的核心立场是，这五个担忧中有三个被市场过度放大。

关于AI云计算与运营商的竞争，GS的回应最为有力。市场担心政府计划在未来五年投入2万亿元建设数据中心，且主要由电信运营商运营，将挤压互联网云厂商的空间。但GS指出，仅中国互联网四大云厂商（百度、阿里、腾讯、字节）未来两年的资本开支就将达到1.5万亿元。在AI Token需求爆发式增长、算力供给短期偏紧的背景下，互联网云厂商在资本开支规模、AI模型能力、以及MaaS（模型即服务）和自研芯片上的先发优势，将使其在中国AI云市场中获得比传统的CPU/IaaS云市场更有利的定价能力和利润结构。这份报告没有完全展开的是，运营商与互联网云厂商在AI云市场可能不是零和博弈，而是各自占据不同的客户层级和应用场景。

关于AI模型层的竞争，GS承认定价竞争比预期更激烈，但指出顶级模型的定价权仍然在维持——GLM和Qwen3.7 Max就是例证。阿里

[... middle omitted ...]

和京东三大平台的日均订单量和单位经济（UE）正在发生微妙变化。阿里的日均订单量从2025年3月的2500万单增长到2026年6月的6600万单，美团的订单量则从6800万增长到8000万单，京东维持在 万单的水平。

更关键的是UE的变化。GS判断，随着反垄断调查的推进以及阿里进一步转向AI/云战略，前三大外卖平台在即时零售领域的市场份额趋于稳定，补贴行为变得更加理性。这带来的直接结果是，美团的食品配送业务（不含平台层面成本）在2Q26和FY26E的EBIT可能分别达到盈亏平衡和亏损17亿元，相较于此前的大幅亏损是一个明显的改善。

对于投资者来说，这个细分领域的含义是双重的。一方面，即时零售的UE改善将为美团和阿里贡献可观的利润增量，这部分在当前的估值中可能没有被充分反映。另一方面，京东在即时零售上的投入规模相对有限，这意味着其利润修复的节奏可能比市场预期的更快。GS将京东列为电商与出行子板块的首选，这个判断的背后逻辑是“竞争缓和+利润修复”，而非单纯的GMV增长。

\subsection{[R012] GS：白酒端午动销预期转弱，市场低估了渠道库存周期的“钝化”}
{\small\color{refgray}归类：地产与消费：一线城市结构性再定价，消费疲软拉长政策等待期}

GS：白酒端午动销预期转弱，市场低估了渠道库存周期的“钝化”

当市场还在讨论“低基数能否带来端午反弹”时，一份来自GS的最新渠道追踪报告给出了一个更冷静的判断：今年的端午节，可能不会成为白酒行业的拐点。

这份于6月中旬发布的研报，通过密集的渠道调研和价格数据跟踪，揭示了一个值得产业决策者关注的信号。尽管2025年同期因政策冲击形成了低基数，但经销商对今年端午的零售增长预期反而从此前的20\%-30\%下修至10\%-20\%。飞天茅台批发价在 元区间窄幅波动，而普五和国窖1573的批价则持续疲软，其中普五在部分批发市场的报价甚至受618补贴压力跌至760元。

这些数字背后，不是一个简单的“旺季不旺”的故事。它指向一个更深层的结构性变化：白酒行业的渠道库存周期正在失去弹性，价格信号对短期消费刺激的反应正在钝化。对于投资者和管理者而言，理解这种“钝化”的成因，比猜测下一个节日的销量数字更重要。

1. 经销商预期的下修并非季节性扰动，而是渠道信心的系统性降温

GS报告中最值得关注的数据点，并非批发价格的绝对水平，而是经销商预期的变化轨迹。从“期待20\%-30\%增长”到“预期10\%-20\%增长”，这个下修发生在2025年低基数背景之下。这意味着，经销商并非因为去年卖得差而悲观，而是因为当前动销环境的疲软超出了他们对“低基数反弹”的原有判断。

渠道备货节奏是另一个佐证。报告指出，品牌方在节前的渠道动作相对平淡。在白酒行业，节前品牌方通常会通过密集的渠道政策、促销活动来推动经销商打款备货。如果连品牌方自身都选择“按兵不动”，这通常意味着它们对终端真实需求的判断同样保守。

这里需要强调的是，端午节在白酒消费周期中的权重远不及中秋国庆。GS在报告中明确提示，中秋国庆才是观察行业拐点更重要的窗口。因此，端午预期的下修本身不构成趋势逆转的证据，但它是一个值得警惕的前瞻信号：如果连低基数都无法激发渠道的备货热情，那么接下来的中秋旺季，需要多大的需求刺激才能扭转预期？

2. 飞天茅台的价格稳定是“量价平衡”的结果，而非需求强劲的证据

飞天茅台批发价在 元区间的稳定，很容易被解读为高端需求的韧性。但仔细拆解报告中的细节，会发现这个稳定背后是供给侧的主动调节。

核心变量来自i茅台。报告显示，i茅台在2026年5月的月活跃用户达到956万，5M26的交易量约713万笔，远高于1Q26的398万笔。这些数据意味着，茅台通过直营平台正在向市场投放更多非标产品——生肖酒、1L装等。这些非标SKU的批发价格在过去两周分别下跌了60元和30元，直接反映了供给增加带来的价格压力。

换言之，飞天标准品的价格稳定，部分是以非标产品价格承压为代价的。茅台正在通过产品结构的调整，在维持核心单品价格体系的同时，消化渠道库存。这是一种主动的“以价换量”，但代价是非标产品的渠道利润被压缩。对于投资者而言，需要追问的是：非标产品的价格持续走软，是否会反过来影响飞天茅台的价格预期？如果非标产品的批价持续低于市场预期，渠道商可能会重新评估整个茅台产品组合的持有价值。

3. 普五和国窖1573的批价分化，揭示了品牌议价能力的根本差异

在高端白酒的价格梯队中，普五和国窖1573的处境比茅台更为严峻。报告显示，普五的批发价格在不同数据源之间存在显著差异：每日酒价报价840元，而百荣批发市场报价已跌至760元。后者受到了618电商补贴的直接冲击。国窖1573则稳定在840元。

这种价格分化不是偶然的。普五在部分批发市场的价格跌破800元，意味着渠道商的实际出货价已经接近甚至低于经销商从厂家拿货的成本。当渠道利润被压缩到极限时，经销商的打款意愿和库存管理策略会发生根本性变化。他们不再愿意为厂家承担库存风险，而是倾向于“零库存”运营，只在接到终端订单后才向厂家要货

[... middle omitted ...]

削弱。过去，渠道商可以通过信息差和库存控制来影响批发价格；现在，i茅台提供了一个透明的、高频的定价参考。

这意味着，茅台正在逐步将定价权从渠道端收回至厂家端。长期来看，这有利于茅台控制价格体系，但短期也会带来阵痛：当渠道商意识到自己不再拥有定价主导权时，它们囤货的意愿会下降，价格波动反而可能加剧。GS报告中没有完全展开这个二阶效应，但这是理解未来茅台价格走势的关键变量。

报告提到，国务院食安办在6月12日宣布查处了75,200箱假冒“特供”酒类产品，关闭了5家公司和36个销售平台。这个新闻在报告中被列为“关键新闻”，但并未被赋予过高的权重。

从短期看，这次打击行动对主流品牌的影响有限，因为被查处的主要是造假企业和非正规渠道。但从长期看，它释放了一个重要信号：监管层正在清理白酒流通环节的灰色地带。过去几年，部分中小酒企通过“特供”“专供”等概念进行差异化营销，绕开了主流品牌的价格体系。如果这种模式被系统性打击，那些依赖概念营销而非品牌力驱动的企业将面临更大的压力。

\subsection{[R013] GS：5月经济数据低于预期，但政策窗口正在打开}
{\small\color{refgray}归类：地产与消费：一线城市结构性再定价，消费疲软拉长政策等待期}

中国经济在二季度中段遭遇了比预期更猛烈的逆风。5月固定资产投资与零售数据双双低于市场预期，工业产出大体符合预期。这份GS研报的核心判断并非停留在“数据不好”这一表层，而是提出了一个值得所有决策者仔细推敲的叙事：短期疲软背后，结构性分化正在加剧，而政策空间与外部环境的边际改善，可能正在为三季度积蓄反弹动能。

这不是一份简单的“数据点评”。它揭示了几个关键信号：服务消费与商品消费的增速缺口在持续扩大；投资端的拖累并非均匀分布，天气与债券发行节奏的扰动被市场低估；而房地产虽然在低位，但一线城市的“绿芽”并未完全枯萎。更重要的是，GS明确指出了7月政治局会议作为潜在政策调整窗口的战略意义。

对于关注中国资产定价的读者而言，这份报告最值得深思的地方，不是5月数据有多差，而是市场是否正在低估供给侧的再定价能力，以及政策应对的灵活性。

5月工业增加值同比增长4.5\%，较4月的4.1\%温和回升，基本符合预期。表面上看，这是好消息。但GS的拆解揭示了一个更复杂的图景：支撑增长的主要是计算机与电子设备、电力热力等行业，而汽车、化工、有色金属冶炼的产出增速却在放缓。

汽车产量同比增速从4月的-2.6\%进一步下滑至5月的-3.2\%。这并非孤立现象。全球能源冲击持续压制化工相关制造业，硫酸和化学纤维产出增速依然疲弱。与此同时，智能手机和计算机产量同比分别下降8.8\%和19.4\%，消费电子链条的压力依然显著。

真正值得关注的是，工业机器人产量同比增长27.9\%，金属切削机床增长10.7\%。这组数据暗示，自动化设备投资仍在加速，企业并非全面收缩，而是在有选择地加大资本开支。这背后反映的，是制造业内部的结构性分化：面向出口的自动化设备需求旺盛，而面向内需的消费品生产仍在承压。

*这意味着，单纯看工业增加值的总量，会掩盖行业间的巨大温差。对于投资者而言，真正需要跟踪的，不是4.5\%这个数字，而是哪些行业正在从出口中获得定价权，哪些行业仍在承受能源成本与内需疲软的双重挤压。**

2. 固定资产投资的下滑并非全是需求问题，天气与债券节奏被低估

5月固定资产投资单月同比大跌10.6\%，较4月的-8.2\%进一步恶化。这个数字看起来触目惊心，但GS特别指出，两个临时性因素可能夸大了实际疲软程度：一是南方和华中地区的持续强降雨以及北方的高温热浪，二是政府债券发行节奏仍然偏慢。

基础设施投资单月同比下滑至-11.2\%，制造业投资也降至-4.1\%。但值得注意的是，粗钢和水泥产出同比降幅在5月有所收窄，分别从-2.8\%和-10.8\%改善至-2.7\%和-8.1\%。GS谨慎地提示，国家统计局偶尔会对此前高报的数据进行“统计修正”，这可能导致近几个季度FAI公布数据的波动性被放大。

这里有一个值得深入推敲的点：如果剔除天气和债券发行节奏的扰动，真实的投资需求究竟有多弱？报告并未给出明确答案，但逻辑指向一个判断：下半年随着财政资金加速落地，基建投资有望出现边际改善。 关键在于，改善的幅度是否足以对冲房地产投资的持续下滑。

*对于产业决策者而言，5月的FAI数据不应被简单解读为需求的全面崩塌。它更像是一个被临时因素放大的信号，真正的考验在于三季度，当天气好转、债券资金到位后，投资能否如期企稳。**

3. 零售数据的结构性裂痕：大企业跑输小企业，服务消费跑赢商品消费

5月名义零售额同比转为-0.6\%，创下2022年12月以来的最低水平。但GS的分析没有停留在总量层面，而是揭示了两组重要的结构性分化。

第一组分化：限额以上企业零售额同比下滑4.9\%，远弱于整体零售增速。这意味着大型零售商的表现明显不如小企业。这背后的逻辑可能是，大企业更依赖线下渠道和传统品类，而小企业更灵活，能够更快适应消费降级和线上化的趋势。

第二组分化更为关键：服务业生产指数

[... middle omitted ...]

意味着，与出行、体验、本地生活相关的产业链仍有支撑，而耐用消费品和传统零售则需要面对更长时间的调整期。**

5月房地产投资单月同比下滑24.3\%，较4月的-20.1\%进一步恶化。销售面积同比下降13.1\%，销售金额下降9.5\%。新开工面积降幅略有收窄至-24.6\%，但仍处于深度收缩区间。

GS在报告中的一个关键表述是：“despite recent green shoots in large cities”。这意味着，尽管总量数据依然低迷，但一线城市确实出现了边际改善的迹象。这一点与市场普遍感知的“一线城市二手房成交量回暖”一致。但报告也同时指出，NBS和民间数据均显示，房价下行压力仍在持续，主要集中在低线城市。

*房地产的真正问题，不在于一线城市能否企稳，而在于三四线城市的库存去化周期。** GS并未在本次报告中给出明确的库存数据，但逻辑上，低线城市的房价下行会进一步抑制购房意愿，形成负反馈循环。这也是为什么政策端虽然在放松限购、降低首付，但效果仍未充分传导。

\subsection{[R014] GS：市场低估了全球LNG供给的结构性改善，而非地缘政治风险}
{\small\color{refgray}归类：能源与大宗：库存重建锚定油价，铜关税赌局6月底见分晓}

GS：市场低估了全球LNG供给的结构性改善，而非地缘政治风险

欧洲天然气市场刚刚经历了一个关键的转折点。美伊协议为霍尔木兹海峡的能源流动恢复铺平了道路，这直接决定了欧洲今年冬季的天然气储存补充前景。但GS这份研报的真正价值，并不在于它对地缘政治事件的判断——毕竟协议的达成与落空都存在变数——而在于它揭示了一个被当前价格所忽视的结构性变化：全球除卡塔尔外的LNG供给正在以超出预期的速度增长，这一变化正在从根本上改变欧洲天然气市场的供需平衡。

这份报告最值得关注的判断是：即便霍尔木兹海峡的恢复时间晚于预期，欧洲天然气价格在2026年下半年仍将稳定在41欧元/兆瓦时左右，而2027年将进一步降至30欧元/兆瓦时。这背后是两个被市场低估的变量：一是非卡塔尔LNG供给的大幅增长，二是市场对欧洲储气库填充目标的主动下调。这两个变量正在重塑欧洲天然气市场的定价锚。

1. 非卡塔尔LNG供给的意外增长正在抵消霍尔木兹中断的紧缩效应

GS在报告中明确指出，全球除卡塔尔外的LNG供给增长超出了他们此前的预期。他们将2026年6月全球除卡塔尔外的LNG供给预测从此前水平上调了10百万吨/年，达到384百万吨/年。这一上调主要来自几个来源：阿曼、马来西亚和挪威等地的维护作业明显减少；尼日利亚的原料气供应出现强劲反弹；以及俄罗斯受制裁的LNG出口量增加。

这一判断的意义在于：市场此前对欧洲天然气紧缺的担忧，很大程度上建立在对卡塔尔LNG供给中断的线性外推上。但GS的数据表明，全球LNG供给的韧性远比市场想象的更强。当卡塔尔的供给出现缺口时，其他地区的产能正在通过维护优化和产能释放来弥补这一缺口。这不是一次性的应急反应，而是全球LNG供给体系在价格信号下的自我调节。

对于产业决策者而言，这意味着欧洲天然气市场的供给弹性正在提升。即便霍尔木兹海峡的恢复进程出现反复，其他地区的供给增量也能在相当程度上缓冲冲击。当然，这种缓冲存在上限——GS也指出，如果霍尔木兹中断持续到2027年，TTF价格可能需要超过100欧元/兆瓦时才能通过抑制亚洲需求来重新平衡市场。但就当前情景而言，非卡塔尔供给的增长已经显著降低了极端情景的概率。

GS报告中另一个被市场忽视的关键判断是：TTF近期的定价表明，市场并未像此前预期的那样推动储气库填充至80\%的满度，而是似乎满足于70\%-75\%的填充范围。这不是市场对供给紧张的无奈接受，而是一种主动的、基于对冬季供给信心增强的策略调整。

GS的分析逻辑是：随着全球LNG供给的增长，欧洲冬季的天然气可得性得到改善，这使得市场对储气库的依赖程度下降。具体而言，他们测算70\%-75\%的储气库填充率仍足以支持一个10年平均气温的冬季后，储气库剩余量接近30\%，这个水平足以抵御一个标准差以上的寒冷冬季（相当于西北欧储气库容量的12\%）。

这一判断的深层含义是：欧洲天然气市场的风险定价正在从“供给恐慌”模式转向“供给充裕”模式。当市场相信冬季有足够的LNG现货可以补充时，对储气库的“战略储备”需求就会下降，进而降低对价格的支撑。这也是GS维持2026年下半年TTF价格预测在40-42欧元/兆瓦时的核心原因——这个价格水平正好处于天然气相对煤炭具有成本优势的区间，能够持续激励天然气对煤炭的替代需求，从而防止储气库过度拥挤。

对于投资者而言，需要关注的是这一填充目标调整的可持续性。如果全球LNG供给增长持续超预期，市场可能进一步下调填充目标，从而对价格形成持续的下行压力。反之，如果供给增长出现意外中断，市场可能重新回到80\%以上的填充目标，推动价格上行。

GS对2027年的展望可能是这份报告中最具前瞻性的部分。他们预计，在更高的2027年3月储气库剩余量和全球LNG供给同比增长47百万吨/年（至488百万吨/年）的共

[... middle omitted ...]

峡的流动能够在未来几周内正常化。如果这一假设成立，那么欧洲天然气市场将从2027年开始进入一个供给充裕的新平衡，价格中枢将显著下移。但即便在霍尔木兹恢复慢于预期的情景下，GS也认为2027年TTF均价不会超过65欧元/兆瓦时——这仍然低于当前远期曲线所隐含的水平。

尽管GS的报告提供了大量有价值的数据和判断，但有一个关键变量并未得到充分讨论：亚洲LNG需求的价格弹性。

GS在分析中假设，如果欧洲需要抑制亚洲需求来平衡市场，TTF需要达到100欧元/兆瓦时以上。但这一假设的基础是亚洲需求在当前价格水平下具有较强的刚性。然而，如果亚洲主要进口国（尤其是中国和印度）在更高价格下能够更快地转向煤炭、可再生能源或国内天然气生产，那么欧洲所需的价格水平可能低于GS的估计。

反过来，如果亚洲经济复苏超预期，或者可再生能源装机进度低于计划，那么亚洲需求的价格弹性可能比GS假设的更小，从而使得欧洲即使在高价下也难以有效吸引LNG现货。这一不确定性在报告中并未得到充分的敏感性分析。

\subsection{[R015] NOM：五月经济数据印证了一个被市场低估的核心矛盾——供给侧的约束比需求侧更顽固}
{\small\color{refgray}归类：地产与消费：一线城市结构性再定价，消费疲软拉长政策等待期}

NOM：五月经济数据印证了一个被市场低估的核心矛盾——供给侧的约束比需求侧更顽固

五月经济数据出炉，零售销售同比转负，固定资产投资大幅下滑，工业增加值反弹微弱。市场习惯性地将这些数字解读为“内需不足”，并等待新一轮刺激政策。NOM这份研报提供的关键洞察是：真正的问题不是需求不够，而是供给端出现了多重结构性约束——从石油供应中断到“反内卷”政策对产能的主动压制，这些约束正在改变经济的运行逻辑。投资者若只盯着需求侧的刺激预期，可能会错过更深刻的定价变化。

以下是我们从这份报告中提炼出的五个关键层次，每一个都在回答同一个问题：为什么五月数据比表面数字更值得警惕。

1. 五月数据的核心特征是“需求弱、供给更紧”，而非简单的经济放缓

零售销售同比转负至-0.6\%，固定资产投资同比暴跌至-10.7\%，这两个数字足以让市场紧张。但NOM的分析指出，这两条线背后的驱动力并不相同。零售的疲软更多来自消费者信心低迷和以旧换新政策退坡后的“还债效应”，而固定资产投资的崩塌则叠加了政策主动压缩产能的影响——尤其是“反内卷”运动在制造业投资上的持续压制。这意味着，即使未来需求有所回暖，供给端也可能因为政策约束和原料短缺而无法同步响应。这种“供给刚性”才是当前经济运行中最容易被忽视的变量。

2. 石油供应冲击对工业生产的传导，远比市场预期的更深、更持久

五月原油加工量同比下滑9.1\%，创下多年新低。NOM明确指出，这不仅仅是短期事件，而是中东能源供应中断后，中国工业体系暴露出的结构性脆弱。硫磺短缺已经传导至硫酸产量下滑，进而影响整个化工产业链。乙烯、化学纤维等关键中间品的产出增速连续数月下滑。这意味着，即使下游需求稳定，上游原料的“断供”也会迫使部分工业活动被动减速。市场对“AI投资拉动工业”的叙事过于乐观，而忽略了能源和化工这些基础环节正在承受的真实压力。这个链条的修复，可能比任何需求刺激政策都来得更慢。

3. 固定资产投资中，基础设施投资的“塌陷”暴露了财政效率问题，而非资金不足

基础设施投资同比下滑10.8\%，远比市场预期的-3.7\%要差。NOM提供了一个关键判断：这不是因为缺钱。五月地方政府专项债发行量与去年同期基本持平，资金供给并未收缩。真正的问题在于资金到位后未能及时转化为实际支出——即财政存款沉淀，而非项目开工。这指向的是项目审批节奏和地方执行效率的瓶颈，而非政策意愿的缺失。对于投资者而言，这意味着短期内不宜押注基建投资快速反弹，除非中央层面明确放松项目审批或要求加速资金拨付。基础设施投资的下滑，更多是制度性摩擦的体现，而非周期性需求的缺失。

4. 房地产市场的分化已经演变为“一线城市局部企稳，全国继续下沉”的二元结构

五月新房价格环比下跌0.20\%，二手房价格下跌0.26\%，跌幅扩大。但一线城市二手房价格仍保持正增长，北京、上海、广州、深圳均录得月度环比上涨。NOM的Iceberg指数显示，全国二手房挂牌价在五月继续下跌0.3\%，六月前两周跌幅甚至扩大至0.4\%。这组数据揭示了一个清晰的图景：AI驱动的财富效应和人口流入主要集中在少数“聪明城市”，而全国大部分中小城市的房地产市场仍在加速下行。市场对“房地产触底”的讨论过于笼统，真正的底部可能只存在于一线城市。对于布局房地产相关资产的投资者而言，区分“城市级别”和“区域分化”比判断“整体见底”更有实际意义。

5. 出口的繁荣并未有效传导至制造业投资和工业产出，政策驱动的产能约束正在“切断”传统传导链条

五月出口增速高达19.4\%，但制造业投资同比下滑4.2\%，工业增加值也仅微增0.4个百分点。NOM的分析暗示了一个反直觉的结论：在“反内卷”政策的约束下，企业即使看到强劲的外部订单，也无法或不愿扩大产能。这打破了“出口好—投资强—就业稳”的传统逻辑。光

[... middle omitted ...]

定性。如果中央层面出台新的项目审批加速机制或地方债务化解方案，基础设施投资可能在两到三个季度内出现反弹。反之，如果执行层面的摩擦持续，那么今年下半年的经济增长将面临更大的下行压力。这个问题的答案，将直接影响我们对政策效果和资产定价的判断。欢迎来星球微信群里继续讨论，我们会分享更多关于财政效率、项目审批动态以及下半年经济走势的跟踪分析。

Personal reading notes and learning share only. Not investment advice.

某外资投行最新研报显示，5月中国经济数据整体偏弱，内需进一步恶化。零售销售同比-0.6\%，为2022年底以来首次转负；固定资产投资同比-10.7\%，远低于预期。工业生产小幅反弹至4.5\%，但动力不足。

1️⃣ 零售“冷”得明显 汽车零售额同比-16.1\%，消费信心低迷 家电、家具、体育娱乐零售持续负增长 餐饮收入增速降至0.6\%，大餐厅甚至转负 石油产品零售额虽改善，但实际消费可能大降30\%

\subsection{[R016] Citi：市场低估了中国AI供应链在下一代GPU架构中的份额跃迁}
{\small\color{refgray}归类：AI/算力：智能体时代CPU复兴，中国供应链份额跃迁}

Citi：市场低估了中国AI供应链在下一代GPU架构中的份额跃迁

这份Citi研报的核心判断并非“AI需求仍在增长”这类共识，而是一个更具体、更具投资含义的结论：中国AI基础设施供应商在英伟达GB200架构中获得的份额有限，但在GB300和Vera Rubin架构中，他们将实现显著的份额跃迁。这不是一个渐进式的改善，而是一个由技术升级、响应速度和产能承诺共同驱动的结构性拐点。对于关注中国制造业升级和AI产业链定价权的投资者来说，这是当前最值得深入理解的变化。

报告基于2026年6月9日至12日的一场产业调研，覆盖了13家中国工业与机器人领域公司。调研中最受投资者关注的公司几乎全部集中在英伟达供应链内，包括PCB钻孔设备商大族数控、全球最大覆铜板（CCL）厂商生益科技、全球最大钻头公司鼎泰高科，以及一家英伟达认证的电源供应商。这些公司过去在GB200中的“钱包份额”远低于日韩台同行，但调研反馈一致指向一个转折点：下一代架构将成为中国供应商证明自身价值的主战场。

这份报告的价值不仅在于揭示了这一趋势，更在于它提供了足够细致的产业链验证——从覆铜板、玻纤布到PCB设备，每一环节的产能瓶颈、技术壁垒和盈利前景都有具体数据支撑。以下是我们从报告中提炼出的五个关键洞察。

1. 中国AI供应商的份额跃迁不是价格战的结果，而是技术曲线和响应速度的胜利

一个常见的误解是，中国供应商的竞争力主要来自价格。但这份报告纠正了这一看法。调研中，生益科技、大族数控、Grace Fabric等公司明确指出，他们赢得更多份额的核心因素依次是：技术曲线从全球视角看已实现升级、对英伟达快速变化的架构有更敏捷的响应时间、以及对产能建设有更坚定的承诺。价格竞争被排在最后。

这一判断有重要的行业含义。它意味着中国AI供应链的竞争力正在从“成本优势”向“系统集成与快速迭代优势”转变。对于英伟达这样产品迭代周期极短的客户而言，供应商能否在架构变更时迅速调整工艺、交付样品、并承诺大规模产能，远比单纯的报价重要。这正是中国供应商在GB200时代积累的经验——他们作为后来者，在经历了地缘政治压力和缺乏台积电供应链接触的劣势后，反而锻炼出了更强的适应能力。

因此，投资者不应简单地将中国AI供应链视为“替代者”，而应将其视为“跟随迭代者”。他们在GB300和Vera Rubin上的份额增长，本质上是对其技术能力和响应速度的重新定价。

2. AI玻纤布是整个CCL/PCB供应链中最紧的瓶颈，且价格弹性远超预期

报告中最令人印象深刻的细节之一，是对AI玻纤布（AI-fabric）产能短缺的量化描述。生益科技和Grace Fabric均指出，整个CCL/PCB供应链中产能最紧张的部分是AI玻纤布，特别是T-glass和Low Dk gen 2。这不是一个笼统的判断，而是有具体数据支撑的：Grace Fabric的AI玻纤布年产能预计将从2025年的约500万米，飙升至2026年的2500万米，再到2027年的至少3300万米，但即便如此，市场仍处于供不应求状态。

更关键的是价格弹性。报告显示，2026年前六个月，KBL（建滔集团）的电子级玻纤布ASP（平均售价）已上涨60\%-95\%。这并非一次性事件，而是结构性供需失衡的结果。由于纺机供应有限，新增产能的扩张速度受到物理约束，这意味着价格上行周期可能比市场预期的更长。

对于投资者而言，这意味着在AI供应链中，上游材料环节（特别是玻纤布）可能比中游的CCL或PCB设备拥有更强的定价权和盈利弹性。报告直接点明了这一偏好：在AI基础设施链条中，Citi的偏好顺序是KBL大于大族激光大于生益科技。KBL目前2027年PE仅为16倍，对应3年EPS CAGR高达90\%，而同类公司Grace Fabric的PE高

[... middle omitted ...]

泰已从去年开始向一家美国领先机器人公司出货，且量级可观。两家公司均预期Optimus Gen 3今年的产量在1万至2万台之间，明年可能增长10倍至10万台以上。基于此，Citi预计恒立液压的机器人业务收入贡献将从2026年的不到1\%，提升至2027年的约4\%，且存在上行空间。

这一判断的深层含义是：在人形机器人从概念走向量产的初期阶段，具有量产能力和客户验证的零部件供应商，比整机厂商更有可能率先实现盈利和估值兑现。整机厂商虽然品牌效应更强，但面临的技术整合难度、资本开支压力和竞争格局恶化风险都更大。对于产业投资者而言，这是一个重要的“抓上游、放下游”的阶段性策略。

4. 挖掘机复苏是结构性的，但整个工程机械行业的隐忧不可忽视

报告对工程机械行业的分析呈现了鲜明的分化。一方面，挖掘机领域出现了可持续的复苏信号。恒立液压的挖掘机零部件产线自年初以来一直满产，主要驱动力来自矿山对大型挖掘机的需求、恒立在挖掘机泵阀市场的份额提升，以及卡特彼勒等外资客户份额的进一步扩大。

\subsection{[R017] GS：中国房地产市场的真实信号并非全面回暖，而是核心城市的结构性再定价}
{\small\color{refgray}归类：地产与消费：一线城市结构性再定价，消费疲软拉长政策等待期}

GS：中国房地产市场的真实信号并非全面回暖，而是核心城市的结构性再定价

市场参与者对5月房地产数据的解读，很容易滑入两个极端：要么被全国层面的同比下滑数字吓退，要么被一线城市的环比改善故事所吸引。GS最新发布的5月中国房地产月度追踪报告提供了一组值得细看的数据，其核心判断并非简单的“回暖”或“继续探底”，而是一个更为精细的结构性叙事——全国总量的疲软掩盖了核心城市正在发生的资产价格再定价过程，而这种再定价的可持续性，取决于供给侧的深度调整而非需求侧的短暂脉冲。

这份报告最值得关注的信号，是一线城市已经连续3至4个月录得二手房价格环比正增长，同时二手房挂牌量同比连续第三个月下降。这两个指标的同步出现，在近两年的市场周期中并不常见。它意味着核心城市的二手房市场，正在从“降价换量”的存量博弈，转向“量稳价升”的供需再平衡。对于产业决策者和高净值投资者而言，真正需要追问的不是“市场到底了没有”，而是“哪些城市的资产正在被重新定价，以及这种定价能走多远”。

1. 一线城市的价格韧性并非偶然，而是供给收缩与需求结构升级共同作用的结果

GS报告显示，5月70城新房价格环比平均下降0.2\%，与4月、3月基本持平；而二手房价格环比降幅则从4月的0.2\%扩大至0.3\%。表面上看，二手房市场似乎比新房更弱。但一线城市的表现提供了完全相反的图景：5月一线城市新房价格环比上涨0.2\%，已是连续第4个月正增长；二手房价格环比上涨0.3\%，是连续第3个月正增长。深圳和上海几乎所有产品类型的二手房成交价格都在改善，北京的高端项目也跑赢大盘。

这里需要回答一个“为什么”的问题。一线城市的价格韧性，不能简单归结为“有钱人还在买房”。GS报告指出，新房价格受到产品结构升级的支撑——核心区域的新房供应本就有限，而开发商在当前环境下更倾向于推出高品质项目，这类项目的单价和去化速度都明显优于普通产品。深圳南山区一宗地块的土地成本达到每平方米10.9万元，杭州上城区一宗地块由绿城和杭州地铁集团联合以21亿元拿下——这些高价地项目未来入市的价格锚点，反过来也会支撑周边二手房的定价预期。

更深层的结构性变化发生在二手房供给侧。报告显示，100城二手房挂牌量同比连续第三个月下降，同时成交周转速度在加快。这意味着，过去两年积累的“抛压盘”正在被市场消化，而新增挂牌的意愿在减弱——业主在经历了深度调整后，对以大幅折价出售的抵触情绪上升。供给侧的主动收缩，是价格企稳的必要条件，这一点在一线城市已经初步成立。

5月全国新房销售面积同比下降13\%，销售额同比下降9\%，均低于GS的预期。新开工和竣工面积分别同比下降25\%和20\%，也弱于预期。这些数字很容易让人得出“市场仍在恶化”的结论。但GS报告在给出这些数字的同时，还提供了另一组关键数据：15城二手房成交量同比增长8\%，高于预期，且增速较4月的4\%有所加快。

这意味着什么？全国新房销售的下滑，很大程度上反映的是供给端的问题——开发商现金流紧张、新开工持续萎缩、可售房源集中在非核心区域。而二手房市场的活跃，则反映了真实居住需求在价格调整到位后的释放。两者并非同一市场，而是同一市场内部的分化。对于投资者而言，把新房和二手房数据简单加总去判断市场方向，可能会错过结构性的机会。

这种分化还体现在租赁市场。50城平均租金水平仍在环比下降，但一线城市的租金已经连续三个月环比上涨，且表现优于2024年和2025年同期。上海、深圳、天津、乌鲁木齐是3至5月表现最好的城市。租金的企稳，是资产价格可持续修复的前提——没有租金支撑的价格上涨，最终都会回到原点。

3. 开发商的行为模式正在从“规模竞赛”转向“利润优先”，这是市场出清的关键信号

报告跟踪的6家覆盖开发商在5月的土地投资强度出现分化：保利发展和绿城加快了拿地节

[... middle omitted ...]

身就在上升——深圳南山地块每平方米10.9万元的土地成本，意味着未来项目的售价必须达到相当高的水平才能实现盈利。如果核心城市的房价上涨预期无法兑现，这些“优质地块”也可能成为新的风险点。

4. GS对6月的预测暗示了一个关键假设：核心城市的价格修复能否扩散到二线城市

报告对6月给出了明确的预测：新房和二手房价格环比降幅都将收窄，一线城市的修复势头将延续，而关键二线城市可能转为环比正增长。同时，新房销售面积和金额预计将保持低两位数百分比的同比下降，二手房成交量预计同比增长中单位数百分比。

这里隐含了一个重要的假设：二线城市正在跟随一线城市的节奏进入价格企稳阶段。5月的数据已经显示，二线城市的新房和二手房价格环比降幅均为0.1\%和0.2\%，与4月基本持平，没有进一步恶化。70城中，有13个城市在5月录得二手房价格环比正增长或持平，其中10个城市已经连续3个月保持这一状态。这些城市是否包括上海、深圳、天津、乌鲁木齐之外的其他城市？报告没有明确点名，但趋势是值得跟踪的。

\subsection{[R018] GS：新能源车周订单回落，但渗透率突破67\%才是市场真正该关注的信号}
{\small\color{refgray}归类：新能源车：渗透率突破67\%，竞争逻辑从抢增量转向抢存量}

GS：新能源车周订单回落，但渗透率突破67\%才是市场真正该关注的信号

6月第二周，中国新能源车市交出了一份看似矛盾的答卷：周订单环比下降13\%，同比下降14\%，但零售与批发的新能源渗透率却双双突破67\%。对于习惯用周订单增速来线性外推行业景气度的投资者而言，这个组合容易引发困惑。

GS本周发布的《中国新能源汽车周刊》提供了理解这一现象的关键框架。报告的核心洞察在于：订单的短期波动主要受到5月底新品集中上市脉冲退潮的影响，属于预期内的周期性回落；而渗透率在6月初持续上行至历史高位，则指向一个更深层的结构性变化——燃油车市场份额正在被不可逆地侵蚀，且这一过程并未因个别品牌的订单波动而放缓。

市场真正需要关注的，不是“下周订单会不会反弹”，而是“当渗透率突破三分之二后，行业竞争的逻辑将从抢增量转向抢存量，这意味着什么”。GS的数据显示，这一临界点已经到来。

GS追踪的9家主要新能源品牌在6月第二周（6月8日至14日）合计订单为15.36万辆，环比下降13\%，同比下降14\%。表面上看，这一数据似乎预示着5月以来的回暖势头正在受阻。

但报告明确指出，这一回落是“coming off the initial pulse of new model launches in late May”——即5月底新品上市带来的首波脉冲消退后，订单回归正常节奏。5月最后一周（W22），合计订单曾达到20.78万辆的年内高位，随后W23回落至17.71万辆，W24进一步降至15.36万辆。这是一个典型的“发布周冲高、随后回归”的节奏，在汽车行业的新品周期中属于正常现象。

从同比角度看，虽然整体下降14\%，但结构分化显著。蔚来、理想、小鹏、小米等品牌仍录得同比正增长，其中蔚来同比增幅高达282\%。这进一步说明，订单回落并非全行业需求疲软，而是个别基数效应和产品周期节奏差异的体现。

GS在报告中使用了“defensive growth”来描述吉利和比亚迪的周订单表现。吉利银河与极氪品牌合计订单环比增长8\%，比亚迪王朝与海洋系列环比仅下降4\%，特斯拉环比下降8\%。在整体订单下滑13\%的背景下，这三家头部企业表现出明显的抗跌性。

这一现象的含义值得深挖。当行业处于新品脉冲退潮期，订单的稳定性实际上反映了品牌的产品力壁垒和渠道韧性。比亚迪凭借庞大的产品矩阵和价格带覆盖，能够在任何市场环境下维持相对稳定的订单流入。吉利则受益于银河系列和极氪品牌的双轮驱动，在15-30万元价格带形成了有效卡位。

从年初至今的累计订单来看，蔚来以同比93\%的增速领跑，华为智选车（HIMA）同比增长31\%，特斯拉同比增长3\%。GS的数据显示，头部品牌的份额正在扩大，而中尾部品牌在新品脉冲退潮后订单下滑更为剧烈。这意味着，2026年的新能源车市已经不再是“人人有肉吃”的增量市场，产品力和品牌力较弱的企业将面临更大的生存压力。

3. 渗透率突破67\%才是本周最值得关注的数字，燃油车进入加速退场阶段

尽管订单数据引发关注，但GS报告中最具战略意义的数字是渗透率。6月1日至7日，新能源零售渗透率达到66.7\%，批发渗透率达到67.2\%，而5月全月这两个数字分别为63\%和61.1\%。短短一周内，渗透率提升了3-6个百分点。

这一跃升的驱动力并非来自新能源订单的暴增——事实上，同期新能源零售同比仍下降14\%——而是燃油车零售量的更快萎缩。6月第一周，乘用车整体零售22.8万辆，同比下降23\%，批发20.4万辆，同比下降25\%。燃油车市场的收缩速度远超新能源，导致新能源在总量中的占比被动提升。

这意味着，新能源渗透率的上升正在从“主动增长”阶段进入“被动替换”阶段。当燃油车需求加速下滑，即使新能源订单增速放缓，渗透率仍会持续攀升。GS的数据实际上在提示一个残酷的

[... middle omitted ...]

及部分品牌主动控制折扣力度以维护品牌形象。但燃油车折扣的持续扩大，则反映出经销商去库存压力仍在加剧。以宝马、奔驰、本田等品牌为例，GS数据显示其终端折扣长期维持在25\%-30\%的高位，且未见明显收窄迹象。

从品牌维度看，比亚迪的经销商折扣仅为4.27\%，远低于行业平均的7.54\%，也低于自身一周前的5.14\%。这显示出比亚迪在定价权上的优势——当竞争对手需要靠降价来维持销量时，比亚迪凭借规模效应和产品力，反而有能力收紧折扣、提升单车利润。

对于投资者而言，定价信号的分歧意味着：新能源行业正在从“价格战”阶段进入“分化战”阶段。拥有定价权的企业将率先改善盈利，而依赖降价维持销量的企业，其利润压力将持续加大。

5. 报告未完全回答的关键问题：电池级碳酸锂价格反弹的可持续性

GS报告在最后一部分提到了上游电池定价动态：电池级碳酸锂价格升至每吨17.45万元，环比上涨3.6\%，而方形LFP电芯和方形NCM电芯价格周环比持平。这是一个值得关注的信号，但报告并未深入展开。

\subsection{[R019] GS：五月新房价格跌幅收窄，但市场真正的信号藏在二手房与库存之间}
{\small\color{refgray}归类：地产与消费：一线城市结构性再定价，消费疲软拉长政策等待期}

GS：五月新房价格跌幅收窄，但市场真正的信号藏在二手房与库存之间

这份来自GS中国经济学团队的最新研报，基于国家统计局70城数据，给出了一个看似温和但内部结构高度分化的判断：五月新房价格环比跌幅收窄，一线城市甚至出现环比正增长。但如果你只读到“企稳”两个字，可能就错过了这份报告真正想传达的信息。

GS团队在报告中反复强调一个关键区分：70城数据仅覆盖新房市场。而二手房市场的价格，根据NBS及第三方平台的数据，过去一年跌幅在5\%至10\%之间。这意味着，当前市场的价格信号存在明显的“双轨”特征——新房市场的边际改善，并不能简单等同于整体房价的触底。

对于产业决策者和高净值读者而言，真正需要追问的不是“市场是否已经见底”，而是“在总量数据改善的表象下，哪些结构性变量正在重塑定价逻辑”。这份报告提供了几个值得深挖的线索。

1. 一线城市新房价格环比转正，但驱动力高度集中且依赖局部政策

五月数据最亮眼的部分是一线城市新房价格环比年化上涨1.4\%，较四月的0.7\%明显加速。其中深圳表现最为突出，环比年化涨幅达到3.9\%，远超四月的1.6\%。上海、杭州、合肥等城市同比依然保持正增长。

然而，GS在报告中明确指出，近期的政策宽松仍然是“局部性的”。广州的收储计划、各地对住房公积金使用的扩大，这些措施并未形成全国性的、系统性的需求提振。换言之，一线城市的回暖更多是局部政策刺激与个别城市供需结构改善的结果，而非全国房地产市场的普遍性复苏。

这意味着，投资者在关注一线城市数据时，需要追问：这种环比改善的可持续性如何？如果政策刺激的边际效应递减，这些城市能否依靠自身的基本面维持价格稳定？目前报告并未给出明确答案，但隐含的判断是，这种改善的广度与深度仍有待观察。

2. 二三线城市跌幅收窄，但“收窄”不等于“止跌”，库存压力才是核心变量

五月二线城市新房价格环比年化跌幅从-2.7\%收窄至-1.1\%，三线城市从-2.9\%收窄至-1.8\%。从方向上看，这是一个积极信号，表明最剧烈的下跌阶段可能已经过去。

但GS同时提供了一个更值得关注的指标：主要城市的库存月数（可售面积除以12个月滚动销售面积）从28.9个月下降至28.5个月。虽然降幅微小，且主要由二线城市驱动，但这表明去库存过程仍在进行中。

这里有一个关键的“所以呢”：库存月数依然处于历史高位。28.5个月意味着即使在没有新增供应的前提下，现有库存也需要超过两年才能消化完毕。在这样的库存压力下，开发商降价促销的动力依然强劲。因此，二三线城市价格跌幅的收窄，更可能是政策托底与开发商主动降价之间的暂时平衡，而非需求端根本性的好转。

对于关注区域市场的读者而言，需要建立一个新的观察框架：不再仅仅盯着价格环比变化，而是将“库存月数的变化速度”作为判断市场是否真正触底的前瞻指标。只有当库存月数出现连续且显著的下降，价格的企稳才更具说服力。

3. 高频交易量数据与价格信号出现背离，市场正处于“量稳价跌”的博弈阶段

GS的高频跟踪数据显示，30城新房交易量在五月和六月初基本与去年同期持平。这是一个容易被忽视但非常重要的信息。

通常，在房地产市场下行周期中，交易量的企稳往往被视为价格见底的前兆。然而，当前的情况是“量稳”但“价未稳”——新房价格虽然跌幅收窄，但仍在下跌；二手房价格跌幅更深。这种背离意味着什么？

一个合理的解释是，当前的市场正处于“以价换量”的博弈阶段。开发商和二手房业主通过主动降价来换取流动性，从而维持了交易量的稳定。但买家的预期并未根本扭转，他们仍在等待更低的价格。这种“量稳价跌”的状态，本质上是一种脆弱平衡。一旦外部环境（如收入预期、信贷政策）出现变化，交易量可能再次萎缩，价格也可能加速下跌。

这份报告没有直接预测这种平衡何时会被打破，但它提供了一个清

[... middle omitted ...]

来房价的长期信心。GS在报告中并未测算这些政策对需求的拉动效果，也没有讨论是否需要更大规模的财政或货币刺激来填补这个缺口。

这是一个值得所有读者深思的问题。从历史经验看，房地产市场的复苏从来不是一个线性过程。政策可以托底，但无法创造需求。只有当居民的收入预期改善、对资产价格的长期信心恢复时，需求才会自发回归。目前，我们看到的更多是政策驱动的边际改善，而非信心驱动的内生复苏。

5. 给决策者的观察框架：从“看价格”转向“看结构”和“看流动性”

基于GS这份报告的核心数据与逻辑，我们建议读者建立一个新的观察框架，以取代过去简单的“涨跌判断”。

第一，区分“新房市场”与“二手房市场”。新房市场的价格受开发商定价策略、推盘节奏和局部政策影响较大，容易失真。二手房市场的价格更能反映真实的供需关系和业主的预期。GS报告明确指出，二手房价格跌幅远大于新房，且部分第三方平台数据已经暂停发布（如贝壳在2023年10月暂停发布二手房价格数据），这本身就说明市场定价的透明度在下降。

\subsection{[R020] 摩根斯坦利：中国房地产市场真正的拐点信号，不是成交量的反弹，而是成交结构的断裂}
{\small\color{refgray}归类：地产与消费：一线城市结构性再定价，消费疲软拉长政策等待期}

摩根斯坦利：中国房地产市场真正的拐点信号，不是成交量的反弹，而是成交结构的断裂

这份摩根斯坦利最新一期中国房地产周度数据库追踪报告，覆盖截至6月14日的数据，释放的信号远比表面数字复杂。市场在寻找“底部确认”的信号，但这份报告揭示的，是一个更值得警惕的结构性分化：新房与二手房、一线与低线城市之间的表现正在走向彻底断裂。

报告最核心的判断并非成交量本身，而是这组数据合在一起指向了一个结论：中国房地产市场的定价机制正在从“新房锚定”转向“二手房锚定”，而这一转变对开发商、投资者和地方财政的含义，远未被市场充分定价。

1. 新房成交的“脉冲式反弹”正在退潮，同比增速从+21\%骤降至-1\%，市场不应将其视为常态波动

上周50城新房成交同比下滑1\%，而此前一周还是+21\%的增长。两周之内，增速落差高达22个百分点。这不是正常波动，而是一个清晰的信号：此前几周的反弹更多依赖于积压需求和政策窗口期的集中释放，而非需求面的根本改善。

从全年累计看，新房成交同比仍为-12\%。一线城市表现最弱，同比下降9\%，而此前一周还是+3\%。三线城市同样从+27\%跌至-4\%。这种同步回落，说明反弹并非个别城市的独立行情，而是整体市场缺乏持续动力。

这组数据意味着什么？如果市场参与者仍用“政策刺激后会有持续反弹”的历史经验来预测未来，可能会持续误判。本轮周期的特殊性在于，需求端的修复速度远慢于供给端的收缩，而政策工具的空间也已大幅收窄。

2. 二手房成交的韧性正在重塑市场定价权，这是报告最容易被忽视的结构性信号

与新房形成鲜明对比的是，10城二手房成交同比增长4\%，全年累计+6\%。一线城市二手房成交同比+3\%，二线城市+3\%。尽管增速较前一周的+34\%和+17\%有所回落，但整体韧性明显强于新房。

这不是一个简单的“二手房比新房好卖”的故事。背后的结构性含义是：二手房正在成为市场定价的基准。当二手房成交量持续稳定、甚至超过新房时，价格发现机制就从开发商定价转向买卖双方博弈。这意味着，开发商未来定价的空间将被二手房成交价牢牢锁定。

对于那些仍试图通过降价促销来拉动新房销售的企业，这个信号尤其值得警惕：不是降价幅度不够，而是定价权已经转移。二手房市场的价格将成为购房者“心理锚点”，新房价格若不能与之匹配，即使降价也难以放量。

3. 销售去化率从100\%暴跌至48\%，这是一个不应被忽视的流动性警示

上周总体销售去化率为48\%，而此前一周为100\%。一线城市从100\%降至45\%，二线城市从100\%降至51\%。去化率腰斩，意味着推盘量在增加，但实际成交并未同步跟上。

这里需要特别指出：去化率100\%本身就是一个非常态数据，可能受单周推盘结构或统计口径影响。但48\%的去化率，结合新房成交同比下滑，说明市场承接力正在快速减弱。对于开发商而言，这意味着库存压力可能在接下来的几周内重新累积。

报告没有展开的一个关键问题是：去化率的下降是暂时的供应端节奏问题，还是需求端持续疲软的反映？如果是后者，那么开发商的现金流压力将再次面临考验，尤其是那些依赖高周转模式的民营企业。

4. 中原领先指数和中介指数同步走弱，二手房价格下行压力尚未触底

中原领先六城二手房挂牌价指数为17.4\%，较前一周的17.5\%小幅下降。一线城市中介指数从54.4降至53.2。这两个指数虽然变动幅度不大，但方向一致，且处于历史低位附近。

报告没有给出价格变动的绝对数值，但指数方向已经足够清晰：二手房价格仍在缓慢下滑，尚未出现企稳迹象。结合前面提到的去化率下降，可以推断：卖方正在通过降价来维持成交量，但买方仍在观望。

这里有一个尚未被充分讨论的二阶效应：如果二手房价格持续下行，将直接冲击银行的抵押物价值评估体系，进而影响开发贷和按揭贷款的审批条件。这不是一个线性风

[... middle omitted ...]

产负债表受损程度远超预期；房价下行预期已经形成自我实现的循环；政策工具本身的空间（如利率、首付比例）已经接近极限。但这些假设需要更多数据验证。

报告中还有一个值得关注的空白：没有讨论不同城市间的库存结构差异。一线城市库存水平相对健康，但低线城市可能面临严重过剩。如果未来政策刺激只能惠及高线城市，而低线城市继续恶化，那么“结构性分化”将演变为“系统性断裂”。

6. 投资者的观察框架应该从“何时触底”转向“谁能在结构分化中生存”

第一，放弃“整体市场底部”的概念。新房和二手房、一线和低线城市、国企和民企，将呈现完全不同的走势。整体数据均值化的意义正在下降。

第二，将“去化率”和“二手房挂牌价指数”作为领先指标，而不是成交量。成交量可能因为推盘节奏而出现脉冲式波动，但去化率和价格趋势更能反映真实供需关系。

第三，关注开发商的“定价权”而非“规模”。在二手房定价主导的市场中，那些能够通过产品力、地段和服务来维持溢价能力的开发商，将比单纯追求规模的企业更具抗风险能力。

\subsection{[R021] Bernstein：数据中心互联的“网络效应”才是真正的护城河}
{\small\color{refgray}归类：AI/算力：智能体时代CPU复兴，中国供应链份额跃迁}

Bernstein：数据中心互联的“网络效应”才是真正的护城河

市场对数据中心行业的关注，多数时候停留在“电力够不够”和“机柜租金涨不涨”这两个维度。但一份来自Bernstein的最新深度报告，提出了一个被严重低估的判断：在AI时代，真正决定一家数据中心运营商长期价值的，不是它拥有多少兆瓦的电力容量，而是它能在多大程度上把物理空间转化为一个“网络枢纽”——也就是互联能力。

这份报告的核心结论相当直接：互联业务（Interconnection）是一个毛利率高达70\%-90\%、客户粘性极强、且具有“网络效应”的利润层。而在这个赛道里，Equinix（EQIX）的领先优势比市场普遍认知的还要大，其护城河不仅深厚，而且正在被AI需求持续加固。这不仅仅是一份关于某家公司的研报，它实际上提供了一个重新评估整个数据中心资产定价逻辑的新框架。

1. 互联不是增值服务，而是数据中心从“仓库”变成“枢纽”的关键

理解这份报告的价值，首先需要跳出传统认知。很长一段时间里，数据中心被视为“数字世界的仓库”——企业租用空间和电力，放置自己的服务器。互联（cross-connects）只是这个仓库里的一个附加功能，用于连接不同网络，避免流量经过公共互联网，以降低延迟、提升安全性。

但Bernstein指出，这种理解已经过时。互联正在从“购买数据中心的考虑因素”演变为“数据中心自身的一种高利润产品”。其商业逻辑非常清晰：一根物理或虚拟的交叉连接线，对运营商来说成本极低，但一旦客户部署了生产流量，这根线就变得极度“粘性”，几乎不可能被替换。更重要的是，当同一个数据中心园区内聚集的网络提供商越多，每一根新增的连接线的价值就越高——这是一个典型的“网络效应”。

这意味着，那些成功从“卖电力和空间”转型为“卖连接和生态”的运营商，能够对同一块物理地产进行多次变现。他们不仅赚取机柜租金，还通过交叉连接、端口和虚拟链路收取高利润的经常性收入。这种模式的毛利率远高于核心的批发托管业务，是真正意义上的利润放大器。

2. Bernstein新数据库揭示：Equinix的密度优势远超想象

为了量化这种“网络效应”，Bernstein的研究团队做了一件此前市场没有做过的事：他们梳理了一个全新的专有数据库，追踪了全球645个最关键的企业级托管设施中的21000多个网络连接，并将其分为云入口、一级全球传输、内容/CDN、ISP等多个类别。

结果一目了然。Equinix拥有222个互联设施，每个站点平均拥有58个互联合作伙伴，支撑着超过51.3万个创收的互联连接（截至2026年第一季度）。相比之下，排名第二的竞争对手（美国铁塔旗下的CoreSite）只有29个设施，平均每个站点有46个提供商。而Digital Realty（DLR）虽然在总网络数量上排名第二，但其每设施平均提供商数量只有34个，这与其相对较晚才转向企业级市场的策略一致。

更惊人的数据是：全球60\%的云入口（Cloud On-Ramps）位于Equinix的设施内。在拥有3个及以上云入口的设施中，57\%属于Equinix，而第二名Digital Realty仅占23\%。Bernstein的结论非常明确：这种密度既是意义重大的，也是极难复制的。

报告还揭示了一个容易被忽视的地理维度。互联的价值并非全球统一，而是高度本地化的。在美国，Equinix的平均每条交叉连接月费约为342美元，而在世界其他地区，这个数字约为197美元。这背后的原因在于，美国市场对低延迟、高安全性的私有互联需求更为旺盛，企业愿意为此支付溢价。

但这并不意味着非美市场没有机会。Digital Realty在欧洲拥有一些异常强大的密集设施——其法兰克福枢纽拥有611个网络，甚至超过了Equinix圣保罗的586个，

[... middle omitted ...]

接利好那些能够提供丰富、低延迟互联生态的运营商。

然而，报告没有完全展开的一个关键问题是：AI工作负载对互联的需求模式，与传统企业级需求有何本质不同？传统企业的互联需求相对稳定，通常是长期合约。但AI训练任务可能具有“爆发性”和“阶段性”特征——一个训练周期内可能需要大量的临时带宽，训练结束后需求又会回落。这种模式是否能被现有的物理和虚拟交叉连接产品完美覆盖，还是会催生出全新的、按需计费的互联产品形态？这将是决定互联业务收入增长天花板的关键变量。

此外，报告虽然强调了Equinix的领先地位，但并未深入探讨“新云”（Neoclouds）——如CoreWeave、Lambda等专为AI设计的云服务商——对互联格局的长期影响。这些新云厂商是否会像传统超大规模云厂商一样，倾向于在Equinix的设施内建立入口？还是会选择与Digital Realty或其他运营商进行更深度的绑定，从而改变现有的互联权力结构？这些问题，是报告留给读者的思考题，也是未来一年最值得跟踪的趋势。

\subsection{[R022] DB：五月社融回暖但需求根基未稳，银行股的选股逻辑正在从规模转向定价权}
{\small\color{refgray}归类：银行与信贷：社融回暖但需求根基未稳，选股逻辑转向定价权}

DB：五月社融回暖但需求根基未稳，银行股的选股逻辑正在从规模转向定价权

五月的金融数据，像一杯温水——比四月的冰水好一些，但远谈不上滚烫。DB最新发布的研报显示，五月新增社融和人民币贷款分别录得2.0万亿和5200亿，较四月的低点有所回升，也略高于市场一致预期。但存量社融和贷款同比增速分别放缓至7.7\%和5.5\%，双双触及历史新低。

这份数据最值得关注的判断不是“回暖”本身，而是回暖的结构揭示了什么：政府债券仍然是社融的绝对主力，企业中长期贷款连续为负，居民贷款延续收缩。换句话说，金融体系正在靠财政和票据“撑场子”，真实的实体融资意愿并未修复。对银行投资者而言，这意味着过去依赖规模扩张的盈利模式正在被挑战，真正需要评估的变量，不再是信贷增速，而是银行在弱需求环境中能否将有限的信贷资源转化为定价权和资产质量优势。

DB在报告中明确表示，偏好具有较大企业业务敞口、经营韧性较强、股息率合理的大型银行，其首选为建设银行H股和中国银行H股。这个判断背后，是对银行股选股逻辑的一次重新定义。

五月社融2.0万亿，看似从四月6250亿的“地量”中反弹，但同比去年五月仍少增约3000亿。拆开来看，有三个信号值得深究。

第一，政府债券发行放缓。五月政府债净融资1.2万亿，低于四月的1.5万亿，说明财政发力的节奏并未加速。在社融构成中，政府债仍是最大单一贡献项，但其增速放缓意味着，如果没有更积极的财政扩张，后续社融的支撑力可能进一步减弱。

第二，企业中长期贷款继续为负。五月企业中长期贷款录得-200亿，虽然比四月的-2000亿有所收窄，但连续为负的态势表明，企业用于固定资产投资的中长期融资需求依然疲弱。这不是季节性波动，而是实体投资意愿的系统性偏弱。

第三，居民贷款连续第六个月为负。五月居民贷款减少1410亿，其中短期贷款减少840亿，中长期贷款减少570亿。虽然降幅较四月有所收窄，但负值本身说明，居民仍在主动去杠杆——消费信贷和按揭需求均未出现趋势性修复。

这三个信号合在一起指向一个判断：当前社融的回暖，更多是低基数效应和票据融资支撑下的“技术性反弹”，而非实体需求的真实改善。DB在报告中使用的措辞是“weak demand persists”，这个判断是审慎而准确的。

2. 票据融资的激增，是银行在弱需求下的被动选择，也是风险的累积

五月企业票据融资新增5570亿，同比去年五月的750亿大幅增长，占当月企业贷款增量的绝大部分。这是整份报告中一个容易被忽视但极其关键的细节。

票据融资的本质是什么？它是银行在找不到合格信贷项目时，通过买入票据来“填充”信贷规模的一种方式。票据融资的激增，通常意味着两件事：一是实体企业的真实贷款需求不足，银行被迫用票据来满足监管的信贷投放要求；二是银行的风险偏好正在下降——相比发放中长期贷款，票据融资期限短、风险低、流动性好，但收益率也低。

从银行盈利的角度看，票据融资的扩张会拉低整体贷款收益率，压缩净息差。从资产质量的角度看，虽然票据本身风险较低，但大量票据融资替代正常贷款，意味着银行对实体经济的信用敞口并未真正扩大，反而可能隐藏了信贷结构的恶化。

DB在报告中提到，票据融资的规模“still below April’s RMB1.2tn”，但同比增幅仍然显著。这个数据背后的问题是：如果剔除票据融资，五月的企业贷款增量几乎为零。银行股投资者需要关注的，不是信贷总量，而是信贷结构——票据占比越高，银行的利息收入质量和可持续性就越差。

3. 居民存款连续两个月下降，背后是“存款搬家”还是消费意愿修复？

五月新增存款1.8万亿，较四月的2700亿大幅反弹，但结构分化明显。非银行金融机构存款增加1.1万亿，财政存款增加7100亿，而居民存款减少1100亿，企业存款减少1700亿

[... middle omitted ...]

在修复，居民短期贷款应该同步回升，但五月居民短期贷款仍然减少840亿。因此，存款减少更可能是资产配置的转移，而非消费意愿的实质性改善。

对企业而言，存款下降1700亿，较四月的1.25万亿降幅大幅收窄，说明企业的现金流状况有所稳定。但企业存款仍为负增长，结合企业中长期贷款的持续为负，说明企业仍在“存钱”而非“花钱”——投资意愿和扩张信心尚未恢复。

4. M1增速回升至5.5\%，但需要警惕这个指标的“失真”风险

五月M1同比增速为5.5\%，较四月的5.0\%有所回升；M2增速稳定在8.6\%。M1的回升，通常被解读为企业活期存款增加、交易性货币需求改善，是经济活跃度上升的信号。

但在这个数据背后，需要留意一个结构性问题。M1的统计口径包括流通中现金和企业活期存款。当前M1增速的回升，部分原因可能是财政存款向企业存款的转移——五月财政存款增加7100亿，意味着政府支出在加速，资金从国库流向企业账户，这会推高M1。但这不是企业主动增加经营性现金，而是财政转移的被动结果。

\subsection{[R023] DB：市场真正低估的不是风险，而是结构性资金洪流对资产定价的重塑}
{\small\color{refgray}归类：宏观与利率：美联储鹰派定价充分，美元区间待突破}

DB：市场真正低估的不是风险，而是结构性资金洪流对资产定价的重塑

当大多数市场参与者还在争论利率走势与信用周期时，一份来自DB证券化研究团队的2026年中展望报告提出了一个更根本的判断：当前证券化市场的核心定价逻辑，已经从“信用风险定价”转向了“需求冲击定价”。

这份报告的核心洞察是：美国证券化市场正在经历一场由结构性资金流入驱动的定价重估。2026年全年证券化发行总量预计首次突破1万亿美元，但净供给仅为2700亿美元。与此同时，仅固定年金和IG债券基金两个渠道，就有约5000亿美元的资金涌入信用市场。供给侧的克制与需求侧的爆发，正在改写几乎所有资产类别的定价方程。

这不是一个关于周期拐点的判断，而是一个关于市场结构永久性变化的判断。理解这一点，比预测下一次降息时点重要得多。

1. 1万亿美元发行量背后，真正的故事是净供给只有2700亿

粗看2026年证券化市场的发行数据，很容易得出“供给过剩”的结论。DB预测全年总发行量将达到1万亿美元，创历史新高。但报告真正有价值的数据点，藏在净供给这个数字里。

全年净供给预计仅为2700亿美元。这意味着，超过70\%的新发行债券，实际上只是在替换到期或提前偿还的存量债券。市场的“净吸水效应”远比表面看起来温和。

以非代理RMBS为例，全年发行量预计达到2150亿美元，但考虑到1350亿美元的预期偿还，净供给只有800亿美元。ABS的情况类似，3600亿发行量对应3050亿偿还，净供给仅550亿。CLO的净供给相对较高，但绝对规模也仅为900亿美元。

这个结构意味着什么？对于投资者而言，新券的稀缺性远比发行量数字所暗示的要高。在需求端持续膨胀的背景下，这种供给结构天然形成了对利差的压制。报告明确判断，AAA级证券化利差将基本保持稳定，这是建立在供给结构支撑之上的结论，而非对宏观风险的忽视。

2. 5000亿资金洪流的源头，不是机构投资者，而是退休人群的资产负债表重构

如果说净供给结构是定价的供给侧逻辑，那么需求侧的故事则更加深刻。DB提出了一个在证券化研究中较少被讨论的视角：人口结构变化正在重塑固定收益市场的需求曲线。

报告指出，60岁以上人口在美国达到8400万，其中超过75\%拥有自有住房。这个群体在过去40年中经历了房价5.5倍的上涨和标普500指数70倍的回报。更重要的是，婴儿潮一代控制了美国60\%到70\%的金融财富。

这些数字的金融含义是：婴儿潮一代在固定收益资产上的配置比例，远低于其财富占比。过去20年，10年期美债平均收益率仅为2.9\%，在这样的利率环境下，退休人群缺乏配置固定收益的动力。但当利率升至当前水平——固定年金收益率可达5.5\%，且享有税收递延优势——需求结构正在发生根本性转变。

DB估算，仅固定年金销售一项，2026年就将达到约3000亿美元。加上IG债券和共同基金渠道的另外2000亿美元，总计5000亿美元的资金正在涌入信用市场。证券化产品在这其中扮演的角色，是帮助保险公司实现承诺的年金收益率。保险公司目前持有超过1万亿美元的证券化资产，这个数字仍在快速增长。

这不是一个短期交易性资金流入，而是退休人群基于生命周期进行的资产配置再平衡。这意味着需求的持续性可能远超市场预期。

对于熟悉信用市场的投资者来说，当前利差水平自然会引发一个疑问：AAA证券化利差普遍处于三年历史区间的最低端，这是否意味着估值已经充分甚至过高？

DB的判断是：利差可能维持低位，但并非没有风险。报告对2026年底各品种利差的预测显示，AAA级CMBS、CLO、信用卡ABS等品种的利差变化预期为零至几个基点。真正可能出现明显走阔的，是BBB级品种——BBB CMBS预计走阔50个基点，BBB CLO走阔10个基点，而BBB级非QM RMBS反而可能收窄

[... middle omitted ...]

 监管变化正在悄悄改变银行的需求曲线，但这一点市场尚未充分定价

在讨论证券化市场需求时，市场通常关注的是保险公司和共同基金。但DB报告揭示了另一个可能被低估的需求来源：银行。

2026年3月，美国监管机构重新提出了巴塞尔III终局规则。与2023年的初始提案相比，新提案对银行持有证券化产品明显更为友好。关键变化有三：P因子维持0.5不变（而非提案中的1.0），高级别证券的风险权重下限从20\%降至15\%，以及允许银行对高级别证券采用“穿透法”计算风险权重。

这些变化意味着什么？在旧规则下，银行持有高级别证券化产品的资本占用可能高于持有底层资产本身，这在经济上是不合理的。新规则修正了这一扭曲，将显著改善银行持有证券化产品的经济性。

报告指出，过去几年银行一直是证券化产品的净卖出方。如果监管变化能够扭转这一趋势，银行将从供给方变为需求方，这将进一步加剧市场的供需失衡。但这个判断仍需要验证——银行是否会真正增加配置，取决于其整体资产负债管理策略，而非仅仅监管资本成本。

\subsection{[R024] 摩根斯坦利：人形机器人市场真正被低估的不是需求，而是供应链的再定价}
{\small\color{refgray}归类：人形机器人：7.5万亿市场今年启动，供应链再定价}

摩根斯坦利：人形机器人市场真正被低估的不是需求，而是供应链的再定价

这份报告的发布时间是2026年6月，但它的判断框架值得今天就读懂。

摩根斯坦利亚太研究团队在最新的《Asia Summer School: Asia Robotics and Humanoids》中给出了一个极其明确的预测：到2050年，全球人形机器人市场规模将达到7.5万亿美元，保有量超过10亿台。这个数字本身已经足够震撼，但真正值得产业决策者和投资者关注的，不是终局数字有多大，而是这个市场正在以什么样的节奏、沿着什么样的路径、在哪些环节率先形成真实的经济性。

报告的核心判断可以提炼为一句话：人形机器人正在从“技术可行性验证”阶段，进入“小规模商业化”阶段。2026年就是这个转折点。而这一转折，将驱动整个供应链——从精密零部件到边缘计算芯片——出现一次前所未有的价值重估。

1. 商业化不是“即将到来”，而是“已经启动”，但信号藏在供应链里

摩根斯坦利明确判断，小规模商业化将在2026年开始。这不是一个模糊的展望。报告给出的证据链条是清晰的。

从资本市场的活动来看，2026年前四个月全球人形机器人领域的风险投资总额已经超过了2025年全年。多家初创公司在这一期间跻身独角兽行列：美国的Figure估值达到390亿美元，德国的Neura Robotics估值70亿美元，中国的银河通用、星海图、星动纪元等也相继突破10亿美元估值门槛。更值得注意的是，多家公司已经提交了IPO申请——Unitree、Deep Robotics、乐聚机器人都在列。

这意味着什么？资本市场的加速不是终点信号，而是起点信号。当一家行业还在VC/PE阶段就出现如此密集的IPO准备，说明产业资本和金融资本已经形成共识：商业化窗口已经打开，先发优势的窗口期可能比大多数人预期的要短。

但这里有一个容易被忽略的细节。报告在展示融资数据时，2026年的数据只覆盖到4月，就已经超过了2025年全年。如果这一节奏延续，2026年全年的人形机器人融资规模可能会是2025年的3到4倍。这个加速度本身，就是比任何终局预测都更值得关注的信号。

报告给出的7.5万亿美元TAM和10亿台保有量，是基于什么样的假设？这个问题决定了这个数字的可信度。

从报告的图表结构来看，这个预测隐含了几个关键假设：人形机器人将逐步渗透工业、商业服务和家庭三个场景；单台成本将从目前的数十万美元下降到2万美元以下；全球劳动力市场对人形机器人的接受度将逐步提升。

这三个假设中，最脆弱也最关键的是成本曲线。报告引用的AlphaWise调查显示，企业认为人形机器人的“甜蜜价格点”在10万到20万人民币之间，约合1.4万到2.8万美元。而目前市面上能看到的商业化产品，价格远高于这个区间。

这意味着什么？人形机器人的大规模商业化，本质上不是一个技术问题，而是一个成本问题。而成本问题的答案，不在整机厂，在供应链。谁能在精密减速器、力矩传感器、边缘计算芯片等核心零部件上实现降本，谁就掌握了这个行业最稀缺的定价权。

3. 供应链的“隐形爆发”正在发生，但大多数投资者还没注意到

报告最值得反复研读的部分，是对硬件BOM成本的拆解和供应链市场空间的预测。

根据摩根斯坦利的测算，到2040年，仅人形机器人相关的核心零部件市场总规模就将达到7800亿美元。其中，增速最快的不是整机，而是零部件。以2025年为基准，边缘计算芯片的市场规模将增长1904倍，达到1.23万亿美元；力矩传感器增长370倍，达到3160亿美元；电机增长170倍，达到2.07万亿美元；减速器增长157倍，达到1.37万亿美元。

这些数字的倍数关系远大于整机市场本身的增速。逻辑很简单：每一台人形机器人需要数十个关节、数十个传感器、一套完整的计算系统

[... middle omitted ...]

的全产业链布局已经形成。

但这并不意味着中国供应链可以高枕无忧。报告在“行业仍处于早期阶段”一节中明确指出，模型路径尚未收敛，数据稀缺且昂贵，硬件在成本、精度、效率和可靠性方面仍有提升空间。更重要的是，随着产量爬坡，数据安全、政府监管、地缘政治约束将成为新的变量。

对于中国的零部件企业来说，这意味着一个两难选择：如果跟随全球头部整机厂的技术路线，可能面临技术路径切换的风险；如果坚持自主技术路线，又可能因为与主流生态不兼容而失去规模化机会。这个矛盾，报告没有给出明确的答案，但它是所有关注这个赛道的投资者必须思考的问题。

这可能是整份报告最大的未解问题，也是我们选择不把报告所有内容讲完的原因——因为这个问题需要持续的跟踪和讨论。

报告提到了两种主流的技术路径：VLA（视觉-语言-动作模型）和World Model（世界模型）。前者强调端到端的感知-决策-执行闭环，后者强调机器人对物理世界的因果理解。两种路径各有优劣，但报告没有给出明确的判断，哪一条路径将最终胜出。

\subsection{[R025] NOM：消费数据跌破疫后低点，市场正在低估“无政策刺激”下的结构性分化}
{\small\color{refgray}归类：地产与消费：一线城市结构性再定价，消费疲软拉长政策等待期}

NOM：消费数据跌破疫后低点，市场正在低估“无政策刺激”下的结构性分化

2026年5月的社零数据，给所有还在观望消费复苏节奏的人一记明确的警钟。根据国家统计局6月16日公布的数据，中国社会消费品零售总额在5月录得同比下跌0.6\%，这是自2022年底中国经济刚走出疫情封锁以来，首次出现月度同比负增长。这一数字不仅远低于市场预期的微增0.1\%，更打破了此前关于“消费正在磨底”的温和叙事。

NOM这份最新研报的核心判断是：消费疲软正在从个别品类扩散为系统性现象，而市场此前对“房地产回暖带动财富效应”和“AI行情支撑消费信心”的期待，均未在数据中得到验证。更值得关注的是，报告明确指出，在没有实质性需求侧政策刺激的背景下，消费情绪在短期内难以逆转。

但这篇报告的价值不在于确认“消费不好”这一已知事实，而在于它提供了一个在总量下行中寻找结构性机会的分析框架。当大多数公司面临收入增长放缓的压力时，哪些企业能够凭借清晰的战略、可见的利润率趋势和有吸引力的估值，成为穿越周期的少数赢家？NOM给出了两个明确的标的，但其背后的选股逻辑，值得所有产业决策者认真拆解。

总量数据是-0.6\%，但真正令人担忧的是内部结构的普遍走弱。限额以上企业商品零售总额同比下跌5.2\%，成为拖累整体社零的主要力量。NOM在报告中指出，虽然部分下滑可以归因于高基数效应和5月不利天气条件——气温偏高、降雨日数增多限制了消费者活动——但这不足以解释如此广泛的疲软。

首先，餐饮服务收入增速从4月的2.2\%进一步放缓至5月的0.6\%，几乎逼近零增长。餐饮作为消费景气度的“温度计”，其持续走弱说明居民外出消费意愿正在快速收缩。

其次，商品零售中，绝大部分品类在5月都出现了增速下滑甚至恶化。家用电器和音像器材类同比下跌15.6\%，建筑及装潢材料类同比下跌13.6\%，这两个与房地产高度相关的品类已经深陷负增长区间，且毫无反转信号。

唯一录得增速加快的品类是饮料类，同比增长6.1\%，高于4月的3.6\%。服装鞋帽针纺织品类增速基本持平于3.8\%。办公用品类的跌幅从6.9\%收窄至1.5\%。这些零星亮点并未改变整体格局。

把这些事实合在一起，含义很清楚：消费的疲软并非某个特定领域的短期扰动，而是需求端信心不足在多个层面的同步映射。NOM将其归因为三个核心因素：缺乏有意义的消费侧政策刺激、房地产市场的回暖仅局限在少数一二线城市且财富效应有限、AI驱动的股市上涨对消费的支持力度不足。这三个因素在短期内都看不到明显改善的迹象。

市场上一度存在一种乐观预期：随着核心城市房地产市场在2025年下半年至2026年初出现企稳迹象，房价回升带来的财富效应将传导至消费端。但5月的数据让这一逻辑链条断裂了。

家用电器和建筑装潢材料这两个品类的持续深度负增长，直接说明房地产市场的回暖并未转化为装修、家电购置等实际消费需求。即便是一线城市二手房成交量的回升，也没有带动下游产业链的实质性复苏。这里需要区分两种财富效应：一种是资产价格上升带来的“纸面财富”增加，另一种是实际交易活跃带来的“变现消费”。当前的情况更接近前者，而后者尚未启动。

NOM报告虽然没有直接展开这一分析，但其数据已经暗示：房地产对消费的拉动作用，可能比市场预期的要弱得多、也慢得多。对于依赖房地产后周期需求的行业——包括家电、建材、家具、家纺等——这意味着2026年下半年的业绩预期可能需要进一步下修。

在消费整体承压的背景下，NOM维持了对两只股票的偏好：运动服饰领域的安踏体育和黄金珠宝零售领域的老铺黄金。这两只股票分属不同的消费子行业，但NOM的选股逻辑有一条清晰的主线：在收入增长放缓的环境中，企业的定价权、利润率趋势和战略清晰度比规模扩张速度更重要。

安踏的推荐逻辑是“清晰的品牌发展战略、可见的利润率趋势和

[... middle omitted ...]

逻辑则更具独特性。NOM认为，如果现货金价能够在更友好的地缘政治环境中企稳，老铺黄金将能够维持相对稳定的利润率趋势并改善运营杠杆。黄金珠宝是一个特殊的消费品类：它既受金价波动影响，又依赖品牌溢价和产品设计。老铺黄金的“古法金”定位使其在一定程度上脱离了传统金饰的原材料定价逻辑，获得了更高的毛利率和品牌溢价能力。NOM给出的目标价1114港元，对应2026财年预期市盈率22.5倍，反映了对其高增长轨迹的认可。

但这里有一个尚未完全回答的问题：如果金价出现大幅下跌，老铺黄金的利润率能否保持稳定？报告将金价大幅走弱列为主要风险之一，但并未深入分析金价下跌对老铺黄金商业模式的具体影响机制。这是一个值得进一步探讨的关键假设。

4. 报告尚未完全回答的关键问题：消费何时见底，以及什么能触发反转

NOM的报告在诊断现状方面做得非常扎实，但关于“消费何时见底”以及“什么能触发反转”这两个核心问题，报告并没有给出明确的结论。这并非报告的缺陷，而是当前数据所能提供的有限信息决定的。

\subsection{[R026] NOM：电商渗透率回升背后，消费结构的分化比总量更值得关注}
{\small\color{refgray}归类：未被主题摘要引用}

5月中国电商零售数据公布后，市场焦点大多落在“增速小幅回暖”这个数字上。NOM在最新发布的Quick Note中捕捉到一个更值得关注的信号：电商渗透率同比回升1.0个百分点至31.6\%，但各品类增长轨迹的分化程度，远超总量数据所能表达的范围。

这不是一个简单的“消费复苏”或“消费降级”的故事。NOM的数据拆解显示，5月电商销售同比增长2.6\%，较4月的0.2\%明显提速，但同期社会消费品零售总额（剔除汽车）增速却从1.8\%回落至1.1\%。线上与线下、必需品与可选品、食品饮料与家电手机——不同维度的增长曲线正在朝着相反方向移动。

对于关注中国消费市场的投资者而言，真正的问题不是“电商是否还能增长”，而是“增长的结构性驱动因素正在发生怎样的根本性转变”。这份报告提供的数据线索，指向一个比表面增速更复杂的判断：消费者在品类间的支出再分配，正在重塑电商平台的竞争逻辑。

1. 必需品与可选品的增速裂口，揭示的不是消费意愿而是消费优先级

NOM的数据显示，5月线上食品饮料销售在1-5月累计同比增长15.5\%，与4月的15.6\%几乎持平。这是一个值得注意的稳定性——在整体零售增速放缓的背景下，食品饮料作为刚需品类，其线上增长几乎没有受到宏观经济波动的侵蚀。

对比之下，可选消费品的表现出现了明显分化。线上服装销售累计增速从4月的6.8\%微升至7.2\%，但化妆品增速从4.7\%回落至2.5\%，家电销售更是录得15.6\%的同比下滑。同一份数据中，线下渠道的饮料销售增速反而从3.6\%加速至6.1\%。

这些数字合在一起意味着什么？消费者并没有停止花钱，而是在重新排列支出优先级。食品饮料和服装等“日常必需的可选品”依然获得稳定的线上增长，而家电、化妆品这类“可推迟的升级消费”则承受更大压力。这不是简单的“消费降级”，而是一种更理性的消费优先级排序：消费者愿意为每日所需的品质改善付费，但对大额、低频的耐用品支出变得更加谨慎。

对于电商平台而言，这意味着品类结构的差异将直接影响GMV增长的质量。以食品饮料和服装为核心品类的平台，可能比依赖家电和3C产品的平台拥有更强的增长韧性。

5月智能手机线上销售同比仅增长0.7\%，较4月的6.2\%大幅回落。NOM分析指出，这主要归因于高基数效应和内存芯片成本上涨引发的终端涨价。

这个判断值得深入推敲。手机销售增速的放缓，并非因为消费者“不想换手机”，而是因为“换机成本上升”。内存芯片价格持续上涨，正在向终端售价传导，而消费者对价格变动的敏感度在当前的宏观环境下被进一步放大。

这是一个关键的结构性信号：上游成本压力正通过供应链向消费端传导，而消费电子品类可能首当其冲。如果这一趋势持续，手机和家电等品类的线上销售可能在未来几个季度继续承压。对于电商平台而言，这意味着高单价品类的增长贡献可能减弱，平台需要更加依赖高频、低客单价的日常消费品来维持整体增速。

NOM报告没有完全展开的一个问题是：这种成本传导是否已经触发了消费者行为的结构性变化？如果消费者因为价格上涨而延长换机周期，这种影响可能持续12-18个月，而不仅仅是月度波动。

3. 渗透率回升的真正含义：线上渠道正在吃掉线下份额，但吃的是“不同”的份额

5月电商渗透率同比上升1.0个百分点至31.6\%，这是报告中最容易被忽视但最具战略意义的数据点。在整体零售增速放缓的背景下，渗透率提升意味着线上渠道正在从线下夺取更多份额。

但关键在于，线上夺取的份额结构正在发生变化。从品类数据看，食品饮料和服装等“高频、低单价”品类的线上渗透正在加速，而家电等“低频、高单价”品类的线上渗透可能已经接近阶段性天花板。这意味着电商平台的增长逻辑正在从“扩品类”转向“深耕高频品类”。

对于平台而言，这个转变有两个含义。第一，高频品类的竞争将更加激烈，价

[... middle omitted ...]

分析家电销售下滑时明确提到了“trade-in subsidies dwindling”——以旧换新补贴正在退坡。这是报告中最具洞察力的一句话，但也是一个没有完全展开的伏笔。

补贴退出的影响不仅仅是“少了一个促销工具”。它可能引发一系列二阶效应：补贴期间被提前透支的需求，将在后续月份形成更高的基数压力；消费者习惯了补贴价格后，对正常价格的接受度可能下降；平台和品牌商为了维持销量，可能被迫自掏腰包补贴，侵蚀利润率。

这些效应的量化难度极高，但恰恰是投资者最需要理解的部分。NOM报告没有给出具体的弹性测算，这可以理解——任何机构在月度快报中都很难完成这种深度分析。但对于决策者而言，这是一个必须持续跟踪的风险变量。如果家电和手机等品类在补贴退出后出现持续性的销售低迷，电商平台的品类结构风险将被放大。

5. 一个值得建立的观察框架：从“总量增速”转向“品类轮动速度”

基于NOM这份报告的数据逻辑，投资者可以建立一个更精细化的观察框架，而不是仅仅盯着电商大盘的月度增速。

\subsection{[R027] Citi：进口车市场正在经历的不是周期性波动，而是品牌溢价能力的根本性分化}
{\small\color{refgray}归类：新能源车：渗透率突破67\%，竞争逻辑从抢增量转向抢存量}

Citi：进口车市场正在经历的不是周期性波动，而是品牌溢价能力的根本性分化

中国进口车市场正在经历一场超出多数人预期的结构性调整。Citi最新发布的基于ThinkerCar保险数据的研报显示，2026年5月，中国进口乘用车零售销量同比下滑38\%，环比零增长，当月销量仅为2.93万辆。前五个月累计销量16.48万辆，同比下跌28\%。

这些数字本身已经足够引人注目。但比总量下滑更值得关注的，是不同品牌之间的表现差异——这种差异指向一个更深层的判断：进口车市场正在从“品牌驱动”向“产品力与价格匹配度驱动”切换，而大多数传统豪华品牌尚未完成这一转变。

Citi这份报告的价值不在于它披露了销量数字——这些数据市场已有感知——而在于它提供了一个清晰的、可横向比较的品牌表现图谱。透过这张图谱，我们可以识别出哪些品牌正在失去议价权，哪些品牌尚能维持基本面，以及更重要的是，哪些结构性变量正在重塑这个市场的竞争逻辑。

1. 雷克萨斯和奔驰的相对韧性说明，品牌护城河正在从“豪华认知”转向“供需纪律”

在几乎所有品牌都录得两位数下滑的背景下，雷克萨斯和奔驰的表现值得单独拿出来讨论。雷克萨斯5月销量10,516辆，同比下滑38\%，但环比持平；前五个月累计销量60,123辆，同比下滑19\%。奔驰5月销量6,271辆，同比下滑19\%，环比增长3\%；前五个月累计34,366辆，同比下滑25\%。

从绝对数字看，这两个品牌同样身处下行通道。但从相对位置看，它们明显优于同行。宝马5月销量3,356辆，同比下滑41\%；保时捷2,582辆，同比下滑32\%；奥迪947辆，同比下滑66\%；路虎2,035辆，同比下滑45\%。

这里的核心差异不在于品牌知名度——这些都是全球顶级豪华品牌——而在于供需管理能力。雷克萨斯在中国市场长期坚持“不压库”策略，经销商库存深度维持在较低水平。奔驰在过去两年中也逐步调整了进口车型的定价和配额策略。这意味着当需求收缩时，这些品牌面临的价格倒挂压力相对较小。

反观奥迪和沃尔沃，5月同比分别下滑66\%和66\%，前五个月累计分别下滑53\%和57\%。这种幅度的下滑已经不能用“消费疲软”来解释。它更接近于一种信号：当品牌在终端市场失去定价权后，经销商体系的信心会迅速瓦解，形成“降价-观望-再降价”的负向螺旋。

这里隐含的判断是：进口车市场的下一轮洗牌，赢家不会是品牌力最强的，而是供需纪律最严的。

2. 大众进口车的逆势增长是一个值得警惕的异类信号，而非趋势反转

在所有品牌中，大众是唯一实现正增长的。5月销量511辆，同比增长84\%，环比增长86\%；前五个月累计1,603辆，同比增长3\%。考虑到大众进口车在中国市场长期处于边缘地位，这个数据很容易被解读为“性价比车型正在崛起”。

但我们认为，这个判断需要谨慎对待。大众进口车目前的销量基数极低——月均仅300余辆——在这个量级上，单一经销商的大宗采购或某个车型的集中交付就足以造成数据的大幅波动。Citi报告本身也没有对这一异常值做出额外解读，这本身就说明该机构认为这不具备趋势意义。

更值得关注的是，大众进口车在中国市场的产品线以途锐等高端车型为主，其定价区间与二线豪华品牌直接重叠。如果大众的逆势增长能够持续三个季度以上，才可能意味着“豪华品牌溢价正在被务实消费逻辑侵蚀”这一判断成立。但在当前阶段，它更可能是一个统计噪音。

真正需要警惕的是：如果大众这种低基数品牌都能在特定月份实现增长，而奥迪、沃尔沃等传统豪华品牌却面临断崖式下滑，说明问题不出在“市场有没有需求”，而在于“消费者愿不愿意为这个品牌的溢价买单”。

3. 保时捷的环比改善掩盖了一个更深层的隐忧：超豪华品牌的客群正在缩圈

保时捷5月销量2,582辆，同比下滑32\%，但环比大幅增长27\%。从环比角度看，

[... middle omitted ...]

暗示：保时捷前五个月累计销量刚刚超过1万辆，而雷克萨斯同期是6万辆。这不仅仅是价格区间的问题——保时捷的客单价远高于雷克萨斯——它反映的是“可买可不买”型消费正在系统性收缩。

对于关注超豪华品牌资产定价的读者来说，这是一个需要持续跟踪的信号。如果保时捷的销量在接下来的三季度仍无法回到同比持平的水平，那么超豪华品牌在中国市场的估值逻辑可能需要重新审视。

4. 报告尚未完全回答的关键问题：进口车下滑有多少是“国产替代”造成的，又有多少是“需求萎缩”造成的

这是Citi这份报告最有价值但也最未充分展开的地方。报告提供了详尽的品牌层面数据，但它没有区分两种完全不同的下滑原因：

一种情况是，消费者仍然有购买该品牌的需求，但转向了该品牌的国产车型。例如，宝马X5国产化后，进口X5的需求自然会下降。这种情况下，品牌在中国的总销量（国产+进口）可能并未显著下滑。

另一种情况是，消费者从整体上减少了对该品牌的需求，无论国产还是进口。这种情况下，品牌在中国市场的基本面正在恶化。

\subsection{[R028] 摩根斯坦利：市场对日本汽车股的悲观定价，忽略了三个尚未被充分理解的变量}
{\small\color{refgray}归类：未被主题摘要引用}

摩根斯坦利：市场对日本汽车股的悲观定价，忽略了三个尚未被充分理解的变量

2026年6月，摩根斯坦利日本汽车团队在纽约和旧金山与机构投资者进行了一轮密集会议。这份调研反馈报告透露出的信号，与当前股价反映的市场共识之间存在一个值得注意的偏差：绝大多数投资者对日本汽车股的判断框架，几乎完全被两个因素主导——中东局势对供应链的冲击，以及中国OEM崛起对市场份额的侵蚀。但这份报告暗示，真正影响个股定价差异的变量，可能并不在这两个已经被充分讨论的叙事里。

当市场对某个行业的担忧高度趋同时，往往意味着定价已经部分吸收了这些风险。而那些尚未被广泛定价的结构性因素，才是超额收益或风险暴露的真正来源。

1. 丰田的盈利指引被误解为保守，但市场低估了HEV带来的结构性利润弹性

报告中最值得关注的细节之一，是摩根斯坦利团队的一个判断：市场对丰田F3/27营业利润指引的保守性质“理解并不充分”。这句话的潜台词是，丰田给出的指引数字很可能远低于公司内部的实际预期。

为什么丰田会选择给出一个明显偏保守的指引？这背后可能涉及日本企业一贯的“低预期管理”策略，但也可能反映了管理层对 年期间多个不确定性的主动对冲。然而，真正让这份报告值得深读的，是投资者对HEV（混合动力车型）盈利贡献的关注度。在纯电动车渗透率增速放缓、充电基础设施瓶颈尚未完全解决的背景下，丰田在HEV领域的技术积累和规模效应，可能正在转化为一种被市场低估的利润缓冲垫。

这意味着：如果中东局势出现任何缓解信号，丰田的盈利修正幅度可能比市场预期的更大。而如果局势持续恶化，HEV较高的利润率也能为丰田提供比纯电动路线更厚的安全垫。这种“不对称的盈利弹性”，目前尚未被充分定价。

2. 美国-墨西哥-加拿大协定2026年7月复审，是日本车企面临的最被低估的政策风险

报告提到，投资者对USMCA（美国-墨西哥-加拿大协定）将于2026年7月进入复审阶段表现出高度关注。这一点值得特别重视，因为很多市场参与者仍然将北美业务的盈利预期建立在现有贸易框架之上。

对于丰田、本田、日产这些在墨西哥拥有大量产能的日本车企而言，USMCA复审可能带来的原产地规则收紧、区域价值含量要求提高，将直接影响其北美供应链的布局成本和利润结构。摩根斯坦利团队在报告中专门提到这一点，说明它已经超出了“地缘政治噪音”的范畴，正在成为影响个股估值的一个可量化因素。

这里的关键在于：不同日本车企在墨西哥的产能占比、本地化采购率、以及北美市场的依赖度各不相同。这意味着USMCA复审对每家公司的影响，将呈现出显著的差异性。目前市场对这一风险的定价，可能过于笼统地反映在整个板块的折价中，而没有区分个股之间的暴露差异。

3. 日产的自由现金流前景，是当前日本汽车股中最值得追踪的“压力测试”案例

报告指出，尽管许多投资者对日产在固定成本削减方面的进展表示肯定，但市场对其能否结束汽车业务亏损并实现正自由现金流，仍抱有显著疑虑。这种疑虑的定价，可能已经在日产相对较低的估值中有所体现，但问题在于：市场是否已经充分考虑了日产在转向电动车过程中利润率恶化的风险？

报告中提到，投资者对日产在电动车业务上的利润率恶化风险表现出“相当大的关注”。这是一个典型的“双重挤压”情景：一方面，传统燃油车业务的利润正在被削减固定成本的压力压缩；另一方面，电动车业务的资本开支和初期亏损正在消耗现金流。如果日产无法在 年之间实现自由现金流转正，其融资成本和战略灵活性将进一步受到制约。

这不仅仅是一个公司层面的问题。日产的情况，实际上是整个日本汽车行业在转型期面临的一个缩影——传统业务的利润正在被结构性因素侵蚀，而新业务的投资回报仍存在高度不确定性。谁的资产负债表和现金流能够支撑这场转型，谁就能在下一轮竞争中占据主动。

报告显示，投资者

[... middle omitted ...]

被削弱。

更重要的是，报告提到投资者对中东局势对印度整体经济的影响表示担忧。这意味着铃木面临的风险，不仅仅是公司层面的经营问题，而是其核心市场正在经历一场宏观层面的需求冲击。这种“国家风险”与“公司风险”的叠加，可能还没有被充分反映在铃木当前的估值中。

5. 斯巴鲁的长期盈利前景，取决于成本削减措施能否填补产品线空白

报告指出，除了最新的股东回报政策调整，投资者对斯巴鲁的长期盈利前景表现出高度兴趣，特别是管理层在2025年经营计划中提出的降低主要成本、充实产品线等措施。斯巴鲁的案例提供了一个观察日本车企转型策略的典型样本：当一家公司的产品线相对集中（如斯巴鲁对SUV的高度依赖），其盈利弹性就高度依赖于新车型的成功和成本结构的优化。

这里的关键问题是：斯巴鲁的成本削减措施，是否足以在现有产品线老化、新车型尚未完全放量的“空窗期”内维持利润率？如果答案是否定的，那么市场对斯巴鲁的估值逻辑就需要重新审视。报告没有给出明确判断，但这恰恰是值得深入阅读完整报告的切入点。

\subsection{[R029] 摩根斯坦利：能源基础设施的定价权正在从宏观叙事转向微观资产}
{\small\color{refgray}归类：油运与能源基建：供给侧弹性被低估，补库周期是真正定价锚}

这份摩根斯坦利本周发布的北美中游与可再生能源基础设施周报，提供了一组看似孤立、实则紧密关联的信号。报告覆盖了从霍尔木兹海峡的地缘政治博弈到Permian盆地管道扩容的微观工程进展，从中国原油进口的急剧收缩到AI数据中心项目的搁浅。将这些信号放在一起，一个更值得关注的判断浮现出来：当前市场对能源基础设施资产的定价，正在从依赖宏观地缘叙事转向聚焦微观资产的运营现实与合约结构。这一转变意味着，过去两年基于“能源安全”和“AI电力需求”的板块普涨逻辑，正在被更精细的资产选择逻辑所取代。

为什么现在重要？因为市场正在经历一个“压力测试”窗口。霍尔木兹海峡可能重新开放的前景正在压低原油价格，而AI相关股票的动量因子正在经历今年以来最剧烈的回撤。在这两个负面因素同时作用下，中游能源基础设施资产的表现分化将比以往任何时候都更能揭示哪些公司真正拥有结构性护城河，哪些只是搭上了主题的便车。

这份报告最值得细读的地方在于，它没有停留在宏观判断，而是提供了多个具体资产层面的数据点——从Waha天然气价格的反弹驱动因素，到单个管道项目的在役时间线，再到AI数据中心项目的客户承诺缺失。这些微观证据，才是判断未来6-12个月能源基础设施投资方向的关键。

1. 霍尔木兹海峡的“投降式”重开并不会终结中游资产的吸引力，反而会加速资本向拥有长期合约的资产集中

报告开篇就抛出了一个大胆的判断：如果美伊达成协议重新开放霍尔木兹海峡，这将被市场视为一个“投降事件”，促使资本从避险资产重新轮动回能源板块。这个判断看似矛盾——地缘风险溢价消失，为什么反而利好能源？但仔细推敲，逻辑是成立的。

关键前提在于时机与不确定性。报告指出，交通流量正常化和全球库存补充需要时间。这意味着即便协议达成，实际供应恢复也不是一蹴而就的。更重要的是，中国过去几个月大量削减原油进口（5月进口量较冲突前水平下降近500万桶/日），但这并非需求崩溃的信号。摩根斯坦利的商品策略师Martijn Rats的分析揭示了一个关键细节：中国在冲突前一直在超量进口、积极建立商业和战略库存。因此，进口削减是在库存充裕的背景下发生的，实际终端需求仅下降了约120万桶/日。

这意味着什么？一旦地缘局势明朗，中国可能迅速从去库存转向补库存，从而对油价形成支撑。对于中游基础设施而言，无论是管道、存储还是出口终端，其收入主要来自长期合约而非现货价格波动。因此，油价因地缘缓和而短期承压，并不会实质性影响这些资产的现金流稳定性。相反，地缘不确定性的降低反而可能释放被压抑的资本支出决策，推动新的管道和LNG项目推进。

这一逻辑的隐含含义是：那些对油价波动暴露度高的上游资产，与那些拥有长期收费合约的中游资产，将在同一地缘事件下面临截然不同的定价逻辑。投资者需要区分的是，哪些公司的收入结构真正独立于油价方向。

2. Permian天然气外输瓶颈的缓解不是利空，而是下一轮增长的催化剂

报告提供了一个非常具体的案例来支撑上述判断。根据Hart Energy的报道，Kinder Morgan的Gulf Coast Express（GCX）管道扩容已经投入服务，增加了570 MMcf/d的压缩能力。这一扩容的直接结果是：Waha（Permian盆地）天然气价格从5月下旬开始出现“急剧且持续的反弹”。

这个案例的关键洞察在于——管道瓶颈的缓解，不是天然气中游资产的利空（因为通常认为新增运力会压缩现有管道的利用率），而是利好。原因在于，Permian盆地的天然气伴生气产量一直在增长，而外输能力不足导致区域气价长期处于深度贴水状态，迫使生产商关井或减少钻探活动。一旦新增运力到位，生产商可以恢复产量，从而为管道系统带来更多的输送量。

报告进一步指出，WhiteWater的2.5 Bcf/d Blackcomb管

[... middle omitted ...]

同样过剩的天然气枢纽的管道。报告没有直接点名，但这一逻辑指向了Targa Resources等拥有Permian集输和加工资产的公司，其业务将直接受益于产量恢复。

3. AI数据中心的电力需求叙事正在经历现实检验，天然气基础设施的长期逻辑反而更坚实

报告提到一个值得高度关注的事件：Crusoe在怀俄明州的AI数据中心项目因未能获得客户承诺而停滞，谷歌对其成本结构和时间表表示担忧。同时，报告指出AI/电力相关股票的动量因子正在经历今年以来最严重的回撤，这直接拖累了Williams Companies、DT Midstream和Kinder Morgan等天然气基础设施股票的表现。

但报告给出了明确的立场：这种弱势表现“过度且没有反映2H26可能出现的新项目公告”，而新项目的支撑来自“加速的超大规模云服务商AI投资”。报告中的图表显示，亚马逊、谷歌、微软、Meta和Oracle的合计AI资本支出预计将从2025年的约2000亿美元增长到2027年的约4000亿美元。

\subsection{[R030] NOM：中国焦煤供应冲击的真正影响，不在中国，而在印度钢厂的利润结构}
{\small\color{refgray}归类：锂电与焦煤：库存认知分裂，供应冲击暴露结构性脆弱}

NOM：中国焦煤供应冲击的真正影响，不在中国，而在印度钢厂的利润结构

市场正在错误地解读近期山西煤矿事故的后续影响。多数投资者关注的是中国国内焦煤价格是否还会继续上涨，或者澳大利亚海运煤出口能否填补缺口。但NOM这份最新报告揭示了一个更值得注意的判断：真正需要重新定价的，不是中国焦煤的短期供应，而是印度钢铁企业的利润韧性。

5月26日山西煤矿事故发生后，中国有137座焦煤矿被暂停生产。截至6月11日，虽然77座矿井已经复产，但恢复进程并不均衡——部分矿井在重启几天后再次停产。国内焦煤现货价格已上涨超过10\%至每吨270美元。这些信息已被市场消化。但报告真正有价值的部分，在于它如何拆解这一事件对不同市场参与者的差异化影响。

1. 这轮供应冲击暴露了中国煤炭监管的深层模式，而非一次性事件

事故后复产速度看起来很快——56\%的停产的矿井在两周内恢复。但NOM的数据显示，恢复路径并非线性。6月11日的数据甚至显示，复产矿井数量从81座回落至77座，同时停产矿井数从56座回升至60座。这不是一次性的生产中断，而是一个反复震荡的过程。

更关键的是，当局已将安全检查扩展至主要产煤区，重点打击非法和隐蔽采掘巷道。这意味着监管压力不会随着事故热度消退而解除。NOM将这一模式称为“监管风险叙事”——中国煤炭供应的核心变量正在从产能本身转向政策执行力度。

对中国市场而言，这意味着国内焦煤价格很难迅速回到事故前水平。每吨270美元的现货价格，可能不是短期情绪溢价，而是市场在为一种新的供应常态定价。

2. 澳大利亚出口激增的背后，是结构性替代正在发生，但规模仍有上限

5月澳大利亚对中国焦煤出口环比暴增114\%，同比增长77\%，达到148万吨。澳大利亚在中国焦煤进口中的份额从约7\%升至约11\%。这组数字看起来很大，但放在中国整体焦煤需求面前，仍然只是边际变化。

NOM的逻辑链值得关注：如果中国国内供应恢复持续不稳定，需求将被迫转向海运市场。而澳大利亚是最直接的受益者。但报告也指出了两个关键约束：一是中国可以通过内陆进口（蒙古、俄罗斯）部分替代，但物流瓶颈和品质差异限制了替代的充分性；二是澳大利亚优质硬焦煤现货价格仅微涨1.5\%至每吨244美元，远低于中国国内价格的涨幅。

这意味着全球海运焦煤市场目前仍处于“等待确认”状态。真正的价格上行压力，取决于中国国内复产进度是否持续低于预期。如果复产节奏在7-8月仍无法加速，全球焦煤定价中枢可能面临系统性上移。

3. 印度钢厂才是这条传导链上最值得关注的变量，但市场可能高估了风险

对于印度钢铁企业而言，焦煤成本是利润的核心驱动力。NOM算了一笔账：焦煤价格每上涨10美元/吨，对综合性钢厂的EBITDA影响约为7-9美元/吨，具体取决于库存周期和配煤比例。

在NOM覆盖的印度钢厂中，JSW Steel和Tata Steel的进口焦煤敞口最高，Jindal Steel因采购结构不同，暴露程度相对有限。这组敏感性分析本身并不令人意外。真正有价值的判断在于：NOM认为，中国焦煤供应的逐步恢复，将使海运焦煤价格保持相对稳定，因此这一事件对印度钢厂近期利润预期不应构成增量负面冲击。

这个结论依赖于一个关键假设：中国国内复产不会出现二次中断。而报告的图表9已经显示，复产进程已经在6月11日出现反复。如果这一趋势延续，印度钢厂的利润预期可能需要重新审视。

4. 报告没有完全回答的问题：中国HRC出口利润能否承受成本上升

报告图10显示，中国出口热轧卷板（HRC）的利润已从2021年的高点每吨550美元大幅回落至目前的约170美元，低于历史中位数水平。如果国内焦煤价格持续高企，而HRC出口价格无法同步上涨，中国钢厂的出口利润将进一步被压缩。

这里存在一个尚未被充分讨论的二阶效应：中国

[... middle omitted ...]

比回落至约36,000卢比，但仍处于健康水平。如果焦煤成本上涨被国内钢价上涨部分对冲，利润压缩幅度将远小于敏感性计算给出的数字。但如果国内需求走弱，钢价无法转嫁成本，情况则完全不同。

NOM目前对Tata Steel、JSW Steel、Jindal Steel和Lloyds Metals均维持买入评级。但这个评级成立的前提，是中国焦煤供应恢复进程不出现新的重大扰动。

这份报告最值得深思的地方，不在于它告诉你了什么，而在于它没有完全回答的问题：中国监管检查的持续强度如何量化？蒙古和俄罗斯的替代供应是否存在隐性瓶颈？印度钢厂的库存周期能否缓冲下一轮价格冲击？这些问题的答案，往往藏在报告的图表细节和假设条件里。

欢迎来星球微信群里继续讨论这些未解问题，我们可以一起拆解报告的完整逻辑、对比不同机构的假设差异、跟踪后续数据变化。

Personal reading notes and learning share only. Not investment advice.

\subsection{[R031] GS：市场低估了资本支出超级周期对资产定价的深远重塑}
{\small\color{refgray}归类：资本市场：资本支出超级周期重塑资产定价}

过去四十年，投资者习惯了低利率、低通胀、低波动和全球化带来的确定性回报。但GS在最新一期全球策略报告中提出了一个核心判断：我们正站在一个全新周期的起点，这个周期的底层逻辑与1980年代以来的“现代周期”截然不同。市场目前对AI和基础设施投资带来的资本支出浪潮有所定价，但真正被低估的，是这一资本支出超级周期将如何系统性重塑股票回报的构成、行业竞争格局以及资产定价的锚。

这份报告的核心洞察在于，投资回报的驱动力正在发生根本性转移。从依赖估值扩张和多重压缩的“金融工程”时代，转向依赖盈利增长和现金流的“实业创造”时代。理解这一转变，比预测下一个季度的GDP数据或美联储的降息时点更为重要。

1. 过去四十年的“现代周期”已不可逆地终结，其底层驱动因素正在全面逆转

GS将1982年至疫情前的时期定义为“现代周期”。这一周期的核心特征是：低通胀、低利率、低波动、全球化、低资本密集度和高利润率。这些因素共同造就了历史上最长的股票牛市，也塑造了整整一代投资经理的思维范式。

但这份报告清晰地指出，支撑这一周期的所有支柱都在瓦解。从低通胀转向再通胀，从量化宽松和零利率转向更高的资金成本，从全球化转向区域化和碎片化，从低资本密集度转向大规模基础设施投资。这些不是短期的周期性波动，而是结构性转变。例如，全球军事支出占GDP的比重在经历了数十年的下降后，自2022年以来已经重新上升。企业税率在经历了长期下降后，也面临上行压力。劳动力成本在全球化退潮和人口结构变化下，其下降趋势已经逆转。

这意味着，投资者不能再简单地将过去四十年的经验外推。那种“买入并持有”就能享受估值扩张和全球化红利的日子，可能已经过去了。未来的投资回报将更加依赖对结构性变化的深度理解和主动选股能力。

2. 资本支出超级周期已全面启动，从资本节约转向资本密集型增长是核心变量

报告中最具洞察力的部分，是对资本支出超级周期的详细阐述。GS指出，这一周期的驱动力来自三个相互强化的方向：AI革命驱动的私营部门投资，能源安全和地缘政治驱动的公共部门投资，以及供应链重构带来的资本密集度系统性提升。

AI的兴起正在催生一个全新的增长引擎，这不仅仅是软件层面的创新，更涉及大规模的计算基础设施、数据中心和能源网络建设。与此同时，对能源安全的重视和地缘政治紧张局势，正在触发一轮全球性的公共投资浪潮。各国政府都在加大对本土制造业、半导体、新能源和国防工业的投资。供应链从追求效率转向追求韧性，这意味着企业需要建立更多的库存、备用的生产线和多元化的采购渠道，所有这些都推高了资本密集度。

这一转变的含义是深远的。过去二十年，企业通过外包和轻资产模式实现了高利润率。而现在，企业需要重新进行大规模的资本投入。这将对企业的自由现金流、资产负债率和估值模型产生系统性影响。那些能够成功将资本支出转化为未来盈利增长和市场份额的企业，将获得市场的奖励；而那些无法有效转化资本投入的企业，则可能面临资本回报率下降的风险。

3. 股票回报的驱动力从估值扩张转向盈利增长，选股逻辑需要根本性调整

GS报告中的一个关键图表演示了这一点：在“现代周期”中，尤其是在金融危机后的零利率时代，股票回报的主要驱动力是估值扩张。经济复苏虽然弱于历史均值，但股市的反弹却异常强劲，因为大量的流动性推高了资产价格，而非实体经济产出。

然而，在“后现代周期”中，GS认为这种模式正在改变。随着利率和资金成本上升，估值进一步扩张的空间受到限制。未来的股票回报将更多地依赖于企业盈利的实际增长。这意味着，投资者需要更加关注企业的基本面，尤其是其单位数增长、利润率趋势和资本配置效率。

这对于投资风格的影响是巨大的。那些依赖低利率环境进行杠杆收购、股票回购和估值套利的策略可能不再有效。相反，那些拥有强大定价权、能够将成本压力转嫁

[... middle omitted ...]

，“上升的横截面离散度增加了阿尔法生成的空间”。这意味着，对于具备深度研究能力的投资者而言，未来的市场环境比过去几年更加友好。被动投资虽然成本低廉，但可能无法捕捉到这种结构性分化带来的超额收益。主动选股的价值将重新凸显。

5. 报告尚未完全回答的关键问题：资本支出转化为盈利的“转化率”如何？

尽管GS的报告对资本支出超级周期的到来给出了强有力的论证，但它也留下了一个核心问题有待市场验证：这些巨额的资本支出，究竟能在多大程度上转化为未来的盈利增长？或者说，企业的资本回报率（ROIC）是否会因为竞争加剧和成本上升而系统性下降？

报告提到了“资本密集型扩张”的转变，但没有详细分析不同行业和企业的资本支出效率差异。例如，AI基础设施的投资热潮，最终是催生出少数几家拥有垄断地位的赢家，还是导致大量重复建设和产能过剩？能源转型领域的投资，其回报周期和风险特征与传统能源有何不同？这些问题的答案，将直接决定哪些公司能够成为资本支出超级周期的受益者，哪些公司可能陷入“投资陷阱”。

\subsection{[R032] Citi：AI产业真正的分水岭不是模型能力，而是推理成本的精细化控制}
{\small\color{refgray}归类：AI/算力：智能体时代CPU复兴，中国供应链份额跃迁}

Citi：AI产业真正的分水岭不是模型能力，而是推理成本的精细化控制

2026年6月的Research and Applied AI Summit（RAAIS）传递了一个清晰但容易被市场噪音掩盖的信号：AI行业的竞争焦点，正在从“谁拥有最大的模型”转向“谁能以最低的token成本交付可用的推理结果”。这不是一个渐进式的优化，而是一个可能重塑产业链价值分配的结构性转折。

Citi在最新发布的研报中系统梳理了RAAIS 2026的核心技术趋势，并将“推理成本效率”和“递归自我改进”列为当前最值得关注的产业变量。这份报告的价值不在于罗列技术名词，而在于揭示了两个正在并行发生的深层变化：第一，前沿模型的边际收益正在递减，而推理成本的高速增长正在倒逼整个行业重新设计技术栈；第二，AI从“语言模型”向“世界模型”的跃迁，其商业化路径远比市场预期的要长，但一旦突破，其影响将远超当前任何单一应用场景。

以下是我们从这份研报中提炼的五个核心洞察，以及它们对产业决策者和投资者的含义。

1. 推理成本正在成为AI商业化的第一约束，而“降本”本身已经催生出一个新的技术层

Citi研报中两个数字值得反复咀嚼。Revolut的ML工程负责人披露，其内部开发的模型路由平台通过“为任务匹配正确规模的模型”，实现了8倍的成本节约且质量无损。ElevenLabs的研究工程师则展示了在固定硬件上实现70倍吞吐量提升的技术组合，其中仅连续批处理（continuous batching）一项就贡献了15倍增益。

这些数字背后是一个正在成型的产业逻辑：当模型能力达到一定阈值后，企业竞争力的核心不再是“能否调用最强模型”，而是“能否在维持质量的前提下，将推理成本压到对手无法跟进的水平”。Revolut的做法尤其具有代表性——它建立了一个中央网关和优雅降级机制的路由平台，其核心洞察是“前沿模型往往是错误的选择”。这意味着，对于大量实际业务场景，使用小型、专用模型不仅成本更低，而且效果可能更好。

这个趋势的直接含义是：一个围绕“推理效率”的新技术层正在形成。它包含模型路由、量化感知训练、多token预测推测解码、KV缓存量化、蒸馏等技术模块。这些模块的供应商——无论是云服务商、芯片公司还是中间件厂商——都有可能在这一轮重构中获得新的定价权。而对于应用端企业，能否快速构建自己的推理效率体系，将在未来12-18个月内成为决定其AI投资回报率的关键变量。

2. “递归自我改进”正在从概念走向工程化，但它的真正瓶颈不在算法，而在基础设施

去年RAAIS上，DeepMind展示了“AI Squared”——AI通过递归自我改进来提升自身能力的概念。今年，这一主题被进一步具体化为一个框架：将科学发现视为一个强化学习问题，从发现到开放式发现，再到加速的开放式发现，其中元学习（meta-learning）加速了发现的速度。

Citi报告提到，支撑这一框架的基础设施已经出现：DiscoGen支持超过4亿个生成式研究任务，MLGym-Bench和Rainbow Teaming则提供了评估和对抗训练的环境。Odyssey的PROWL更进一步，它训练了一个强化学习代理来改进世界模型本身，而不仅仅是决策策略。其哲学是：“你的代理永远不会比你学到的环境模型更好。”

这里的核心洞察是：递归自我改进不是一种算法优化，而是一种新的研发范式。它需要大规模的基础设施来生成、评估和筛选改进方案。这意味着，能够负担这种基础设施投入的玩家将越来越少。对于大多数AI公司来说，参与这一赛道的门槛正在急剧升高，而真正的受益者可能是那些提供“研究基础设施即服务”的平台。

但报告没有完全展开的一个问题是：递归自我改进的“自我”边界在哪里？当改进的对象从模型参数扩展到世界模型，再到研究流程本

[... middle omitted ...]

印象深刻的进展：Odyssey-2 Max在PAI-Bench的Physics维度上获得了93分，Starchild-1实现了实时的像素和音频联合生成，Agora-1则能够在单个H100上支持约17个并发玩家的多代理会话。

但Citi报告也明确指出，世界模型的商业化仍处于早期阶段。其当前最被看好的应用场景是作为数据生成引擎，打破机器人领域的遥操作瓶颈。DeepMind的Genie 3和他们的双部分机器人建模方案展示了向教育、代理训练和灵巧操作延伸的可能性。

这里的核心问题是：世界模型的“GPT-2时刻”到“ChatGPT时刻”之间需要多长时间？如果参照语言模型的发展轨迹，这个窗口期可能长达2-3年。在这期间，哪些基础设施投资是必要的，哪些是过早的？对于投资者而言，现在需要关注的不是某个世界模型公司的估值，而是支撑世界模型训练和推理的硬件、数据和仿真平台是否正在形成壁垒。

4. AI在医疗健康领域的价值被语言模型框架严重低估，生物学数据的模型化可能是一个更大的机会

\subsection{[R033] JEF：锂电产业链的真实信号不是价格下跌，而是库存认知的分裂}
{\small\color{refgray}归类：锂电与焦煤：库存认知分裂，供应冲击暴露结构性脆弱}

JEF：锂电产业链的真实信号不是价格下跌，而是库存认知的分裂

这份JEF锂电产业链调研报告，覆盖了中国市场约77\%的锂盐转化企业和73\%的电池材料生产企业，是当前最有代表性的产业链一线感知采样。读完整份报告，最值得关注的核心判断不是价格还会跌多少——市场对此已有充分预期——而是一个正在悄然形成的结构性裂痕：上游转化企业和下游电池企业对库存水平的认知出现了方向性分歧。这种分歧，正在重塑整个锂电产业链的议价机制和利润分配格局。

报告显示，转化企业认为电池厂商的库存“过低”，但电池厂商却认为自己的客户——整车厂和储能集成商——存在“适度过剩”的库存。同一组数据，两个环节读出了完全相反的结论。这意味着什么？意味着产业链的信息传导已经不再是对称的，每一环都在根据自身订单感受做出独立判断，而这些判断正在驱动截然不同的采购和备货行为。

这份调研发生于一个微妙的时点。锂价在经历了2022年的狂飙和2023年的回调后，市场普遍认为底部区域已经探明，但去库存周期何时结束、需求何时真正企稳，始终缺乏高频、一手的数据支撑。JEF的这份调研，恰好填补了这一空白。它不是基于模型推演，而是基于实际负责采购和排产的一线管理者的直接反馈。

调研中一个最直接的信号是价格预期。转化企业预计未来三个月锂盐价格将出现超过10\%的下跌，而电池企业的预期则温和得多，仅为个位数跌幅。值得注意的是，从2024年初至今，转化企业的价格预期曲线始终在电池企业下方，且差距在2025年下半年开始明显扩大。

这意味着什么？转化企业正在承受比下游更重的价格压力。原因不难理解：转化环节产能过剩更为严重，且产品同质化程度高，价格竞争几乎是唯一的竞争手段。电池企业虽然也面临降价压力，但其产品差异化程度更高、客户粘性更强，价格传导机制相对缓和。

但更值得关注的不是当前价格预期的绝对值，而是趋势。从调研数据看，转化企业的价格预期在2025年11月曾短暂收窄至接近零，但2026年初再次恶化，5月又回落至-12\%的深度负值区间。这种反复波动说明，市场并未真正找到价格底部，任何短期企稳都可能被新一轮供给释放或需求波动打破。

对于投资者而言，这意味着锂盐板块的盈利底部仍难以确认。即便价格已经跌至许多企业的现金成本以下，但只要产能出清没有实质性发生，价格反弹的空间和持续性就始终存疑。

如果说价格预期是情绪指标，那么订单增速就是更硬的基本面信号。调研中，电池企业的新订单同比增速自2024年以来始终维持在0\%-5\%的窄幅区间，表现相当稳定。而转化企业的新订单增速则持续为负，直到2025年底才勉强回到零附近。

更关键的是两者的交叉验证。从“预期销售增速减新订单增速”这一差值指标来看，转化企业自2025年下半年以来持续为正，意味着他们预期销售增速高于新订单增速。这通常是一个危险信号——预期高于实际，往往意味着后续的库存被动累积和更激进的价格让步。电池企业的这一差值则持续为负，说明他们的预期更为保守，实际订单可能好于预期。

这组数据合在一起指向一个清晰的判断：产业链的定价权和增长动能正在从上游向下游转移。电池企业凭借更强的客户关系和更稳定的需求端，正在成为产业链中更强势的一环。转化企业则陷入“量价双杀”的困境——量在萎缩，价在下跌，利润空间被双重挤压。

对于产业决策者而言，这意味着战略重心需要向下游延伸。纯上游资源或转化环节的商业模式，在当前的供需格局下面临的结构性挑战远不止是周期性的。那些能够向下游电池材料、甚至电芯环节延伸的企业，将获得更强的抗风险能力和议价空间。

这可能是整份报告最值得深入解读的部分。转化企业和电池企业对“客户库存水平”的判断出现了系统性的方向性分歧。

转化企业认为电池厂商的库存“过低”——这一比例在2026年3月降至仅10\%，是调研周期内的最低点。而电

[... middle omitted ...]

身订单感受，容易受到订单波动的影响而放大判断偏差。

但这里有一个更值得警惕的可能性：如果电池企业的判断是正确的，即终端客户确实存在库存过剩，那么当前电池企业的稳定订单增速可能只是“假象”——终端客户正在消化库存，新采购需求并未真正释放。一旦去库存完成，电池企业的订单增速可能面临一轮更明显的下滑，届时上游转化企业将承受更猛烈的冲击。

这个假设尚未被验证，但它指向一个对投资者至关重要的观察窗口：未来一到两个季度，终端整车和储能项目的实际出货数据，将比任何调研问卷都更有说服力。

JEF的调研清晰地描绘了产业链各环节的当前状态，但它没有——也不可能——回答一个更根本的问题：产能出清何时真正开始，以及什么事件会触发它。

从调研数据看，转化企业的库存同比增速自2025年下半年以来持续回升，到2026年初已经转正。这意味着，尽管价格在跌、订单在萎缩，但企业并没有真正削减库存，反而在累积。这通常不是主动补库的信号，而是被动积压——企业无法以当前价格卖出产品，库存只能被动上升。

\subsection{[R034] Bernstein：液冷CDU的真正分化不在技术路线，而在推理与训练的两极定价}
{\small\color{refgray}归类：未被主题摘要引用}

Bernstein：液冷CDU的真正分化不在技术路线，而在推理与训练的两极定价

当市场还在争论液冷与风冷谁将主导下一代数据中心时，一份来自Bernstein的研报揭示了一个更微妙的信号：技术路线的选择不是核心变量，真正决定CDU制造商未来估值分化的，是AI工作负载的“推理”与“训练”两极分化，以及由此引发的CDU产品分层定价。

这份报告讨论的是谷歌近期发布的“Brazos”项目——一款面向开放计算项目的液冷到空气CDU参考设计。表面看，这只是又一个技术规格的公布。但Bernstein的分析指向一个更深层的判断：这不是一次技术迭代，而是一次对CDU市场结构的重新定义。它意味着市场正在被切割为两个截然不同的价格带——一个属于高溢价、高技术壁垒的训练CDU，另一个属于可标准化、可商品化的推理CDU。

对于投资者而言，理解这个分化的含义，远比追踪下一个液冷技术突破更为紧迫。

1. 谷歌Brazos的真实意图不是竞争，而是为推理负载开辟一条低成本通道

研报明确指出，Brazos并非谷歌要亲自制造的CDU产品。它是一套开放的技术规格，供Vertiv、nVent、Boyd等OCP生态成员生产。其冷却容量仅为60kW，甚至无法支撑一个Blackwell机架（需要约120kW）。从参数上看，它并不“惊艳”。

但Bernstein的洞察在于：这种“不够强”恰恰是设计意图的体现。Brazos被设计为直流供电，适配已有直流基础设施的旧有数据中心。60kW的冷却能力虽然无法支持前沿训练模型，却恰好覆盖推理负载的常见功率范围。

这背后有一个更宏大的逻辑：尽管训练模型需要新建高密度数据中心，但推理负载的增长速度预计在2030年前将超过训练。而推理负载的部署面临的最大瓶颈不是技术能力，而是数据中心容量的交付速度。如果能通过简单的L2A CDU改造现有设施来承载推理任务，将大幅缩短交付周期。Brazos正是为此而生。

所以，这不是一个技术竞争的故事，而是一个市场分层的故事。训练和推理需要不同的基础设施策略，而谷歌正在用Brazos为推理负载铺一条低成本、快部署的路径。

2. 推理CDU面临真实的商品化风险，但训练CDU的护城河仍在加深

这是Bernstein报告中最具投资洞察力的部分。研报提出了两个中期风险，其中第一个尤为关键：推理CDU生态系统的商品化风险。

Brazos的技术规格比之前的Deschutes项目（液冷到液冷，2MW容量）要简单得多。这意味着，满足Brazos规格的CDU不需要像Vertiv或nVent这样的顶级工程能力。更多中小型制造商可以参与竞争，价格压力将随之而来。虽然服务附件的价值依然存在，但整体利润率将远低于旗舰训练CDU。

这里的核心变量是推理机架的功率密度。如果超大规模客户的专用推理芯片机架功率密度维持在60kW/机架以下，那么相当一部分推理负载可以用Brazos风格的L2A CDU来冷却，商品化风险将加速兑现。反之，如果推理密度持续上升，则需要更大容量的Deschutes类CDU，这些产品的利润率和服务粘性预计更高。

换句话说，推理CDU的利润率取决于一个尚未明确的参数——推理机架的未来功率密度。这个变量目前仍存在不确定性，但它的走向将直接决定CDU制造商的价值分布。

3. 绿色field到改造的迁移风险：一个尚未被数据证实但值得警惕的信号

第二个中期风险来自需求结构的潜在变化。如果超大规模客户因项目延迟而将部分需求从新建数据中心转向现有设施的改造，那么Brazos这类标准化L2A CDU的需求将上升。Bernstein的数据中心容量追踪器尚未显示“搁浅容量”的异常，但已有来自产业链公司的传闻，称部分客户因准备不足或项目延误而推迟交付。

这个信号之所以重要，是因为它可能改变CDU制

[... middle omitted ...]

谜。

第一个未解问题是推理机架的未来功率密度。这个参数将决定Brazos类CDU的市场规模，进而决定CDU制造商在推理领域的利润空间。但Bernstein坦承，目前尚不清楚超大规模客户专用推理芯片的机架密度会走向何方。

第二个未解问题是，商品化是否会从推理CDU蔓延至训练CDU。报告认为，目前训练CDU仍享有创新溢价，且随着顶级液冷技术的演进，这种溢价预计将持续。但问题在于，如果Brazos式的标准化路径在推理领域获得成功，超大规模客户是否会尝试将类似的标准化逻辑应用于训练领域？目前没有证据支持这种推断，但这是一个值得持续跟踪的方向。

第三个未解问题是，服务附件的价值在推理和训练两个领域究竟有多大差异。报告暗示，训练CDU的故障成本更高，因此服务粘性更强。但推理CDU的服务价值是否会被低估？如果推理负载的规模足够大，服务附件是否能弥补硬件利润的下降？这些问题的答案，需要更深入的产业链调研。

5. 投资者需要建立的观察框架：从“谁做得好”到“谁在哪个市场做得好”

\subsection{[R035] GS：日本银行股正在定价的利率路径，可能比市场意识到的更陡}
{\small\color{refgray}归类：区域市场：日本央行鸽派加息，银行股定价偏保守}

日本央行在6月议息会议上将政策利率从0.75\%上调至1.0\%，这本身并不令人意外。真正值得关注的信号，藏在加息之后的两个细节里：主要银行在存款利率上调中保持了与市场预期一致的存款贝塔系数，以及横滨金融集团率先披露了加息对下财年利润的正面影响。GS这份报告的核心判断是，市场对日本银行股的定价仍然偏保守——当前股价隐含的政策利率假设大约在0.75\%至1.0\%区间，而央行传递的加息周期延续信号，意味着这个区间需要重新校准。

这不是一个关于“加息利好银行股”的简单叙事。真正的问题在于，哪些银行能够将加息转化为可持续的利润增长，以及市场何时会开始为更深的利率正常化定价。报告提供了几个关键的分析锚点，值得逐层拆解。

1. 存款利率上调和贷款利率上调之间的利差，决定了加息的实际利润转化率

加息对银行利润的影响，从来不是政策利率变动的简单函数。关键在于净息差的变化，而净息差取决于资产端和负债端利率调整的节奏差。

从报告披露的数据看，三菱UFJ金融集团将普通存款利率从0.30\%上调至0.40\%，上调幅度10个基点，同时将短期优惠利率从2.125\%上调至2.375\%，上调幅度25个基点。这意味着存款利率对政策利率的贝塔系数维持在40\%，与市场此前预期一致。三井住友金融集团和Mizuho金融集团随后跟进，将普通存款利率同样上调至0.40\%。

这个数字很关键。40\%的存款贝塔系数意味着，每100个基点的政策利率上调，银行只需要将其中40个基点传导给存款人，剩余60个基点则转化为净息差的扩张。报告中的测算显示，对于以0.75\%政策利率假设编制指引的银行，一次25个基点的加息对净利润的影响在1.1\%至13.5\%之间，差异巨大。Mizuho金融集团的利润弹性约6.5\%，而Resona控股高达13.5\%。

这种差异的背后，是各家银行存款结构和贷款定价能力的根本不同。存款粘性更强的银行、贷款定价更灵活的银行，在加息周期中的利润弹性更大。报告没有完全展开的一个问题是，存款贝塔系数是否会随着加息周期的深入而上升——当政策利率从1.0\%进一步向1.5\%或更高水平移动时，存款人的议价能力会增强，银行的存款成本上升速度可能加快。这个假设如果被验证，目前基于40\%贝塔系数的利润弹性测算就需要向下修正。

2. 多家银行以0.75\%政策利率假设编制指引，加息至1.0\%后存在系统性的上调空间

报告列出了一个关键表格：覆盖银行的盈利指引及其隐含的政策利率假设。三菱UFJ金融集团以1.0\%的政策利率假设编制了2026财年指引，但三井住友、Mizuho、静冈银行、横滨金融集团、千叶银行、福冈金融集团等多数银行仍以0.75\%作为假设。

这意味着什么？这些银行在编制2027财年指引时，有相当大的可能性需要上调政策利率假设至1.0\%或更高。每一次假设上修，都意味着盈利预测的上调。横滨金融集团已经率先行动，披露了加息对2028财年税前利润约140亿日元、税后约100亿日元的影响，ROE提升约0.7个百分点。

报告暗示，这种盈利上修不是一次性的。随着加息周期延续，银行可能会像横滨金融集团一样，逐次披露加息对各财年利润的影响。如果2027财年中期计划中隐含的政策利率假设是0.75\%至1.0\%，那么当前1.0\%的政策利率已经触及这个区间的上沿。如果央行继续加息至1.25\%或1.5\%，这些中期计划将面临系统性上调。

这里还需要注意一个尚未完全回答的问题：银行在披露加息影响时，是否已经充分考虑了存款成本上升的二阶效应？横滨金融集团的披露是基于当前存款贝塔系数的静态测算，如果存款贝塔系数随利率上升而上升，实际利润弹性会低于披露数字。这个动态关系，报告没有完全展开，但值得投资者自行跟踪。

3. 市场持仓偏轻叠加加息周期确认，银行股可能经历类似2025年12月加息后的

[... middle omitted ...]

续，市场有理由开始为更深的利率正常化定价。

报告附了一张关键图表，展示了Topix银行股在加息前后30天的表现。从图表数据可以看出，加息至0.50\%时，银行股在加息前30天至加息后30天基本持平；但加息至0.75\%时，银行股在加息后30天内从基准1.05上涨至1.20，涨幅约14\%。如果历史规律成立，当前加息至1.0\%后，银行股可能重复类似的走势。

但这个类比有一个需要谨慎对待的前提：前两次加息发生在利率绝对水平较低的环境中，市场对加息周期的持续性存在分歧。当前政策利率已升至1.0\%，继续加息的阻力可能更大，尤其是在日本经济增速放缓或全球衰退风险上升的情况下。报告对此没有给出明确判断，而是将问题留给了读者。

4. 报告尚未完全回答的关键问题：利率正常化的终点在哪里，以及哪些银行会掉队

GS这份报告的核心贡献，是提供了一个清晰的定价框架：银行股当前价格隐含的利率路径，与央行实际传递的路径之间存在差距。但报告没有回答的一个关键问题是，这个差距最终会以什么方式收敛。

\subsection{[R036] 摩根斯坦利：国防产业正经历一场“去空心化”的结构性重塑，而非简单的周期上行}
{\small\color{refgray}归类：资本市场：资本支出超级周期重塑资产定价}

摩根斯坦利：国防产业正经历一场“去空心化”的结构性重塑，而非简单的周期上行

国防产业的拐点，往往不是由军费预算的增减来定义的。真正值得关注的信号，是资本结构、技术扩散路径和产业组织形态同时发生系统性变化。摩根斯坦利在其首届国家安全创新峰会的总结报告中，捕捉到了这样一个时刻。这份报告的核心判断并非“国防开支在增长”，而是：美国国防产业正在经历一场从“冷战后的空心化”向“新中层的重建”的深层转型。驱动这一转型的，不是单一的政策指令，而是资本、技术和地缘竞争三重力量在供给侧的协同作用。对于关注这一领域的投资者和决策者而言，理解这场结构性重塑，远比追逐季度订单数字更为关键。

1. 冷战后的“和平红利”正在逆转，产业空心化催生了新中层的崛起

冷战结束后，美国国防产业经历了一轮大规模的整合浪潮。其直接后果是，中型国防企业（Mid-tier Defense Assets）被大量兼并或消失，产业格局呈现出“两头大、中间小”的哑铃型结构：一端是少数几家巨头（如洛克希德·马丁、诺斯罗普·格鲁曼），另一端是大量小而专的技术初创公司。摩根斯坦利指出，这种“空心化”在相当长一段时间内制约了产业的灵活性和创新效率。而现在，随着大国竞争回归，国防需求从“维持存量”转向“快速迭代与规模化生产”，这种结构性缺陷变得不可持续。

报告揭示了两股力量正在共同“填补”这个中空地带。其一是国防硬件的商业化。许多原本仅用于军事领域的传感器、通信和动力技术，开始在商业市场找到应用场景，这降低了单一客户依赖风险，也吸引了更广泛的资本。其二是风险资本的介入。传统上，国防领域因其长周期、高门槛和客户单一性，对风险资本缺乏吸引力。但现在，一批愿意承担早期技术和市场风险的“颠覆性资本”正在涌入，它们不追求短期回报，而是押注于能够定义未来作战范式的平台级公司。这两股力量的汇合，正在催生一个全新的国防产业“中层”。这个新中层的公司，既不是巨头，也不是作坊，而是具备规模化生产能力和技术领导力的“新精锐”（Neoprimes）。这不仅仅是产业结构的调整，更意味着估值逻辑的迁移：当国防投资者基础从行业专家向成长型投资者扩展时，那些能够证明“品类领导者”地位和可扩展生产能力的公司，将获得更接近高增长科技公司的估值倍数。

2. 从“技术优越性”到“工业韧性”：生产能力本身已成为战略资产

摩根斯坦利在多个专题讨论中反复强调一个观点：在当代冲突中，工业韧性（Industrial Resilience）的重要性正在超越单纯的技术优越性。乌克兰冲突已经提供了一个清晰的案例——能够快速、低成本、大规模生产武器的能力，与拥有最先进的武器系统同样重要。这意味着，国防制造商的竞争优势来源正在发生根本性转移。

报告明确指出，下一代国防制造商的差异化优势，将不再主要来自产品层面的创新，而是来自“压缩资格认证、产品认证和生产周期”的能力。过去，国防采购流程以“定制化”为核心，每个项目都有独特的需求和漫长的认证流程，这天然地与规模化生产相矛盾。而现在，面对快速演变的威胁环境，适应性和速度成为新的评价标准。那些能够将“工厂”本身打造成一个敏捷、可复用的战略资产的公司，将获得更大的市场份额。这个判断对投资者的含义是：评估一家国防企业时，其工厂的柔性、供应链的纵深以及压缩生产周期的能力，可能比其最新武器的技术参数更具长期价值。同时，这也在政策层面提出了一个尖锐的问题：美国五角大楼基于一年预算周期且频繁出现持续决议（CR）的拨款模式，是否能够与这种对长期、稳定、可预测生产环境的需求相匹配？报告认为，政治层面的解决方案是解锁多年度采购路径、为产业提供必要可见性的关键前提。

报告中关于“AI赋能的战场”和“战场自主性”的讨论，揭示了一个远比“无人机取代有人战机”更深刻的趋势。摩根斯坦利认为，自主性本身

[... middle omitted ...]

同和生产可视化的软件层。它们将成为新战场生态的“操作系统”。同样引人深思的是，在“太空情报”和“情报融合”领域，摩根斯坦利观察到，限制因素越来越不是技术能力，而是“组织效率”。未来价值创造可能更多流向那些能够简化数据访问和决策工作流的公司，而非仅仅是改进收集硬件的公司。这意味着，在国防科技投资中，软件定义、平台化思维和组织设计能力，其重要性正在超越传统的硬件工程能力。

4. 报告尚未完全回答的关键问题：新中层的可持续性与“去耦合”的现实成本

摩根斯坦利的这份报告提供了一个极具洞察力的顶层框架，但它也留下了一些悬而未决的关键问题，值得深度探讨。首先，新国防“中层”的可持续性如何？这些“新精锐”公司能否在维持技术创新的同时，应对国防采购固有的官僚主义和长周期？它们是否会重蹈过去“和平红利”时期中型公司的覆辙，在下一轮整合中被巨头收购，从而再次导致产业空心化？报告提到了“专业化与去整合”的潜力，但这需要资本市场的持续耐心和国防部采购流程的实质性改革，两者都充满不确定性。

\subsection{[R037] 摩根斯坦利：市场低估的不是油运需求，而是供给侧的再定价}
{\small\color{refgray}归类：油运与能源基建：供给侧弹性被低估，补库周期是真正定价锚}

过去三周，油轮运输板块股价跑输大盘17至27个百分点。市场给出的理由很清晰：霍尔木兹海峡关闭时间超出预期，全球原油库存持续消耗，市场担忧油运需求萎缩，进而导致运力过剩、即期运价承压。

但摩根斯坦利在最新发布的一份研报中提出了一个与市场共识截然相反的判断——如果海峡真的重新开放，油轮运输板块非但不是利空出尽，反而可能迎来一轮持续性更强的运价支撑。这份报告的核心价值不在于预测运价能涨多少，而在于指出了市场定价中一个被系统性忽视的结构性变量：供给侧的真实弹性远低于市场假设。

市场在过去三周抛售油轮股的逻辑链条是：海峡关闭导致原油运输需求下降，运力过剩，运价承压。这条逻辑在表面上看没有问题，但它隐含了一个关键假设——油运需求的下降是永久性的。

摩根斯坦利的报告给出了完全不同的解读。海峡关闭期间，全球原油库存持续消耗，进口国补库意愿被压制，出口国在波斯湾的储油也无法顺利释放。一旦海峡重新开放，这些被压制的需求不会缓慢恢复，而是可能在短时间内集中释放。进口国需要重建库存，出口国需要清空储油，两股力量叠加，油运需求可能在一两个季度内出现脉冲式增长。

这意味着，过去三周市场定价的“需求萎缩”很可能只是暂时的流动性收缩，而非趋势性恶化。真正值得关注的不是需求是否会恢复，而是当需求以超预期速度回归时，供给端能否接得住。

2. 即期运价不会一次冲高后迅速回落，真正的定价锚在补库周期的长度

市场对“运价冲高后迅速回落”的担忧有其合理性。历史上，突发事件后的运价脉冲往往难以持续。但摩根斯坦利在这份报告中提出了两个更值得关注的变量，它们决定了运价均衡水平而非短期波动。

第一个变量是全球石油补库周期的长度。报告明确判断，这一轮补库可能持续数个季度，甚至更长。如果补库需求确实持续两到三个季度，那么即期运价在冲高后不会迅速回到事件前水平，而是会在一个更高的均衡区间内震荡。这与市场普遍预期的“脉冲式上涨后快速回落”存在本质差异。

第二个变量是伊朗石油制裁的走向。报告没有展开讨论，但将其列为可能支撑需求的关键因素。如果制裁出现松动，伊朗原油重新进入全球市场，将额外增加油运里程需求。这一点在报告的风险因素部分被反复提及，说明分析师认为其概率并不低。

这两个变量的共同点是：它们都是供给侧的结构性因素，而非需求侧的周期性波动。市场目前定价的更多是短期情绪，而非这些结构性变量的潜在影响。

市场另一个常见的误判在于对波斯湾运力的评估。直观上看，海峡关闭期间大量油轮被困在波斯湾，一旦重新开放，这些运力将迅速释放，压制运价。但摩根斯坦利的分析揭示了一个更细致的现实：目前波斯湾内的油轮大多是满载状态，短期内不会变成可用运力。

更关键的是，部分油轮在关闭期间已经绕行至大西洋市场。这意味着即使海峡重新开放，中东地区的有效运力供给也不会立即恢复到事件前水平。运力从大西洋重新调配回中东需要时间，而这个时间窗口恰好与补库需求的高峰期重叠。

这种运力错配是市场定价中最容易被忽视的结构性因素。市场看到的是“波斯湾有船”，但没有区分“有船”和“有可用船”之间的本质差异。

摩根斯坦利对两家中国油运龙头二季度业绩的预期对比，提供了一个非常清晰的验证场景。

对于招商轮船，报告预期二季度净利润环比增长超过50\%，主要受益于3月VLCC即期市场在事件后的冲高以及干散货市场的同步走强。这是一个典型的“事件受益型”标的，其核心优势在于资产布局与事件冲击的方向高度一致。

而对于中远海能，报告预期二季度净利润环比仅小幅增长，理由包括部分船舶被困波斯湾导致有效运价被稀释、中东成品油出口减少拖累国际成品油运输、以及中国原油进口放缓压制内贸利润。

这两家公司的对比说明了一个关键问题：在油运板块，公司层面的业绩分化不是规模决定的，而是由资产敞口的方向和灵活性决定的。招商轮船

[... middle omitted ...]

运价的支撑将是有限的，当前股价的反弹空间也会被压制。但如果补库能够持续3到4个季度甚至更长，那么油运板块的盈利中枢可能被系统性上调。

这个问题之所以关键，是因为它直接决定了当前估值是否合理。摩根斯坦利给出的目标价和估值倍数已经隐含了补库持续性的假设，但这一假设并未在报告中得到充分论证。投资者在参考这份报告时，需要自行评估补库周期的长度及其对盈利的影响。

另一个未解问题是伊朗制裁的走向。报告将其列为上行风险，但没有给出概率判断。如果制裁真正出现松动，全球油运需求将面临一次更大幅度的结构性增长，这可能是当前市场完全没有定价的变量。

6. 投资者的观察框架：从“事件驱动”转向“结构性供需再平衡”

基于这份报告的分析，投资者在评估油轮运输板块时，可能需要调整自己的观察框架。

第一，关注运力供给的真实弹性。市场习惯性地认为运力是弹性的、可快速调配的，但这次事件表明，有效运力的恢复速度可能远低于市场预期。投资者需要更细致地跟踪波斯湾运力的实际可用性，而非仅仅看总量数据。

\section{Disclaimer}
{\scriptsize\color{refgray}
This community is a paid learning and information-sharing community created for general educational purposes only. The membership fee is charged solely for access to community discussions, educational content, organization, and operational costs. It is not an advisory fee, management fee, performance fee, brokerage fee, or compensation for personalized financial, investment, legal, tax, or professional advice.\par
I participate and share information solely as an individual and for educational discussion among community members. I am not acting as a financial adviser, investment adviser, broker, dealer, portfolio manager, legal adviser, tax adviser, accountant, fiduciary, or any other licensed professional adviser. Nothing shared in this community constitutes, and should not be interpreted as, financial advice, investment advice, legal advice, tax advice, a recommendation, solicitation, offer, endorsement, instruction, or guarantee to buy, sell, hold, short, trade, or invest in any security, cryptocurrency, fund, derivative, strategy, or other financial product.\par
All content shared in this community, including messages, discussions, links, files, charts, examples, market commentary, watchlists, AI-generated or AI-polished summaries, and third-party materials, is general, impersonal, and educational in nature. It does not take into account any individual's financial situation, investment objectives, risk tolerance, investment horizon, liquidity needs, tax status, legal restrictions, or suitability requirements. No content should be relied upon as suitable for any specific person, account, portfolio, or financial decision.\par
Information shared in this community may be incomplete, inaccurate, outdated, speculative, AI-generated, AI-edited, or based on third-party internet sources that have not been independently verified. I make no representation or warranty as to the accuracy, completeness, timeliness, reliability, or suitability of any information shared.\par
Financial markets involve substantial risk, including the possible loss of principal. Past performance, historical data, hypothetical examples, simulated results, backtests, personal opinions, or educational examples do not guarantee future results. Each member is solely responsible for conducting their own research, exercising independent judgment, and making their own decisions. Before making any financial, investment, legal, or tax decision, members should consult qualified and licensed professionals.\par
Participation in this community, payment of any membership fee, or communication with me or other members does not create any adviser-client, fiduciary, professional, contractual, agency, partnership, employment, or other advisory relationship. By joining or participating in this community, you acknowledge and agree that you are solely responsible for your own decisions, actions, transactions, profits, losses, risks, and consequences, and that I shall not be liable for any direct, indirect, incidental, consequential, financial, legal, tax, or other loss or damage arising from or related to any content, discussion, information, or material shared in this community.\par
}
\end{document}
